\chapter{Determinant norms and finite descent}
\index{determinant norm}\index{finite descent!envelope theorem}
\label{chap:bj-finite-descent}

\section{Orientation}
\label{sec:ch41-orientation}

\subsection{What has to be transported, and in which direction}
\label{sec:ch41-what}

Let
\[
                         h:T\longrightarrow X
\]
be a finite surjective morphism of normal integral projective varieties over
the trivially valued field $k$, and let
$\theta\in N^1(X)_{\mathbf R}$.  This chapter proves one implication:
if every bounded-above family of $h^*\theta$-psh functions on $T$ has a
psh regularized supremum, then the same is true for every bounded-above
family of $\theta$-psh functions on $X$.  In the notation of
\Cref{def:bj-envelope-property} the assertion is
\[
 \operatorname{EP}_T(h^*\theta)
 \ \Longrightarrow\
 \operatorname{EP}_X(\theta),
\]
and it is \Cref{thm:bj-finite-normal-descent}.

The shape of the argument is dictated by an asymmetry that is worth naming
before any construction.  \emph{Going up} is free.  Pullback along
$h^{\An}$ sends a Fubini--Study presentation to a Fubini--Study
presentation and a decreasing net to a decreasing net, so it maps
$\PSH_X(\theta)$ into $\PSH_T(h^*\theta)$; this is
\Cref{prop:ch41-finite-pullback}, and follows by checking the section presentation, class convergence and
nontriviality in
\Cref{def:bj-psh}.  \emph{Coming down} is not free.  An
$h^*\theta$-psh function on $T$ is a limit of sections of line bundles
on $T$, and there is no reason at all for those sections to descend: the
pullback map $H^0(X,mL)\to H^0(T,mh^*L)$ need not be surjective.
This explains the obstruction to descending a chosen section
presentation; descent of the resulting function also requires constancy
on analytic fibres and cannot be characterized by that cokernel alone.

What the chapter constructs is not a section-by-section descent but an
\emph{averaging} operator.  Given a weighted system of sections upstairs,
one produces a weighted system of sections downstairs whose Fubini--Study
function is the multiplicity-weighted fibre average
\[
 z\longmapsto\frac1d\sum_{w\mid z}e(w/z)\,\varphi(w),
 \qquad
 \sum_{w\mid z}e(w/z)=d,
\tag{41.1}\label{eq:ch41-average-shape}
\]
of the given function $\varphi$ upstairs.  The averaging is performed by a
\emph{determinant norm}: the coefficients of the determinant of
multiplication by a universal linear combination $\sum_jT_js_j$ on the
flattened finite locally free algebra $q_*\cO_W$ are the sections one needs.
The original finite map $h$ need not be flat.  When $\varphi$ is
already pulled back from $X$, the fibre average returns the function it
was pulled back from, and the operator becomes a descent.

Two features of \eqref{eq:ch41-average-shape} are what make the chapter
technically delicate, and both are demanded by the theorem this chapter
serves.

\emph{The weights are not multiplicities of geometric points.}  They are
dimensions of the local Artinian factors of the fibre algebra over the
completed residue field, equivalently sums of their geometric Artin lengths
after an isometric extension to an algebraically closed field.  In positive
characteristic they retain the full field degree, including inseparability
and possible defect; over a point with nonreduced fibre they retain
nilpotents.  Counting analytic lifts alone omits their field degrees.
After algebraic closure, counting distinct geometric points without their
Artin lengths omits nonreduced structure.  Neither count supplies the
unconditional identity $\sum_we(w/z)=d$ needed below.

\emph{The construction must survive at points of proper support.}  The
family is bounded above but its members may take the value $-\infty$.
For a single Fubini--Study stage, the polar locus is the analytification
of a common zero locus, by
\Cref{prop:ch34-generic-finiteness}\,\cref{itm:ch34-gf:polar};
this does not imply the same description for an arbitrary psh limit.
The scheme-theoretic fibre at a proper-support point may also contain
nilpotents.  \Cref{ex:ch41-nonreduced-fibre} exhibits both phenomena in
one calculation.  Neither phenomenon is confined to proper support in
general: inseparable extensions may acquire nilpotents after completion,
and a general psh polar locus need not be algebraic.

\subsection{Why linear averages and unweighted sheet products do not suffice}
\label{sec:ch41-naive}

Since the goal is an average, it is natural to look first for an averaging
operator on sections.  The trace and the Reynolds operator do not supply the required
construction.  A product over embeddings is a valid separable-field norm,
but a fibrewise count without multiplicities does not cover all fibres.
The determinant construction treats these cases uniformly.

\emph{(i) The trace.}  For $q:W\to Z$ finite locally free of rank $d$
there is a trace map $\Tr_q:q_*\cO_W\to\cO_Z$, and one might hope that
$\frac1d\Tr_q$ implements the average.  It does not, for two independent
reasons.  First, even in the best case the trace is a \emph{sum}, and the
ultrametric inequality bounds the absolute value of a sum by the maximum of
the absolute values, so $\log|\Tr_q(s)|$ is bounded above by
$\max_{w\mid z}\log|s(w)|$ and admits no lower bound whatsoever; the
maximum is not the average, and there is no inequality in the useful
direction.  Second, and fatally, the trace can vanish identically.  If
$k$ has characteristic $p$ and $F:\mathbf P^1\to\mathbf P^1$ is the
$p$-th power map, then in the basis $1,u,\ldots,u^{p-1}$ of
$k[u]$ over $k[u^p]$ every element $u^i$ with $0\le i\le p-1$ has
trace $0$: for $i=0$ because $\Tr(1)=p=0$, and for $1\le i\le p-1$
because multiplication by $u^i$ shifts the basis indices modulo $p$
and has no diagonal entry when $1\le i<p$.  The trace map is the
zero map.  This is computed in full in \Cref{ex:ch41-frobenius-average}.

\emph{(ii) The Reynolds operator.}  If the cover is Galois with group
$G$, one may try $|G|^{-1}\sum_{\sigma\in G}\sigma^*$.  This requires
$|G|$ to be invertible in $k$, which is exactly what fails when the
characteristic divides the degree, and it requires the cover to be Galois,
which a general finite surjective morphism is not.  Both hypotheses are
unavailable: \Cref{chap:bj-alterations} feeds this chapter the finite part
of a Stein factorization of an alteration, over a field of arbitrary
characteristic, with no control on the Galois group or on the degree modulo
$p$.

\emph{(iii) The product over sheets.}  For a separable field extension,
the symmetric product over all embeddings is the field norm; it needs
no ordering and is invariant under relabelling.  It can also extend
over a normal target.  The limitation is a different one: counting
only distinct points of an arbitrary fibre omits inseparability and
Artin lengths.  At a nonreduced fibre the underlying points still exist,
but they do not record the full algebra.  The finite locally free
determinant supplies both those multiplicities and arbitrary
base-change compatibility.

The determinant has none of these defects.  It is defined for an arbitrary
finite locally free algebra, of arbitrary rank, in arbitrary
characteristic, with arbitrary fibres; it is multiplicative rather than
additive, so it converts the ultrametric \emph{inequality} for sums into an
ultrametric \emph{identity} for products; and, as a determinant, it is
manifestly independent of any choice of basis or ordering.  The price is
that it is not a linear map but a homogeneous polynomial law of degree
$d$, and that its interaction with a \emph{finite maximum} of weighted
sections has to be organized: this is the content of
\Cref{prop:bj-weighted-coefficient-norm}, and it is why the universal
variables $T_j$ are introduced.

\begin{remark}[The one place division by $d$ occurs]
\label{rem:ch41-division-by-d}
The number $d$ appears twice in \eqref{eq:ch41-average-shape} and in
\eqref{eq:bj-weighted-average-definition}: once as a real normalizing factor
$1/d$ in front of a real-valued weight, and once as a rational
normalizing factor $1/d$ in front of a numerical class.  Neither is a
division inside $k$, inside $\cO_Z$, or inside the finite algebra.  In
particular nothing in this chapter degenerates when
$\operatorname{char}k\mid d$, which is the situation of
\Cref{ex:ch41-frobenius-average}.  A reader who has met the classical
"average over the fibre'' construction in characteristic zero should keep
this distinction in view; it is the whole difference between an argument
that works and an argument that works only for $p\nmid d$.
\end{remark}

\subsection{The firewall}
\label{sec:ch41-firewall}

This chapter is one of the inputs of the envelope theorem, and the danger of
circularity is concrete rather than theoretical: \Cref{chap:bj-alterations}
applies \Cref{thm:bj-finite-normal-descent} to reduce the normal case to the
smooth case, and \Cref{chap:bj-envelope-statement} lists it in its proof map.
A single appeal here to the envelope property as an established fact would
close a loop through \Crefrange{chap:bj-cartier-realization}{chap:bj-envelope-consequences}.

\begin{warning}{The envelope theorem is not available in this chapter}
\label{warn:ch41-firewall}
Nothing in \Cref{chap:bj-finite-descent} may use, directly or transitively:
\begin{enumerate}
\item \Cref{thm:bj-final-envelope}, in any of its formulations, nor
      \Cref{def:bj-envelope-property-final};
\item the envelope property $\operatorname{EP}$ of
      \Cref{def:bj-envelope-property} as a \emph{theorem} for any class on
      any variety --- with the single exception of
      \Cref{thm:bj-smooth-envelope} on a smooth source, and then only
      inside illustrative examples such as \Cref{ex:ch41-descent-cone};
      the descent proof otherwise consumes EP only as a displayed
      hypothesis on the source; EP on the target is its conclusion;
\item closure of $\PSH_X(\theta)$ or of $\PSH_T(h^*\theta)$ under
      suprema, regularized or not;
\item the obstacle envelope $P_\theta$ of
      \Cref{def:bj-obstacle-envelope} regarded as an operator with values in
      a psh class;
\item \Cref{thm:ch35-cartier-psh-criterion} in its direction
      $\Leftarrow$, and its consequences
      \Cref{thm:bj-nef-bdivisor-correspondence} and
      \Cref{cor:ch35-monotone-approximation};
\item any statement of \Cref{chap:bj-envelope-consequences}.
\end{enumerate}
Every limit taken below is a decreasing limit.  Every supremum
\emph{constructed} below is a supremum over a \emph{finite} index set ---
the finitely many multi-indices $I$ with $|I|=d$, or the finitely many
weights $\lambda_j$ of a Fubini--Study presentation.  The two regularized family suprema consumed by the descent proof
are both handled under the assumed
$\operatorname{EP}_T(h^*\theta)$: the given family of
\Cref{def:bj-envelope-property}, pulled back in Step~0 of the proof of
\Cref{thm:bj-finite-normal-descent}, and the canonical carrier
\eqref{eq:bj-canonical-carrier}, invoked through
\Cref{lem:bj-canonical-carrier}.
\end{warning}

\begin{remark}[Why the firewall is not vacuous here]
\label{rem:ch41-firewall-not-vacuous}
The conclusion is precisely the envelope property downstairs, so
invoking it in advance would bypass the work.  Step~0 instead uses only
the assumed envelope property on $T$ and analytic openness.  Step~2
uses that same source hypothesis to obtain a nonempty canonical carrier.
Step~3 applies the already established carrier reflection and identifies
functions by analytic surjectivity.  It does not invoke the later
envelope theorem or a divisorial comparison with an uncertified psh
function.
\end{remark}

\subsection{The six tasks, and where each is discharged}
\label{sec:ch41-tasks}

\begin{enumerate}
\item\label{itm:ch41-task:open}
      \emph{Analytic openness.}  A finite dominant morphism of normal
      integral projective varieties induces an open surjection on
      analytifications, including at points of nonzero support.  This is
      \Cref{lem:bj-finite-normal-open}, and its corollary
      \eqref{eq:bj-finite-open-usc} is what allows an upper-semicontinuous
      regularization computed downstairs to be recognized upstairs.  Proved
      in \Cref{sec:ch41-openness}.
\item\label{itm:ch41-task:flat}
      \emph{Flattening.}  A finite morphism need not be flat, and the
      determinant norm needs a locally free algebra.  Raynaud--Gruson
      flattening replaces $h$ by a finite \emph{flat} morphism $q$ after
      birational changes of both source and target.  This is
      \Cref{lem:bj-finite-flatification}, proved in
      \Cref{sec:ch41-flatification}; the failure of flatness is genuine and
      is computed in \Cref{ex:ch41-nonflat-cone}.
\item\label{itm:ch41-task:norm}
      \emph{Determinant norms.}  For a finite locally free algebra one
      constructs the norm of an element, the norm of a line bundle, and the
      norm of a section; these are multiplicative and compatible with
      arbitrary base change.  \Cref{sec:ch41-determinant}.
\item\label{itm:ch41-task:average}
      \emph{Weighted averaging.}  The coefficients of the universal norm
      polynomial assemble into a Fubini--Study function downstairs whose
      value at every analytic point is the multiplicity-weighted fibre
      average \eqref{eq:ch41-average-shape}.  No separability, no
      reducedness, no bound on the characteristic.
      \Cref{prop:bj-weighted-coefficient-norm}, in
      \Cref{sec:ch41-averaging}.
\item\label{itm:ch41-task:descent}
      \emph{Descent.}  Combining the four preceding items with birational
      carrier reflection from \Cref{chap:bj-reflection} gives
      \Cref{thm:bj-finite-normal-descent}, in \Cref{sec:ch41-descent}.
\item\label{itm:ch41-task:law}
      \emph{Universal polynomial laws and base change.}  The norm is not a
      map of modules; it is a homogeneous polynomial law of degree $d$ in
      the sense of \Cref{def:ch41-polynomial-law}.  Its base-change
      compatibility, \Cref{prop:ch41-norm-base-change}, is the exact
      statement that permits the fibrewise computation of
      \Cref{sec:ch41-averaging} to be performed after an isometric field
      extension and then transported back.  This is the item singled out for
      audit in the plan under the heading ``universal properties of
      determinant norms''.
\end{enumerate}

\subsection{Standing hypotheses and notation}
\label{sec:ch41-hypotheses}

Throughout the chapter $k$ is an algebraically closed field, of arbitrary
characteristic, equipped with the trivial absolute value, and all schemes
are of finite type over $k$.  The standing geometric datum is
\[
 \begin{aligned}
  h&:T\longrightarrow X &&\text{finite surjective},\\
  T,X& &&\text{normal, integral, and projective},\\
  d&:=[K(T):K(X)],
 \end{aligned}
\tag{41.2}\label{eq:ch41-standing}
\]
together with a class $\theta\in N^1(X)_{\mathbf R}$.  Since $T$ and
$X$ are integral and $h$ is finite and dominant, $d$ is finite and
positive, and $h$ is surjective as a map of schemes.

Analytifications are those of \Cref{def:trivial-analytification}: a point of
$X^{\An}$ is given on affine charts by a semivaluation with values
in $\mathbf R\cup\{+\infty\}$, trivial on $k$; it is nonnegative on
the local ring of its centre.
Equivalently, by \Cref{lem:ch41-point-pair}, it is a pair consisting of a scheme
point $\xi\in X$ and a real valuation of $\kappa(\xi)$ trivial on $k$.
We write $s(v)$ for the support, $c(v)$ for the centre,
$X^{\divis}$ for the divisorial points and $v_{Y,\triv}$ for the trivial
semivaluation of an integral closed subvariety $Y\subseteq X$, as in
\Cref{chap:trivial-analytification}.  For a section $s$ of a line bundle,
$\log|s|$ is the function \eqref{eq:section-log-value}, with values in
$[-\infty,0]$.

The Fubini--Study spaces $\cH_\Lambda(L)$ and
$\cH^{\mathrm{gf}}_\Lambda(L)$ are those of
\Cref{def:gf-fs-functions}, and $\PSH_X(\theta)$ is
\Cref{def:bj-psh}, in its corrected form in which the approximating
classes are allowed to vary.  We record once and for all that this
correction is used below: the numerical classes of the averaged presentations
constructed in \Cref{prop:bj-weighted-coefficient-norm} converge to
$\beta^*\theta$, without being required to equal it at finite stages, and
\Cref{thm:bj-basic-closure}\,(5) is the statement that permits the passage
to the limit.  Neither $\theta$ nor $h^*\theta$ is assumed nef, big, or
even pseudoeffective; by \Cref{lem:ch34-psh-nonempty} the statements are
vacuous unless $\theta$ is pseudoeffective, and nefness would be needed
only to guarantee \emph{bounded} psh functions, which we never require.

Finally, $\operatorname{EP}_X(\theta)$ always abbreviates the envelope
property of \Cref{def:bj-envelope-property} for the pair $(X,\theta)$.

\subsection{Dependencies}
\label{sec:ch41-dependencies}

From \Cref{chap:trivial-analytification} we use the description of
$X^{\An}$, the support and centre maps
(\Cref{thm:ch32-support-center-maps}), compactness
(\Cref{thm:triv-berkovich-compactness}), the scaling action
(\Cref{prop:ch32-scaling-action}) and the maximum principle
(\Cref{prop:ch32-maximum-principle}).  From
\Cref{chap:bj-fs-pl-psh} we use \Cref{def:gf-fs-functions},
\Cref{prop:ch34-generic-finiteness}, \Cref{def:bj-psh},
\Cref{lem:ch34-psh-nonempty} and \Cref{thm:bj-basic-closure}.  From
\Cref{chap:bj-homogeneity-bdivisors} we use the scaling calculus
\Cref{lem:ch35-action-basic}, the homogeneous psh space
\eqref{eq:homogeneous-psh-space}, the normalization
\Cref{prop:ch35-homogeneous-normalization} and the homogeneous
regularization \Cref{prop:homogenization-gf}.  From
\Cref{chap:bj-reflection} we use exactly two statements:
\Cref{thm:bj-carrier-reflection} for carrier reflection, and
\Cref{lem:bj-canonical-carrier} for the canonical carrier under an assumed
envelope property.  We do \emph{not} use \Cref{thm:ch39-pullback}, whose
standing hypotheses require a normal source; the birational model $W$
produced by flattening is only integral, and the elementary substitute
\Cref{lem:ch41-proper-gf-pullback} is proved below.

The algebraic inputs are collected in the appendices and are not reproved
here.  \Cref{sec:appA-determinant-norms} contains the elementary identities
for norms of elements, \Cref{sec:appB-flattening} contains the
Raynaud--Gruson theorem \Cref{thm:appB-flattening} together with its
specialization \Cref{cor:appB-flatify-finite-algebra}, and
\Cref{app:stacks} audits every Stacks Project tag used below: 035S, 0BRI,
0BRJ, 0815, 02K9, 01OF, 0BD2 and 0CNF are all discussed there with the exact
statement imported and the hypotheses to be checked.  The reader who wants
to verify that no tag is being asked to prove more than it does should read
\Cref{sec:stacks-flattening} and \Cref{sec:stacks-finite} beside
\Cref{sec:ch41-flatification} and \Cref{sec:ch41-determinant}, and the
summary chain in \Cref{sec:stacks-dependency-chains}.

\section{Semivaluations under a finite morphism}
\label{sec:ch41-finite-analytic}

\subsection{Points, supports and centres}
\label{sec:ch41-points}

The first thing to fix is a usable description of a point of $X^{\An}$.
\Cref{def:trivial-analytification} presents such a point as a multiplicative
seminorm on the structure sheaf; for the computations of this chapter the
following equivalent presentation, as a scheme point together with a genuine
valuation of its residue field, is the one to use.  It is nothing more than
a change of packaging, but the packaging matters: the naive habit of
identifying $X^{\An}$ with the space of valuations of the function field
$K(X)$ discards all proper-support points.  This would leave the
all-point openness and fibre statements unproved, regardless of where
nonreduced fibres occur.

\begin{lemma}
\label{lem:ch41-point-pair}
Let $X$ be a $k$-variety.  The assignment
\[
 v\longmapsto
 \bigl(\xi_v,\ |\cdot|_v\bigr),
 \qquad
 \xi_v:=\text{the generic point of }s(v),
 \quad
 |f|_v:=\rme^{-v(f)},
\tag{41.3}\label{eq:ch41-point-pair}
\]
is a bijection from $X^{\An}$ onto the set of pairs $(\xi,|\cdot|)$ in
which $\xi\in X$ is a scheme point and $|\cdot|$ is a multiplicative
non-Archimedean absolute value on the residue field $\kappa(\xi)$ that is
trivial on $k$.  Under \eqref{eq:ch41-point-pair} the support of $v$ is
$\overline{\{\xi_v\}}$; the point $v$ lies in $U^{\An}$ for a Zariski
open $U\subseteq X$ if and only if $\xi_v\in U$; and $v$ is the
trivial semivaluation $v_{Y,\triv}$ of an integral closed subvariety
$Y\subseteq X$ if and only if $\xi_v$ is the generic point of $Y$ and
$|\cdot|_v$ is trivial.
\end{lemma}

\begin{proof}
On an affine chart $U=\Spec A$, an analytic point is a multiplicative
seminorm, or equivalently a semivaluation
$v:A\to\mathbf R\cup\{+\infty\}$ trivial on $k$.
Its kernel $\mathfrak p_v=\{f:v(f)=+\infty\}$ is prime.
The induced valuation on $A/\mathfrak p_v$ extends to
$\operatorname{Frac}(A/\mathfrak p_v)=\kappa(\xi_v)$.
Conversely an absolute value on $\kappa(\xi)$ gives the seminorm
$f\mapsto|f(\xi)|$ on every affine neighbourhood of $\xi$.
These constructions are inverse and agree on overlaps.  Nonnegativity
is a centre condition; it is not imposed on every affine chart.

The statement about supports is the definition of $\xi_v$.  The statement
about $U^{\An}$ is
\Cref{thm:ch32-support-center-maps}\,\cref{itm:ch32-scm:s}.  The last
statement is immediate from the description: the trivial absolute value on
$\kappa(\xi)$ corresponds to the semivaluation taking only the values
$0$ and $+\infty$, which is exactly $v_{Y,\triv}$ for
$Y=\overline{\{\xi\}}$.
\end{proof}

\begin{remark}[Two loci, and why both are needed]
\label{rem:ch41-two-loci}
\Cref{lem:ch41-point-pair} exhibits $X^{\An}$ as fibred over the
underlying scheme $X$ by the support map, the fibre over $\xi$ being the
space of $k$-trivial absolute values on $\kappa(\xi)$.  The fibre over
the generic point is the space of valuations of $K(X)$, which is what a
purely birational treatment would retain; every other fibre lies in the
proper-support locus of \Cref{def:ch32-proper-support}.  Both matter here.
Proper-support points may be poles of psh functions and may have
nonreduced fibre algebras.  These phenomena also have broader scopes,
as noted in \Cref{sec:ch41-what}.  The proper-support locus is dense
under the dimension hypotheses of
\Cref{prop:ch32-proper-support-dense}; it is not dense on a projective
curve.  In all dimensions, a claim established only on $X^{\val}$
does not by itself establish the required neighbourhood statement at
every analytic point.
\end{remark}

The next two lemmas record how the two structural maps of
\Cref{thm:ch32-support-center-maps} interact with a finite morphism.  They
are used constantly and are proved once here.

\begin{lemma}
\label{lem:ch41-support-centre-finite}
Let $g:W\to Z$ be a finite morphism of projective $k$-varieties and let
$w\in W^{\An}$, $z:=g^{\An}(w)$.  Then, in the notation of
\Cref{lem:ch41-point-pair}:
\begin{enumerate}
\item\label{itm:ch41-sc:supp}
      $\xi_z=g(\xi_w)$, and $|\cdot|_z$ is the restriction of
      $|\cdot|_w$ along the field extension
      $\kappa(\xi_z)\hookrightarrow\kappa(\xi_w)$, which is finite;
\item\label{itm:ch41-sc:centre}
      $c(z)=g(c(w))$;
\item\label{itm:ch41-sc:sections}
      for every line bundle $M$ on $W$ and every
      $s\in H^0(W,M)$, the value $\log|s|(w)$ may be computed with a
      frame of $M$ defined on a neighbourhood of $g^{-1}(c(z))$.
\end{enumerate}
\end{lemma}

\begin{proof}
\cref{itm:ch41-sc:supp}  By construction $g^{\An}(w)$ is the semivaluation
$f\mapsto w(g^\sharp f)$, whose support prime is the contraction of the
support prime of $w$; this is the assertion $\xi_z=g(\xi_w)$, and the
description of $|\cdot|_z$ as the restriction of $|\cdot|_w$ is the same
statement read through \eqref{eq:ch41-point-pair}.  The extension
$\kappa(\xi_z)\subseteq\kappa(\xi_w)$ is finite because $g$ is finite:
the fibre of $g$ over $\xi_z$ is the spectrum of a finite
$\kappa(\xi_z)$-algebra, of which $\kappa(\xi_w)$ is a residue field.

\cref{itm:ch41-sc:centre}  Let $R_w$ be the valuation ring in
$\kappa(\xi_w)$.  The centre is the image of the closed point under
the unique extension $\Spec R_w\to W$ of $\Spec\kappa(\xi_w)\to W$.
Composing with $g$ gives a map to $Z$.  Put
$R_z=R_w\cap\kappa(\xi_z)$.  The inclusion $R_z\subset R_w$ is
local.  Composing $\Spec R_w\to\Spec R_z$ with the centre map to $Z$
has the same generic-point map as the preceding composite, so
separatedness identifies them.  Their closed-point images give
$c(z)=g(c(w))$.  This argument uses maps of local rings; it does not
inject a local ring with a nonzero support kernel into a field.

\cref{itm:ch41-sc:sections}  By \cref{itm:ch41-sc:centre} the centre $c(w)$
lies in $g^{-1}(c(z))$, and by \Cref{lem:ch34-frame-independence} the
value $\log|s|(w)$ is computed from any frame of $M$ near $c(w)$.
A frame on a neighbourhood of the whole finite fibre can be chosen:
$M$ is free over the finite algebra over the local ring at $c(z)$,
which is semilocal, and that generator spreads to a neighbourhood of
the fibre.  Closedness of $g$ allows the base to be shrunk so that its
full preimage lies in this neighbourhood.  This uses finiteness,
not flatness.
\end{proof}

\begin{lemma}
\label{lem:ch41-finite-surjective}
Let $g:W\to Z$ be a finite surjective morphism of projective
$k$-varieties.  Then $g^{\An}:W^{\An}\to Z^{\An}$ is surjective, closed,
and proper; in particular it maps compact sets to compact sets and its
fibres are compact.
\end{lemma}

\begin{proof}
For surjectivity fix $z\in Z^{\An}$ and write $z=(\xi,|\cdot|)$ as in
\Cref{lem:ch41-point-pair}.  Since $g$ is surjective and closed, there is
a point $\eta\in W$ with $g(\eta)=\xi$; indeed the fibre
$W_\xi=W\times_Z\Spec\kappa(\xi)$ is a nonempty affine scheme, finite over
$\kappa(\xi)$.  Choose a point $\eta$ of $W_\xi$; then
$\kappa(\eta)/\kappa(\xi)$ is a finite field extension.  It remains to
extend the absolute value $|\cdot|$ from $\kappa(\xi)$ to
$\kappa(\eta)$.  Let $\cO\subseteq\kappa(\xi)$ be its valuation ring and
let $C$ be the integral closure of $\cO$ in $\kappa(\eta)$.  The ring
$C$ is integral over $\cO$, so by lying over there is a prime
$P\subset C$ with $P\cap\cO=\mathfrak m_{\cO}$, and the classical
structure theory of the integral closure of a valuation ring in an
algebraic extension shows that $C_P$ is a valuation ring of
$\kappa(\eta)$ dominating $\cO$; see
\cite[Chapter~VI, \S4 and \S11]{ZariskiSamuelII}.  Its associated absolute
value is trivial on $k$, because it is trivial on $k$ already on
$\kappa(\xi)$ and $k$ is a subfield of $\kappa(\xi)$.  The
resulting pair $(\eta,|\cdot|_{C_P})$ is a point of $W^{\An}$ mapping to
$z$.

Closedness and properness are formal: $W^{\An}$ is compact and
$Z^{\An}$ is Hausdorff by \Cref{thm:triv-berkovich-compactness}, and a
continuous map from a compact space to a Hausdorff space is closed and
proper.
\end{proof}

\begin{remark}[What surjectivity does and does not give]
\label{rem:ch41-surjective-not-open}
\Cref{lem:ch41-finite-surjective} holds for every finite surjective
morphism, with no normality assumption, and it already yields the
inequality $\sup_{W^{\An}}g^*a=\sup_{Z^{\An}}a$ for any function $a$ on
$Z^{\An}$, together with injectivity of $g^*$ on functions.  What it
does \emph{not} give is the interchange
\eqref{eq:bj-finite-open-usc} of pullback with upper-semicontinuous
regularization.  That interchange needs $g^{\An}$ to be \emph{open}, and
openness is where normality of the target enters.
\Cref{ex:ch41-node-not-open} shows a finite surjective morphism onto a
non-normal curve whose analytification is not open, so the hypothesis is not
decorative.
\end{remark}

\subsection{Extensions of an absolute value}
\label{sec:ch41-value-extension}

Openness of $h^{\An}$ is proved by exhibiting $X^{\An}$ as a quotient of
the analytification of a Galois cover.  The group-theoretic input is the
following classical transitivity statement, which we prove from the
prime-transitivity tag audited in \Cref{sec:stacks-finite} so that the
hypotheses are visible.

\begin{lemma}
\label{lem:ch41-value-transitivity}
Let $F\subseteq E$ be a finite normal field extension, let $|\cdot|$ be
a non-Archimedean absolute value on $F$, and let $|\cdot|_1$,
$|\cdot|_2$ be absolute values on $E$ restricting to $|\cdot|$ on
$F$.  Then there is $\sigma\in\Aut(E/F)$ with
$|\cdot|_2=|\cdot|_1\circ\sigma$.  In particular $\Aut(E/F)$ acts
transitively on the set of extensions of $|\cdot|$ to $E$, and that set
is finite.
\end{lemma}

\begin{proof}
Let $E_0:=E^{\Aut(E/F)}$.  Since $E/F$ is normal, $E_0/F$ is purely
inseparable and $E/E_0$ is Galois with group $G:=\Aut(E/F)$; see
\cite[Chapter~II]{ZariskiSamuelII}.

\emph{Step 1: uniqueness over the purely inseparable part.}  If
$E_0/F$ is purely inseparable of exponent $p^r$, then for $b\in E_0$
one has $b^{p^r}\in F$, so any absolute value $\|\cdot\|$ on $E_0$
extending $|\cdot|$ satisfies
\[
 \|b\|=\|b^{p^r}\|^{1/p^r}=|b^{p^r}|^{1/p^r},
\tag{41.4}\label{eq:ch41-inseparable-formula}
\]
which determines $\|\cdot\|$ uniquely.  Hence $|\cdot|_1$ and
$|\cdot|_2$ agree on $E_0$; call the common restriction $\|\cdot\|$,
with valuation ring $\cO_0$.

\emph{Step 2: transitivity over the Galois part.}  Let $C$ be the integral
closure of $\cO_0$ in $E$.  It is stable under $G$, and
$C^G=C\cap E_0=\cO_0$, the last equality because $\cO_0$ is integrally
closed in $E_0$, being a valuation ring.  The valuation rings
$\cO_1,\cO_2\subseteq E$ of $|\cdot|_1,|\cdot|_2$ both dominate
$\cO_0$ and are integrally closed, so each contains $C$ and is a
localization of $C$ at a prime lying over $\mathfrak m_{\cO_0}$; this is
again the classical structure theory of \cite[Chapter~VI,
\S11]{ZariskiSamuelII}, applied to the algebraic extension $E/E_0$.  Write
$\cO_i=C_{P_i}$.  Both $P_1$ and $P_2$ lie over
$\mathfrak m_{\cO_0}=\mathfrak m_{\cO_0}\cap C^G$, so by
\cite[\href{https://stacks.math.columbia.edu/tag/0BRI}{Tag 0BRI}]
{stacks-project}, applied to the action of the finite group $G$ on the
ring $C$ with $C^G=\cO_0$, there is $\sigma\in G$ with
$\sigma(P_1)=P_2$.  Then $\sigma(\cO_1)=\cO_2$, so
$|\cdot|_2=|\cdot|_1\circ\sigma^{-1}$.
The normalization is fixed by the restriction to $F$: if that absolute
value is nontrivial, an equivalent real valuation cannot be rescaled;
if it is trivial, every algebraic extension is also trivial.
Replacing $\sigma$ by its inverse gives the stated formula.

Finiteness of the set of extensions follows, the set being a quotient of the
finite group $G$.
\end{proof}

\begin{lemma}
\label{lem:ch41-radicial-homeo}
Let $g:W\to Z$ be a finite surjective morphism of projective
$k$-varieties which is radicial, that is, universally injective in the
sense of \Cref{def:appA-universal-homeomorphism}.  Then
$g^{\An}:W^{\An}\to Z^{\An}$ is a homeomorphism.  This applies in
particular when $W$ and $Z$ are integral and normal and
$K(W)/K(Z)$ is purely inseparable.
\end{lemma}

\begin{proof}
The map $g^{\An}$ is continuous and surjective by
\Cref{lem:ch41-finite-surjective}, and $W^{\An}$ is compact while
$Z^{\An}$ is Hausdorff, by \Cref{thm:triv-berkovich-compactness}.  It
therefore suffices to prove injectivity.

Let $z\in Z^{\An}$, written as a pair $(\xi,|\cdot|)$.  Since $g$ is
universally injective, the fibre $W_\xi$ has exactly one point $\eta$,
and $\kappa(\eta)/\kappa(\xi)$ is purely inseparable.  If
$\operatorname{char}k=0$ this forces $\kappa(\eta)=\kappa(\xi)$, and the
absolute value extends in exactly one way.  If
$\operatorname{char}k=p>0$, let $p^r$ be the exponent of
$\kappa(\eta)/\kappa(\xi)$; then any extension is given by
\eqref{eq:ch41-inseparable-formula}, hence is unique.  Thus $g^{\An}$ is
injective.

For the final assertion, let $b\in\cO_W(V)$ for an affine open
$V\subseteq W$ with $V=g^{-1}(U)$, $U=\Spec A\subseteq Z$ affine.
If $\operatorname{char}k=0$, a purely inseparable extension is trivial, so
$b\in K(Z)$ is integral over the integrally closed ring $A$, whence
$b\in A$; thus $A\to\cO_W(V)$ is an isomorphism, $g$ is an
isomorphism, and there is nothing to prove.  Assume therefore
$\operatorname{char}k=p>0$.
Since $K(W)/K(Z)$ is purely inseparable of some exponent $p^r$, one has
$b^{p^r}\in K(Z)$; and $b^{p^r}$ is integral over $A$, because $b$
is.  As $Z$ is normal, $A$ is integrally closed in $K(Z)$, whence
$b^{p^r}\in A$.  Since $A\to\cO_W(V)$ is injective, both rings being
domains with the same fraction field up to a purely inseparable extension,
its kernel is zero and in particular locally nilpotent, and
$A[1/p]=0=\cO_W(V)[1/p]$ because $p=0$ in characteristic $p$.
Therefore $g$ satisfies the criterion of
\cite[\href{https://stacks.math.columbia.edu/tag/0CNF}{Tag 0CNF}]
{stacks-project}, recorded as
\Cref{thm:appA-universal-homeomorphism-criterion}, and is a universal
homeomorphism; in particular it is radicial.
\end{proof}

\begin{remark}[Where each hypothesis is used]
\label{rem:ch41-radicial-uses}
The argument for \Cref{lem:ch41-radicial-homeo} uses normality of $Z$
only in the final paragraph, to convert a purely inseparable extension of
function fields into a scheme-theoretic universal homeomorphism.  Without
target normality, equality of function fields need not imply a universal
homeomorphism: the normalization map of a nodal curve is a counterexample.
The normalization of a cuspidal cubic is a universal homeomorphism and
therefore illustrates the positive case.  This distinction is related to
the unibranch condition studied in
\Cref{chap:bj-subintegral-normalization}, and the analytic consequences are
drawn there, not here.
\end{remark}

\subsection{Openness}
\label{sec:ch41-openness}

\begin{lemma}\label{lem:bj-finite-normal-open}
Let
\[
                         h:T\longrightarrow X
\tag{41.5}\label{eq:bj-finite-map}
\]
be a finite dominant morphism between normal integral projective
varieties.  Then
\[
                         h^{\An}:T^{\An}\longrightarrow X^{\An}
\]
is open and surjective.  Consequently, for every bounded-above function
\[
                         a:X^{\An}\longrightarrow
                         \mathbf R\cup\{-\infty\},
\]
one has
\[
                         (h^*a)^*=h^*(a^*).
\tag{41.6}\label{eq:bj-finite-open-usc}
\]
\end{lemma}

\begin{proof}
The proof has three steps.  \emph{Step 1} treats a separable extension by
comparing $X^{\An}$ with the analytification of a Galois cover.
\emph{Step 2} removes the separability hypothesis.  \emph{Step 3} deduces
\eqref{eq:bj-finite-open-usc} from openness by a purely topological
argument.  Surjectivity of $h^{\An}$ is
\Cref{lem:ch41-finite-surjective} and will not be repeated.

\medskip
\noindent\emph{Step 1: the separable case.}
Assume first that $K(T)/K(X)$ is separable.  Choose a finite Galois
closure $M/K(X)$, let $S\to X$ be the normalization in $M$, and put
\[
                         G:=\operatorname{Gal}(M/K(X)).
\]
The integral closure in the finite extension $M/K(X)$ is finite
because finite-type schemes over a field are Nagata
(\Cref{sec:stacks-flattening}); Tag 035S records the corresponding
finiteness statement for normalization in the function field.
Thus $S$ is a normal integral projective variety, finite
and surjective over $X$, and the tautological action of $G$ on $M$
extends uniquely to an action on $S$ by $X$-automorphisms, by
functoriality of the relative normalization.
Normality of $X$ gives $X=S/G$: on a
$G$-stable affine chart corresponding to normal domains $A\subset B$,
the ring $B^G$ is integral over $A$ and lies in
$M^G=K(X)$, hence $B^G=A$.
Such charts exist and cover $S$: take $B=\cO_S(g^{-1}(U))$ for
$U=\Spec A\subseteq X$ affine open and $g:S\to X$ the structure
morphism, which is affine because it is finite.

We must verify the analytic quotient on all points, including points with
nonzero support kernel.  Write an analytic point as
\[
 x=(\mathfrak p,|\cdot|_x),\qquad
 |\cdot|_x\text{ an absolute value on }\kappa(\mathfrak p).
\tag{41.7}\label{eq:bj-analytic-point-pair}
\]
A lift to $S^{\An}$ is a pair
\[
 y=(\mathfrak q,|\cdot|_y),\qquad
 \mathfrak q\cap A=\mathfrak p,\qquad
 |\cdot|_y|_{\kappa(\mathfrak p)}=|\cdot|_x.
\tag{41.8}\label{eq:bj-analytic-lift-pair}
\]
We claim that $G$ acts transitively on the set of such lifts.  There are
two things to prove: transitivity on the support primes, and, for a fixed
support prime, transitivity on the extensions of the absolute value.

The primes of $B$ above $\mathfrak p$ form one $G$-orbit by
\cite[\href{https://stacks.math.columbia.edu/tag/0BRI}{Tag 0BRI}]
{stacks-project}, applied to the action of the finite group $G$ on the
ring $B$ with $B^G=A$; this is the tag whose hypotheses are checked in
\Cref{sec:stacks-finite}.  We may therefore assume that the two lifts
under comparison have the same support prime $\mathfrak q$.  Let
\[
 D:=\{\sigma\in G\mid\sigma(\mathfrak q)=\mathfrak q\}
\]
be the decomposition group.  For a fixed $\mathfrak q$, the residue
extension $\kappa(\mathfrak q)/\kappa(\mathfrak p)$ is finite and normal,
and the decomposition group
surjects onto the automorphism group of the normal residue extension by
\cite[\href{https://stacks.math.columbia.edu/tag/0BRJ}{Tag 0BRJ}]
{stacks-project}; this group acts transitively on extensions of the absolute
value through the normal part, by \Cref{lem:ch41-value-transitivity} applied
to the finite normal extension
$\kappa(\mathfrak p)\subseteq\kappa(\mathfrak q)$. The purely inseparable
residue part has a unique extension,
characterized by
\[
                              |b|=|b^{p^r}|^{1/p^r}.
\tag{41.9}\label{eq:bj-purely-inseparable-value}
\]
Combining the two transitivity statements: given lifts $y_1,y_2$ of $x$,
first move $y_1$ by an element of $G$ so that its support prime becomes
that of $y_2$, and then move it further by an element of $D$, lifting an
element of $\Aut(\kappa(\mathfrak q)/\kappa(\mathfrak p))$ supplied by
\Cref{lem:ch41-value-transitivity}, so that the two absolute values agree.
Thus every fibre of $S^{\An}\to X^{\An}$ is a single $G$-orbit.
This argument begins with $\kappa(\mathfrak p)$, not with $K(X)$, and
therefore does not omit proper-support semivaluations.

The natural continuous bijection
\[
                             S^{\An}/G\longrightarrow X^{\An}
\tag{41.10}\label{eq:bj-galois-analytic-quotient}
\]
is a homeomorphism because both spaces are compact Hausdorff.  Let us spell
this out.  The map $S^{\An}\to X^{\An}$ is continuous, surjective by
\Cref{lem:ch41-finite-surjective}, and $G$-invariant, so it factors
through the orbit space; the factorization is injective by the previous
paragraph and continuous for the quotient topology.  The orbit space of a
compact Hausdorff space by a finite group acting through homeomorphisms is
again compact Hausdorff --- compactness because it is a continuous image,
Hausdorffness because the orbits are finite and may be separated by
intersecting finitely many separating neighbourhoods and saturating.  A
continuous bijection from a compact space to a Hausdorff space is a
homeomorphism.  The quotient
map is open: the saturation of an open set $V\subset S^{\An}$ is
\[
                              \bigcup_{\sigma\in G}\sigma V,
\]
a finite union of open sets, each $\sigma$ acting as a homeomorphism of
$S^{\An}$ by functoriality of analytification.  Composing, the finite
morphism $S\to X$ has open analytification.

There is a finite surjection $\pi:S\to T$: indeed $S$ is normal with
function field $M\supseteq K(T)$ and is finite over $X$, hence finite
over $T$, and integrally closed in $M$, so $S$ is the normalization of
$T$ in $M$; it is surjective because it is finite and dominant.  For
every open
$V\subset T^{\An}$,
\[
 h^{\An}(V)
   =(S^{\An}\to X^{\An})\bigl((\pi^{\An})^{-1}V\bigr),
\tag{41.11}\label{eq:bj-galois-image-formula}
\]
which is open.  The inclusion $\supseteq$ in
\eqref{eq:bj-galois-image-formula} holds because $h\circ\pi$ is the
structure morphism $S\to X$; the inclusion $\subseteq$ holds because
$\pi^{\An}$ is surjective by \Cref{lem:ch41-finite-surjective}, so every
$y\in V$ is $\pi^{\An}(s)$ for some $s$, necessarily in
$(\pi^{\An})^{-1}V$.  Notice that this argument does not claim that
$\pi^{\An}$ itself is open; that would be a statement about the pair
$(S,T)$, which happens to be true here by the same lemma applied to
$\pi$, but is not needed and is not used.

\medskip
\noindent\emph{Step 2: the inseparable factor.}
For a general finite extension, take the maximal separable intermediate
field $K_s$, so that $K(X)\subseteq K_s\subseteq K(T)$ with
$K_s/K(X)$ separable and $K(T)/K_s$ purely inseparable.  Let $X_s$ be
the normalization of $X$ in $K_s$.  The same Nagata property makes
this relative normalization finite over $X$, hence a normal integral
projective variety, and $h$ factors as
\[
 T\xrightarrow{\ h_2\ }X_s\xrightarrow{\ h_1\ }X .
\]
Both factors are finite and surjective, and $T$ is normal with function
field $K(T)$, integrally closed in it, so $T$ is the normalization of
$X_s$ in $K(T)$.
The remaining finite purely inseparable normalization map $h_2$ is
radicial and surjective, hence a universal homeomorphism, by the last
assertion of \Cref{lem:ch41-radicial-homeo}.  It induces a
homeomorphism on analytifications: each support prime has a unique lift,
and the residue absolute value extends uniquely by
\eqref{eq:bj-purely-inseparable-value}.  This is exactly the first
assertion of \Cref{lem:ch41-radicial-homeo}.  Composing this homeomorphism
with the separable case, applied to $h_1$, proves openness in arbitrary
characteristic: $h^{\An}=h_1^{\An}\circ h_2^{\An}$, a homeomorphism
followed by an open map.

\medskip
\noindent\emph{Step 3: regularization.}
Let $y\in T^{\An}$ and $x=h^{\An}(y)$.  Both regularizations are
computed by the neighbourhood formula \eqref{eq:appE-usc-regularization}.
For the inequality $\ge$, let $V\subseteq T^{\An}$ be an open
neighbourhood of $y$.  Then $h^{\An}(V)$ is an open neighbourhood of
$x$, by Steps~1 and~2, and
$\sup_V h^*a=\sup_{h^{\An}(V)}a$ because the function $a\circ h^{\An}$
takes on $V$ exactly the values that $a$ takes on $h^{\An}(V)$.
Hence $\sup_Vh^*a\ge\inf_{U\ni x}\sup_Ua=a^*(x)$, and taking the infimum
over $V$ gives $(h^*a)^*(y)\ge a^*(x)$.  For the inequality $\le$,
let $U\subseteq X^{\An}$ be an open neighbourhood of $x$ and put
$V:=(h^{\An})^{-1}(U)$.  This is an open neighbourhood of $y$ because
$h^{\An}:T^{\An}\to X^{\An}$ is continuous;
then $\sup_Vh^*a\le\sup_Ua$, and taking infima gives
$(h^*a)^*(y)\le a^*(x)$.  Therefore
\[
 \begin{aligned}
 (h^*a)^*(y)
   &=\inf_{V\ni y}\sup_Vh^*a\\
   &=\inf_{U\ni x}\sup_Ua
    =a^*(x).
 \end{aligned}
\]
The computation remains valid when the common value is $-\infty$: the
suprema are taken in $\Rbar$, which is a complete lattice, and no
subtraction occurs.  Since $y$ was arbitrary and $h^{\An}$ is
surjective, this is \eqref{eq:bj-finite-open-usc}.
\end{proof}

\begin{remark}[Openness is a statement about all of $T^{\An}$]
\label{rem:ch41-openness-scope}
Step~1 checks both support primes and residue absolute values,
so it proves the quotient description at every analytic point.
A proof confined to valuations of the function field would cover only
full support and would not establish openness at the remaining points.

Normality of the source is unnecessary for this analytic assertion.
For a finite dominant map from an integral projective $T$ to a normal
$X$, normalize $T$ to $\widetilde T$; this is finite.  The composite
$\widetilde T\to X$ has open analytification by the proved case.
If $V\subset T^{\An}$ is open, its image in $X^{\An}$ equals the
image of the open inverse image of $V$ in $\widetilde T^{\An}$,
because normalization is analytically surjective.  Thus the image is
open.  The descent theorem below retains normality of $T$ for its
source psh and carrier inputs.
\end{remark}

\begin{warning}{Target normality in the openness argument}
\label{warn:ch41-openness-normality}
\Cref{lem:bj-finite-normal-open} is false without normality.  Target normality enters in the quotient and radicial arguments: in Step~1, to know $B^G=A$, which is what makes
$X^{\An}$ the orbit space rather than merely a continuous image; and in
Step~2, to know that the purely inseparable part is a universal
homeomorphism.  \Cref{ex:ch41-node-not-open} gives a finite surjective
morphism onto a nodal curve, from a normal source, whose analytification is
not open, and \Cref{ex:ch41-node-not-open} also shows that
\eqref{eq:bj-finite-open-usc} genuinely fails there.  This is the reason
\Cref{chap:bj-alterations} arranges the Stein factorization so that the
finite step has \emph{normal} target, and the reason the non-normal case is
postponed to
\Cref{chap:bj-universal-homeomorphisms,chap:bj-subintegral-normalization},
where it is handled by an entirely different mechanism.
\end{warning}

\subsection{Pullback along a finite morphism}
\label{sec:ch41-finite-pullback}

Openness is one half of what the analytic map has to supply.  The other
half is that pullback along it is compatible with the Fubini--Study and psh
formalisms; this is the analogue for finite morphisms of
\Cref{thm:ch39-pullback}, and it is proved the same way.

\begin{proposition}
\label{prop:ch41-finite-pullback}
Let $g:W\to Z$ be a finite surjective morphism of integral projective
$k$-varieties and let $\eta\in N^1(Z)_{\mathbf R}$.  Write
$g^*u:=u\circ g^{\An}$.  Then:
\begin{enumerate}
\item\label{itm:ch41-pb:fs}
      $g^*$ maps $\cH^{\mathrm{gf}}_\Lambda(L)$ into
      $\cH^{\mathrm{gf}}_\Lambda(g^*L)$ and $\cH_\Lambda(L)$ into
      $\cH_\Lambda(g^*L)$, for every $\mathbf Q$-line bundle $L$ on
      $Z$ and every additive subgroup
      $\Lambda\subseteq\mathbf R$; explicitly, if
      $\phi=\frac1m\max_j\{\log|s_j|+\lambda_j\}$ then
      $g^*\phi=\frac1m\max_j\{\log|g^*s_j|+\lambda_j\}$;
\item\label{itm:ch41-pb:psh}
      $g^*$ maps $\PSH_Z(\eta)$ into $\PSH_W(g^*\eta)$, and
      homogeneous functions to homogeneous functions;
\item\label{itm:ch41-pb:order}
      $g^*$ is injective and reflects the order: for arbitrary functions
      $u_1,u_2$ on $Z^{\An}$, one has $g^*u_1\le g^*u_2$ if and only
      if $u_1\le u_2$;
\item\label{itm:ch41-pb:ops}
      $g^*$ commutes with the addition of real constants, with arbitrary
      suprema and infima, and with the scaling action, and it preserves
      upper semicontinuity;
\item\label{itm:ch41-pb:sup}
      $\sup_{W^{\An}}g^*u=\sup_{Z^{\An}}u$, and
      $g^*u(v_{W,\triv})=u(v_{Z,\triv})$.
\end{enumerate}
\end{proposition}

\begin{proof}
\cref{itm:ch41-pb:fs}  Let
$\phi=\frac1m\max_j\{\log|s_j|+\lambda_j\}$ with
$0\ne s_j\in H^0(Z,mL)$.  Then $g^*s_j\in H^0(W,mg^*L)$ is nonzero:
$g$ is dominant and $W$ is integral, so the generic point of $W$ maps
to the generic point of $Z$, where $s_j$ does not vanish.  For
$w\in W^{\An}$ with $z=g^{\An}(w)$, choose by
\Cref{lem:ch41-support-centre-finite}\,\cref{itm:ch41-sc:centre} a frame
$\tau$ of $mL$ near $c(z)$; then $g^*\tau$ is a frame of
$mg^*L$ near $c(w)$, and writing $s_j=f_j\tau$ we get
$g^*s_j=(f_j\circ g)g^*\tau$, whence
$\log|g^*s_j|(w)=-w(f_j\circ g)=-z(f_j)=\log|s_j|(z)$ by
\Cref{lem:ch41-support-centre-finite}\,\cref{itm:ch41-sc:supp}.  Taking the
weighted maximum gives the displayed formula.  Generic finiteness is
preserved because $W$ is integral and some $g^*s_j$ is nonzero; and if
the $s_j$ have no common zero on $Z$, the $g^*s_j$ have none on
$W$, because $g$ is surjective.

\cref{itm:ch41-pb:psh}  Let $u\in\PSH_Z(\eta)$ with approximating net
$u_i\in\cH^{\mathrm{gf}}_{\mathbf R}(L_i)$, $c_1(L_i)\to\eta$,
$u_i\searrow u$, as in \Cref{def:bj-psh}.  By \cref{itm:ch41-pb:fs} the
functions $g^*u_i$ lie in $\cH^{\mathrm{gf}}_{\mathbf R}(g^*L_i)$; the
net $(g^*u_i)$ is decreasing and converges pointwise to $g^*u$; and
$g^*c_1(L_i)\to g^*\eta$ because
$g^*:N^1(Z)_{\mathbf R}\to N^1(W)_{\mathbf R}$ is linear on a
finite-dimensional space.  Nontriviality of $g^*u$ follows directly from
$g^*u(v_{W,\triv})=u(v_{Z,\triv})>-\infty$, by
\cref{itm:ch41-pb:sup}.  Hence
\Cref{def:bj-psh} applies with the class $g^*\eta$.  Homogeneity is
preserved by \cref{itm:ch41-pb:ops}.

\cref{itm:ch41-pb:order}  Immediate from surjectivity of $g^{\An}$
(\Cref{lem:ch41-finite-surjective}): for $z\in Z^{\An}$ choose $w$ with
$g^{\An}(w)=z$ and evaluate.

\cref{itm:ch41-pb:ops}  Composition with a fixed map commutes with all
pointwise lattice operations and with the addition of constants, and
preserves upper semicontinuity because $g^{\An}$ is continuous.  For the
scaling action, note that $g^{\An}(tw)=t\,g^{\An}(w)$ for $t>0$, which
is immediate from \eqref{eq:valuation-scaling} and the definition of
$g^{\An}$; then
$g^*(t\cdot u)(w)=t\,u(t^{-1}g^{\An}w)=t\,g^*u(t^{-1}w)
=(t\cdot g^*u)(w)$.

\cref{itm:ch41-pb:sup}  The first identity is surjectivity of $g^{\An}$.
For the second, $g^{\An}(v_{W,\triv})=v_{Z,\triv}$ by
\Cref{lem:ch41-support-centre-finite}\,\cref{itm:ch41-sc:supp}: the generic
point of $W$ maps to the generic point of $Z$, and the trivial absolute
value restricts to the trivial absolute value.
\end{proof}

\begin{remark}[Pullback needs no normality and no openness]
\label{rem:ch41-pullback-cheap}
\Cref{prop:ch41-finite-pullback} uses only that $g$ is finite, surjective,
and a morphism of integral projective varieties.  Neither normality nor
\Cref{lem:bj-finite-normal-open} is invoked.  This is the precise sense in
which ``going up is free'' in \Cref{sec:ch41-what}, and it is worth keeping
the two facts apart: the openness lemma is used only to identify a
regularized supremum, never to move a psh function upwards.
\end{remark}

\subsection{Two worked examples}
\label{sec:ch41-openness-examples}

\begin{example}[A degree-two cover of the line]
\label{ex:ch41-openness-p1}
Let $\operatorname{char}k\ne2$, let $T=X=\mathbf P^1$ and let
$h(u)=u^2=x$.  Then $h$ is finite of degree $2$ between smooth
projective curves, and $K(T)/K(X)=k(u)/k(u^2)$ is separable, so Step~1 of
\Cref{lem:bj-finite-normal-open} applies with $S=T$ and
$G=\{\pm1\}$ acting by $u\mapsto-u$.

Let us see the orbit description in every stratum of the support map.
Over the generic point of $X$, the residue field is $k(x)$, and a
$k$-trivial absolute value on it is a valuation of $k(x)$; its
extensions to $k(u)$ form one orbit under $u\mapsto-u$, so there are one
or two of them.  For
$v=t\ord_{x=a}$ with $a\ne0$ one has two distinct lifts
$t\ord_{u=c}$ and $t\ord_{u=-c}$, where $c^2=a$; for $v=t\ord_{x=0}$
there is one lift, namely $\tfrac t2\ord_{u=0}$, the factor $\tfrac12$
being forced by $2\ord_{u=0}(u)=\ord_{u=0}(u^2)$.  Over a closed point
$x=a\ne0$ of $X$, the fibre of the support map is the single trivial
semivaluation $v_{\{x=a\},\triv}$, whose lifts are the two trivial
semivaluations at $u=\pm c$; over $x=0$ there is one lift.

Openness may be checked directly here.  A basic open neighbourhood of
$\tfrac t2\ord_{u=0}$ in $T^{\An}$ is
$\{|u|\in(\rme^{-t/2-\varepsilon},\rme^{-t/2+\varepsilon})\}$, whose image
is $\{|x|\in(\rme^{-t-2\varepsilon},\rme^{-t+2\varepsilon})\}$, an open
neighbourhood of $t\ord_{x=0}$.  The point of the general argument is that
this computation is uniform in the stratum: the image of an open set is a
finite union of images under the group, and each of those is open.
\end{example}

\begin{example}[Failure of openness over a nodal curve]
\label{ex:ch41-node-not-open}
Normality of the target cannot be dropped.  Let $X$ be the nodal plane
cubic, let $\nu:T\to X$ be its normalization, so that $T\cong\mathbf P^1$
and $\nu$ is finite, surjective and birational, with
$\nu^{-1}(p)=\{q_1,q_2\}$ over the node $p$.  Then $\nu^{\An}$ is
\emph{not} open.  Since $K(T)=K(X)$, the map $\nu^{\An}$ restricts to a
bijection on the full-support loci, while
$\nu^{\An}(v_{q_1,\triv})=\nu^{\An}(v_{q_2,\triv})=v_{p,\triv}$.  Choose
disjoint open sets $V_1\ni v_{q_1,\triv}$ and $V_2\ni v_{q_2,\triv}$ in
$T^{\An}$, which exist because $T^{\An}$ is Hausdorff.  For $t>0$ let
$t\ord_{q_i}\in T^{\val}$ be the scaled divisorial valuation at $q_i$;
by \Cref{prop:ch32-scaling-limits}, $t\ord_{q_2}\to v_{q_2,\triv}$ as
$t\to+\infty$, so $t\ord_{q_2}\in V_2$ for all large $t$, and hence
$t\ord_{q_2}\notin V_1$.  Its image $\nu^{\An}(t\ord_{q_2})$ therefore
does not lie in $\nu^{\An}(V_1)$, the preimage of a full-support point
being unique.  On the other hand $\nu^{\An}$ is continuous, so
$\nu^{\An}(t\ord_{q_2})\to v_{p,\triv}$.  Thus every neighbourhood of
$v_{p,\triv}$ contains points outside $\nu^{\An}(V_1)$, and
$\nu^{\An}(V_1)$, which contains $v_{p,\triv}$, is not a neighbourhood
of it.  Hence $\nu^{\An}$ is not open.

The regularization identity \eqref{eq:bj-finite-open-usc} fails as well.
Let
\[
 A:=\bigl\{\nu^{\An}(t\ord_{q_1})\ :\ t>0\bigr\}\subset X^{\An},
\]
and put $a:=\one_A$, a bounded function on $X^{\An}$.  Every
neighbourhood of $v_{p,\triv}$ meets $A$, since
$\nu^{\An}(t\ord_{q_1})\to v_{p,\triv}$ as $t\to+\infty$; hence
$a^*(v_{p,\triv})=1$.  On the other hand
$\nu^*a=\one_{\{t\ord_{q_1}:t>0\}}$, because a full-support point of
$X^{\An}$ has exactly one preimage and a proper-support point of
$T^{\An}$ is never mapped into $A$.  The closure of
$\{t\ord_{q_1}:t>0\}$ in $T^{\An}$ is obtained by adding the two limits
$v_{q_1,\triv}$ and $v_{T,\triv}$, again by
\Cref{prop:ch32-scaling-limits}, and does not contain $v_{q_2,\triv}$.
Therefore $(\nu^*a)^*(v_{q_2,\triv})=0$, while
$\nu^*(a^*)(v_{q_2,\triv})=a^*(v_{p,\triv})=1$.  Hence
$(\nu^*a)^*\ne\nu^*(a^*)$.
This is a bounded-function example of failure to commute with
regularization.  It does not exhibit a psh-family counterexample to
the envelope property.  The branch-asymmetric normalization example
in \Cref{sec:ch37-node-cusp} has the same proof-scope distinction.
\end{example}

\section{Flattening the finite algebra}
\label{sec:ch41-flattening}

The determinant construction is simplest for a finite locally free map.
Raynaud--Gruson flattening produces such a map after a birational
modification of the target.

\subsection{Why a finite morphism need not be flat}
\label{sec:ch41-nonflat}

Let $h:T\to X$ be as in \eqref{eq:ch41-standing} and put
$\mathcal A:=h_*\cO_T$, a coherent sheaf of $\cO_X$-algebras, finite as
an $\cO_X$-module.  Since $T$ and $X$ are integral and $h$ is
dominant, $\mathcal A$ is torsion free and its generic rank is
$d=[K(T):K(X)]$.  By generic freeness there is a dense open
$U\subseteq X$ over which $\mathcal A$ is locally free of rank $d$.
Over the complement, however, torsion-freeness is not enough.

\begin{example}[A finite degree-two morphism that is not flat]
\label{ex:ch41-nonflat-cone}
Let $\operatorname{char}k\ne2$, let
\[
 X:=\{AB=C^2\}\subset\mathbf P^3
\]
be the projective quadric cone, that is, the cone over a smooth conic, with
vertex $o=[0:0:0:1]$, and let $T:=\mathbf P^2$.  Identify $X$ with the
weighted projective plane $\mathbf P(1,1,2)$, with weighted coordinates
$a,b$ of weight $1$ and $z$ of weight $2$, so that
$A=a^2,\ B=b^2,\ C=ab$ and the fourth coordinate is $z$.  Let
\[
 h:T=\mathbf P^2\longrightarrow X=\mathbf P(1,1,2),
 \qquad
 [u:v:w]\longmapsto[u:v:w^2],
\]
which is the quotient by the involution
$\iota[u:v:w]=[u:v:-w]$.  It is a finite surjective morphism of degree
$d=2$ between normal integral projective surfaces --- $\mathbf P^2$ is
smooth, and $\mathbf P(1,1,2)$ is normal with a single singular point at
the vertex $o=[0:0:1]$.

We claim $h$ is not flat at $o$.  Work in the affine chart
$w\ne0$ upstairs, with coordinates $s=u/w,\ t=v/w$; the involution is
$(s,t)\mapsto(-s,-t)$, and the corresponding chart downstairs is
\[
 R:=k[s,t]^{\iota}=k[s^2,st,t^2]\cong k[A,B,C]/(AB-C^2),
\]
the affine quadric cone, with the vertex corresponding to the maximal ideal
$\mathfrak m=(A,B,C)=(s^2,st,t^2)$.  As an $R$-module,
\[
 k[s,t]=R\oplus M,
 \qquad
 M:=R\,s+R\,t,
\]
the decomposition into $\iota$-invariants and $\iota$-anti-invariants.
The module $M$ has rank one, since
$M\otimes_R\operatorname{Frac}R=\operatorname{Frac}R\cdot s$.  It is not
free.  To see this, grade everything by total degree in $s,t$: then
$R$ is spanned by the monomials of even degree, $M$ by the monomials of
odd degree, and $\mathfrak m=(s^2,st,t^2)$ is spanned by the monomials of
even degree $\ge2$.  Consequently $\mathfrak mM$ is spanned by the
monomials of odd degree $\ge3$, and
\[
 \dim_k M/\mathfrak mM=\dim_k\langle s,t\rangle=2 .
\]
By Nakayama's lemma, $M_{\mathfrak m}$ needs exactly two generators over
$R_{\mathfrak m}$, whereas a free module of rank one needs one.  Hence
$M$, and with it $k[s,t]=R\oplus M$, is not free at $\mathfrak m$; a
finite module over a noetherian local ring is flat if and only if it is
free, so $h$ is not flat at $o$.

The class-group description makes the same point conceptually:
the rank-one reflexive module $M$ represents the nontrivial element of $\operatorname{Cl}(R)\cong\mathbf
Z/2$, and a rank-one reflexive module over a normal domain is free at a
point exactly when its local divisor class is trivial.  The flat locus of $h$ is
therefore exactly $X\setminus\{o\}$, a dense open set, as the general
theory predicts.
\end{example}

\begin{remark}[What flatness would buy, and what replaces it]
\label{rem:ch41-why-flat}
Every construction in \Cref{sec:ch41-determinant} needs a locally free
algebra: the determinant of multiplication by an element is defined on a
locally free module, the rank must be constant for the polynomial law to be
homogeneous of a fixed degree, and the norm line
\eqref{eq:bj-determinant-norm-line} is formed from determinants of
locally free sheaves.  \Cref{ex:ch41-nonflat-cone} shows that these are not
available for $h$ itself.  A generic field norm may still exist on a normal target even when
the finite algebra is not locally free:
\cite[\href{https://stacks.math.columbia.edu/tag/0BD3}{Tag 0BD3}]
{stacks-project}.  Norms of units can then glue a line-bundle norm
using frames on full finite preimages.  This observation does not
supply the arbitrary-base-change determinant package or a fibre
algebra of constant dimension $d$.
The response adopted here, and the one
the plan prescribes, is to change the base: flatten $\mathcal A$ by a
birational modification of $X$, perform the determinant construction on
the flattened algebra, and return to $X$ at the very end by birational
reflection.  The price is that the averaging is performed on a birational
model, which is why \Cref{chap:bj-reflection} is a prerequisite for this
chapter and not the other way round.
\end{remark}

\subsection{The flattening square}
\label{sec:ch41-flatification}

\begin{lemma}\label{lem:bj-finite-flatification}
Let $h:T\to X$ be as in \eqref{eq:bj-finite-map}, and set
\[
                              d=[K(T):K(X)].
\]
There is a commutative diagram
\[
\begin{tikzcd}[column sep=large,row sep=large]
W \arrow[r,"\alpha"] \arrow[d,"q"'] &
T \arrow[d,"h"]\\
Z \arrow[r,"\beta"'] &
X
\end{tikzcd}
\tag{41.12}\label{eq:bj-flatification-square}
\]
in which:
\begin{enumerate}
\item $Z$ is normal and projective, and $\beta$ is projective
      birational;
\item $W$ is integral and projective, and $\alpha$ is projective
      birational;
\item $q$ is finite faithfully flat of constant rank $d$.
\end{enumerate}
For every $\theta\in N^1(X)_{\mathbf R}$,
\[
                         q^*\beta^*\theta=\alpha^*h^*\theta.
\tag{41.13}\label{eq:bj-flatification-class}
\]
\end{lemma}

\begin{proof}
\emph{Step 1: flatten the algebra.}
Let $\mathcal A=h_*\cO_T$, and let $U\subset X$ be its dense flat locus.
Over $U$, the finite module $\mathcal A$ is locally free of rank $d$.
Such a $U$ exists and is dense by generic freeness, and it is open because
the flat locus of a finitely presented module is open; on the integral base
$X$ the rank there is the generic rank, namely $d$.
Flattening by blowup
\cite[\href{https://stacks.math.columbia.edu/tag/0815}{Tag 0815}]
{stacks-project} produces a $U$-admissible projective blowup
\[
                         \beta_0:Z_0\longrightarrow X
\]
such that the strict transform $\mathcal A_0$ is flat and finitely presented.
This is \Cref{thm:appB-flattening}, applied with base $S=X$, which is
noetherian and hence quasi-compact and quasi-separated, with
$f=\operatorname{id}_X$, and with the coherent module $\mathcal A$; the
hypotheses of the tag are checked in \Cref{sec:stacks-flattening}, and the
distinction between this tag and the domination tag 081T is
\Cref{rem:stacks-two-blowup-theorems}.
The subsheaf removed in taking the strict transform is a two-sided ideal,
so $\mathcal A_0$ remains a finite algebra.  Concretely, the strict
transform of a module is the quotient of $\beta_0^*\mathcal A$ by the
subsheaf of sections supported on the exceptional locus, and for a sheaf of
algebras that subsheaf is closed under multiplication by
$\beta_0^*\mathcal A$ on either side, since the support of a product is
contained in the support of each factor; hence the quotient inherits a unital
algebra structure and the surjection
$\beta_0^*\mathcal A\twoheadrightarrow\mathcal A_0$ is a morphism of
algebras.  It is therefore finite locally
free of rank $d$, by the finite-flat criterion
\cite[\href{https://stacks.math.columbia.edu/tag/02K9}{Tag 02K9}]
{stacks-project}: it is finite as a module, flat, and of finite
presentation, and its rank is locally constant, hence constant on the
connected base $Z_0$ and equal to the rank over $U$, which is $d$.
Since $d>0$, $q_0:\Spec_{Z_0}\mathcal A_0\to Z_0$ is faithfully flat.

\emph{Step 2: normalize the base.}
The blowup $Z_0$ is integral by the integral blowup result in
\cite[\href{https://stacks.math.columbia.edu/tag/01OF}{Tag 01OF}]
{stacks-project}; the centre of $\beta_0$ is disjoint from the dense open
$U$, hence is a nonzero ideal on the integral scheme $X$, which is the
hypothesis of that result.  Normalize it:
\[
                         \nu:Z\longrightarrow Z_0.
\]
Normalization is finite for a finite-type integral scheme over a field by
\cite[\href{https://stacks.math.columbia.edu/tag/035S}{Tag 035S}]
{stacks-project}; every such scheme is excellent, hence Nagata, as recorded
in \Cref{sec:stacks-flattening}.
Thus $Z$ is normal projective --- projective because a finite morphism to
a projective scheme is projective --- and
$\beta=\beta_0\circ\nu$ is projective birational: $\beta_0$ is birational
because a blowup with nonzero centre on an integral scheme is, and $\nu$
is birational by construction.  Define
\[
                         W:=\operatorname{Spec}_Z(\nu^*\mathcal A_0).
\tag{41.14}\label{eq:bj-flatification-W}
\]
Then $q:W\to Z$ is finite faithfully flat of rank $d$, flatness and the
rank being preserved by the base change $\nu$.

\emph{Step 3: $W$ is integral with function field $K(T)$.}
Let $V=\Spec R\subseteq Z$ be a nonempty affine open, and let
$B:=(\nu^*\mathcal A_0)(V)$, a finite locally free $R$-module of rank $d$ and an
$R$-algebra.  Flatness of $B$ over the domain $R$ makes $B$
torsion free, so $B$ injects into
$B\otimes_R\operatorname{Frac}R=B\otimes_R K(X)$.  The latter is the
generic fibre of the algebra, which over $U$ is $\mathcal A|_U$; by
construction the strict transform agrees with the pullback over $U$, and
the generic fibre of $\mathcal A$ is $K(T)$, the function field of
$T$, because $h$ is finite and $T$ is integral.  Hence
$B\hookrightarrow K(T)$ is a subring of a field, therefore a domain.  The point represented by the generic field is dense in every such
affine inverse image.  These inverse images cover $W$, so the same point
is dense in $W$, and the domain computations show that $W$ is reduced.
Thus $W$ is integral with $K(W)=K(T)$.

\emph{Step 4: the horizontal maps and the class identity.}
The quotient
$\beta_0^*\mathcal A\twoheadrightarrow\mathcal A_0$, followed by base change along
$\nu$, gives a surjection
$\beta^*\mathcal A\twoheadrightarrow\nu^*\mathcal A_0$ of
$\cO_Z$-algebras, hence a
closed immersion
\[
                         W\hookrightarrow \Spec_Z(\beta^*\mathcal A)
                         =T\times_XZ.
\]
Composing with the first projection defines a morphism $\alpha:W\to T$,
which is projective because $W$ is projective over $Z$, hence over
$X$, hence over $T$ after using that $T$ is separated; and it is
birational because $K(W)=K(T)$ by Step~3 and $\alpha$ is dominant.  The
square \eqref{eq:bj-flatification-square} commutes by construction, both
composites $W\to X$ being the structure morphism of $W$ over $X$.
Finally \eqref{eq:bj-flatification-class} follows from functoriality of
pullback on $N^1$: both sides equal the pullback of $\theta$ along the
common composite $W\to X$.
\end{proof}

\begin{remark}[Flatification safeguard]
\label{rem:ch41-flatification-safeguard}
The flattening square does not make $h$ flat.  It replaces $h$ by the
finite flat map $q$ after simultaneous birational changes.  Nor does the
construction make $W$ normal; the later argument uses only
Fubini--Study data on $W$ and no psh closure theorem there.
\end{remark}

\begin{warning}{Three things the flattening square does not provide}
\label{warn:ch41-flattening-scope}
First, $W$ is not asserted to be normal, and in general it is not; every
statement applied to $W$ below --- \Cref{prop:ch41-finite-pullback} for
Fubini--Study functions, and the determinant construction of
\Cref{sec:ch41-determinant} --- is stated for integral, not normal, sources
for exactly this reason.  In particular one may not apply
\Cref{lem:bj-finite-normal-open} to $q$, and we never do.

Second, $\alpha$ and $\beta$ are birational but need not be
isomorphisms anywhere outside the flat locus; in
\Cref{ex:ch41-flatten-cone} they contract a curve.

Third, the square is not canonical.  The blowup produced by
\Cref{thm:appB-flattening} depends on choices, and no minimality is claimed.
Nothing below uses uniqueness; every conclusion is drawn for one square,
fixed once at the beginning of the proof of
\Cref{thm:bj-finite-normal-descent}.
\end{warning}

\subsection{Ranks, pushforward, and the projection formula}
\label{sec:ch41-projection}

The class bookkeeping of \Cref{sec:ch41-descent} uses a pushforward
on numerical Cartier classes induced by the constant-rank norm.
We prove that this map is well defined and satisfies the projection
identity, and record its compatibility with divisor-cycle pushforward.
The cycle identity $q_*[W]=d[Z]$ alone would not define a pushforward
on numerical Cartier classes.

\begin{lemma}
\label{lem:ch41-projection-formula}
Let $q:W\to Z$ be finite faithfully flat of constant rank $d$ between
integral projective varieties, with $Z$ normal.  Define $q_*$
on line-bundle classes by the determinant norm.  Then $q_*:N^1(W)_{\mathbf R}\to
N^1(Z)_{\mathbf R}$ is well defined, and
\[
 q_*q^*\eta=d\,\eta
 \qquad\text{for every }\eta\in N^1(Z)_{\mathbf R}.
\tag{41.15}\label{eq:ch41-projection-formula}
\]
Moreover, for every line bundle $M$ on $W$ and every nonzero rational
section $s$ of $M$,
\[
 q_*\bigl(\operatorname{div}s\bigr)
 =\operatorname{div}\bigl(\mathrm N_q(s)\bigr),
\]
where $\mathrm N_q$ is the determinant norm of
\Cref{def:ch41-norm-line}.
\end{lemma}

\begin{proof}
The norm is multiplicative, so it defines an additive map on Picard
groups.  To descend it to numerical classes, take an integral curve
$C\subset Z$ and normalize it.  Base change of the norm gives
\[
 \deg_C\mathrm N_q(M)
 =\chi(W_C,M|_{W_C})-\chi(W_C,\cO_{W_C})
 =\sum_\Gamma \ell_\Gamma\,\deg_\Gamma M,
\]
where $W_C=W\times_Z C^\nu$, and the sum is over its integral
one-dimensional components with their generic lengths.  Each
component $\Gamma$ maps finitely onto an integral curve $\Gamma'$
in $W$, and its degree term is
$[K(\Gamma):K(\Gamma')]\deg_{\Gamma'}M$.
The first equality is the determinant degree formula on the smooth
curve $C^\nu$, since $q_C$ is finite locally free.  The second is the
degree formula for a line bundle on a projective curve, including a
nonreduced curve, or equivalently its intersection with the fundamental
one-cycle \cite{Ful}.  Zero-dimensional contributions cancel.
If $M$ is numerically trivial on $W$, every term is zero.
Thus $\mathrm N_q(M)$ is numerically trivial on $Z$, and this
defines the linear map $q_*$ on $N^1$.

For a line bundle $L$ on $Z$,
$\mathrm N_q(q^*L)=L^{\otimes d}$ by
\Cref{prop:ch41-norm-line-properties}\,\cref{itm:ch41-nl:pullback}.
This proves \eqref{eq:ch41-projection-formula} and extends by
linearity to real classes.  It uses the constant-rank norm functor;
no general pushforward of Cartier numerical classes for arbitrary
proper maps is asserted.  On the present normal target its compatibility
with divisor-cycle pushforward is the final assertion, proved in
\Cref{prop:ch41-norm-line-properties}\,\cref{itm:ch41-nl:divisor}.
\end{proof}

\subsection{The flattening square computed}
\label{sec:ch41-flatten-example}

\begin{example}[Flattening the double cover of the quadric cone]
\label{ex:ch41-flatten-cone}
We compute the square \eqref{eq:bj-flatification-square} for the morphism
$h:\mathbf P^2\to\mathbf P(1,1,2)$ of \Cref{ex:ch41-nonflat-cone}, with
$\operatorname{char}k\ne2$ and $d=2$.  Here the flattening blowup can be
written down: it is the blowup of the vertex.

Let $\beta:Z\to X$ be the blowup of $X=\mathbf P(1,1,2)$ at the vertex
$o$, and let $\alpha:W\to T=\mathbf P^2$ be the blowup at the point
$\tilde o=[0:0:1]=h^{-1}(o)$.  Then $W\cong\mathbf F_1$, with
exceptional curve $F$ satisfying $F^2=-1$.  The involution $\iota$
lifts to $W$; in the blowup chart $(s,\rho)$ with $t=s\rho$, the lift
acts by $(s,\rho)\mapsto(-s,\rho)$, so it fixes the exceptional curve
$\{s=0\}$ pointwise.  Consequently
\[
 q:W\longrightarrow Z:=W/\iota
\]
is a finite morphism of degree $2$ between \emph{smooth} surfaces: the
quotient of a smooth surface by an involution whose fixed locus is a smooth
divisor is smooth.  By miracle flatness --- a finite surjective morphism
from a Cohen--Macaulay scheme to a regular scheme with equidimensional
fibres is flat, \cite[Theorem~23.1]{Mat89} --- the morphism $q$ is finite
flat of rank $2$.  It is therefore already the flattened cover, and the
square
\[
\begin{tikzcd}[column sep=large,row sep=large]
W=\mathbf F_1 \arrow[r,"\alpha"] \arrow[d,"q"'] &
T=\mathbf P^2 \arrow[d,"h"]\\
Z \arrow[r,"\beta"'] &
X=\mathbf P(1,1,2)
\end{tikzcd}
\]
commutes: both composites are the quotient of $\mathbf F_1$ by $\iota$
followed by the contraction, in either order.

The numbers are as follows.  Write $E:=q(F)\subset Z$ for the image of the
exceptional curve.  The involution fixes $F$ pointwise, so $q$ is
ramified along $F$ with index $2$, and therefore
$q^*E=2F$; by the projection formula for the degree-two map $q$,
\[
 2\,(E^2)=(q^*E)^2=4\,(F^2)=-4,
 \qquad\text{so}\qquad
 (E^2)=-2 .
\]
Thus $Z$ carries a $(-2)$-curve contracted by $\beta$.  Since $Z$ is
smooth, $\beta$ is proper birational and an isomorphism away from the
$A_1$ point $o$, the surface $Z$ is the minimal resolution of $X$,
and hence $Z\cong\mathbf F_2$ with $E$ its negative section.  The
identity
\eqref{eq:ch41-projection-formula} reads $q_*q^*=2$ on
$N^1(Z)_{\mathbf R}$, which for the class of $E$ says
$q_*(2F)=2E$, that is, $q_*F=E$; this is the statement that $q$ maps
$F$ isomorphically onto $E$, which it does, being ramified there.

Two features of this computation are typical.  First, the flattening blowup
is supported exactly on the non-flat locus $\{o\}$ of
\Cref{ex:ch41-nonflat-cone}, as it must be, since
\Cref{thm:appB-flattening} produces a $U$-admissible blowup for $U$ the
flat locus.  Second, the source is modified as well:
$\alpha$ is not an isomorphism, and $W=\mathbf F_1$ is not $T$.  A
reader who expects the flattening to change only the target should note that
$W$ is defined as $\Spec_Z$ of the flattened algebra, and that its
comparison with $T$ is a separate birational statement, Step~4 of the
proof of \Cref{lem:bj-finite-flatification}.
\end{example}

\section{Determinant norms and universal polynomial laws}
\label{sec:ch41-determinant}

Throughout this section $q:W\to Z$ is a finite locally free morphism of
schemes of constant rank $d\ge1$, and $A\to B$ denotes the corresponding
map of rings on an affine open over which $q_*\cO_W$ is free.
The element, polynomial-law and line-bundle constructions require no
variety, algebraic-closure or characteristic hypothesis, so they apply in
\Cref{sec:ch41-averaging} after an arbitrary isometric field extension.
The separate divisor assertion in
\cref{itm:ch41-nl:divisor} assumes integral schemes and a normal target.

\subsection{Norms of elements}
\label{sec:ch41-element-norm}

Recall from \eqref{eq:appA-element-norm} that for $A\to B$ finite locally
free of rank $d$ and $b\in B$ one sets
$\mathrm N_{B/A}(b):=\det(m_b)$, the determinant of multiplication by
$b$, computed on the rank-$d$ locally free $A$-module $B$.  The
elementary identities are \Cref{prop:appA-norm-identities}.  We record the
two further properties that are used below and that are not in the
appendix.

\begin{lemma}
\label{lem:ch41-element-norm}
Let $A\to B$ be finite locally free of rank $d$.
\begin{enumerate}
\item\label{itm:ch41-en:scalar}
      For $a\in A$ and $b\in B$, one has
      $\mathrm N_{B/A}(ab)=a^d\,\mathrm N_{B/A}(b)$; in particular
      $\mathrm N_{B/A}(a)=a^d$.
\item\label{itm:ch41-en:unit}
      $b\in B$ is a unit if and only if $\mathrm N_{B/A}(b)$ is a unit
      in $A$.
\item\label{itm:ch41-en:zero}
      If $A$ is a field and $B\ne0$, then
      $\mathrm N_{B/A}(b)=0$ if and only if $b$ is a zero divisor in
      $B$; if moreover $B$ is local Artinian with residue field $A$,
      then $\mathrm N_{B/A}(b)=\bar b^{\,d}$, where
      $\bar b\in A$ is the residue of $b$.
\item\label{itm:ch41-en:product}
      If $B=B_1\times\cdots\times B_n$ as an $A$-algebra, with each
      $B_s$ finite locally free of rank $d_s$, then $\sum_sd_s=d$ and
      \[
       \mathrm N_{B/A}(b_1,\dots,b_n)
       =\prod_{s=1}^n\mathrm N_{B_s/A}(b_s).
      \]
\end{enumerate}
\end{lemma}

\begin{proof}
\cref{itm:ch41-en:scalar}  Multiplication by $ab$ is $a$ times
multiplication by $b$ as an $A$-linear endomorphism of a rank-$d$
module, and $\det(a\,\varphi)=a^d\det\varphi$.  Taking $b=1$ gives the
second formula.

\cref{itm:ch41-en:unit}  If $b$ is a unit then
$\mathrm N(b)\mathrm N(b^{-1})=\mathrm N(1)=1$ by
\Cref{prop:appA-norm-identities}.  Conversely, if $\det m_b$ is a unit
then $m_b$ is an isomorphism of $A$-modules, so there is $c\in B$
with $bc=1$, namely $c=m_b^{-1}(1)$.

\cref{itm:ch41-en:zero}  Over a field, $m_b$ is an endomorphism of a
finite-dimensional vector space; its determinant vanishes exactly when it
has a kernel, that is, when $bc=0$ for some $c\ne0$.  If $B$ is local
Artinian with maximal ideal $\mathfrak n$ and residue field $A$, write
$b=\bar b+n$ with $n\in\mathfrak n$; then $m_n$ is nilpotent, because
$\mathfrak n$ is, so the characteristic polynomial of
$m_b=\bar b\cdot\operatorname{id}+m_n$ is $(T-\bar b)^d$ and its
determinant is $\bar b^{\,d}$.

\cref{itm:ch41-en:product}  Multiplication by $(b_1,\dots,b_n)$ preserves
each factor, so it is block diagonal with respect to the decomposition
$B=\bigoplus_sB_s$; the determinant of a block diagonal endomorphism is
the product of the determinants of the blocks, and ranks add.
\end{proof}

\begin{remark}[Norm versus trace, once more]
\label{rem:ch41-norm-vs-trace}
\Cref{lem:ch41-element-norm}\,\cref{itm:ch41-en:zero} is the exact point at
which the norm outperforms the trace.  On a local Artinian algebra with
residue field $A$ the norm reads off the residue and raises it to the
$d$-th power --- so it sees the residue exactly, and records the length as
an exponent.  The trace of the same element is
$d\,\bar b+\Tr(m_n)=d\,\bar b$, because the trace of a nilpotent
endomorphism vanishes; this is $0$ whenever $\operatorname{char}A\mid d$,
and it carries no information about $\bar b$ at all in that case.  The
exponent $d$ in $\bar b^{\,d}$ is what becomes the weight $e(w/z)$ in
\eqref{eq:bj-finite-flat-average}.
\end{remark}

\subsection{Polynomial laws}
\label{sec:ch41-polynomial-laws}

The norm is multiplicative and generally not additive.
In a local basis it is a homogeneous polynomial; the polynomial-law
formulation records its naturality under all base changes without choosing
that basis as part of the definition.  The correct notion is due to Roby, and we use
it in the following form.

\begin{definition}[Polynomial law]
\label{def:ch41-polynomial-law}
Let $A$ be a ring and let $V,U$ be $A$-modules.  A \emph{polynomial
law} $\Phi:V\rightsquigarrow U$ is a family of maps of sets
\[
 \Phi_R:V\otimes_AR\longrightarrow U\otimes_AR,
\]
one for every commutative $A$-algebra $R$, such that for every morphism
$\rho:R\to R'$ of $A$-algebras the square
\[
\begin{tikzcd}[column sep=large]
V\otimes_AR\arrow[r,"\Phi_R"]\arrow[d,"\operatorname{id}\otimes\rho"']&
U\otimes_AR\arrow[d,"\operatorname{id}\otimes\rho"]\\
V\otimes_AR'\arrow[r,"\Phi_{R'}"']&U\otimes_AR'
\end{tikzcd}
\]
commutes.  The law is \emph{homogeneous of degree $d$} if
$\Phi_R(rx)=r^d\Phi_R(x)$ for all $A$-algebras $R$, all $r\in R$ and
all $x\in V\otimes_AR$.  If $V$ and $U$ are $A$-algebras, the law is
\emph{multiplicative} if every $\Phi_R$ is multiplicative and
$\Phi_R(1)=1$.
\end{definition}

The definition is exactly designed so that ``compatible with arbitrary base
change'' is not an extra theorem but part of the data.  The content of the
next proposition is that the determinant norm is such a law, and that the
naive description by a polynomial in coordinates is recovered from it by
evaluating on a polynomial algebra.

\begin{proposition}
\label{prop:ch41-norm-law}
Let $A\to B$ be finite locally free of rank $d$.  Then
\[
 \mathrm N_{B/A}:B\rightsquigarrow A
\]
is a multiplicative polynomial law, homogeneous of degree $d$, in the
sense of \Cref{def:ch41-polynomial-law}.  Explicitly, for every
$A$-algebra $R$ the map
$\mathrm N_R:B\otimes_AR\to R$ is
$b\mapsto\det_R(m_b)$, and for every $A$-algebra map $R\to R'$,
\[
 \mathrm N_{R'}\bigl(b\otimes1\bigr)
 =\text{image of }\mathrm N_R(b)\text{ in }R' .
\]
Moreover $\mathrm N_{B/A}$ is determined by its \emph{universal} values:
if $\Phi:B\rightsquigarrow A$ is a polynomial law such that, locally on
$\Spec A$ and for one local basis $b_1,\dots,b_d$ of $B$,
\[
 \Phi_{A[\mathbf T]}\Bigl(\sum_iT_ib_i\Bigr)
 =\mathrm N_{A[\mathbf T]}\Bigl(\sum_iT_ib_i\Bigr)
 \quad\text{in }A[T_1,\dots,T_d],
\]
then $\Phi=\mathrm N_{B/A}$.
\end{proposition}

\begin{proof}
For every $A$-algebra $R$, the ring $B_R:=B\otimes_AR$ is finite
locally free of rank $d$ over $R$, because finite local freeness is
stable under base change.  Define $\mathrm N_R(x):=\det_R(m_x)$ for
$x\in B_R$, the determinant of multiplication by $x$ on the rank-$d$
locally free $R$-module $B_R$; this is well defined because the
determinant of an endomorphism of a locally free module of constant rank is
a globally defined element of the base, computed locally by any basis and
independent of it.

\emph{Base change.}  Let $\rho:R\to R'$.  Then
$B_{R'}=B_R\otimes_RR'$, and multiplication by the image of $x$ is the
base change of $m_x$.  The determinant of an endomorphism of a locally
free module commutes with base change, since it does so for matrices; hence
$\rho(\mathrm N_R(x))=\mathrm N_{R'}(x\otimes1)$.  This is the required
commuting square.

\emph{Homogeneity and multiplicativity.}  For $r\in R$ and $x\in B_R$,
$m_{rx}=r\,m_x$, so $\det(m_{rx})=r^d\det(m_x)$.  Multiplicativity and
$\mathrm N_R(1)=1$ are $m_{xy}=m_xm_y$ and $m_1=\operatorname{id}$.

\emph{Determination by the universal value.}  The assertion is local on
$\Spec A$, so assume $B$ free with basis $b_1,\dots,b_d$ and suppose
$\Phi_{A[\mathbf T]}(\xi)=\mathrm N_{A[\mathbf T]}(\xi)$, where
$A[\mathbf T]=A[T_1,\dots,T_d]$ and
$\xi:=\sum_ib_i\otimes T_i\in B\otimes_AA[\mathbf T]$.  Let $R$ be any
$A$-algebra and $x\in B\otimes_AR$.  Since $B\otimes_AR$ is free over
$R$ on $b_1\otimes1,\dots,b_d\otimes1$, we may write
$x=\sum_ib_i\otimes r_i$ with $r_i\in R$, and the $A$-algebra map
$\rho:A[\mathbf T]\to R$ with $T_i\mapsto r_i$ carries $\xi$ to
$x$.  Naturality of $\Phi$ and of $\mathrm N$ for $\rho$ then gives
\[
 \Phi_R(x)=\rho\bigl(\Phi_{A[\mathbf T]}(\xi)\bigr)
 =\rho\bigl(\mathrm N_{A[\mathbf T]}(\xi)\bigr)
 =\mathrm N_R(x).
\]
Independence of the chosen basis follows, since two bases differ by an
element of $\operatorname{GL}_d(A)$ and the corresponding change of
variables is an $A$-algebra automorphism of $A[\mathbf T]$.
\end{proof}

\begin{warning}{A polynomial law is not determined by its values on $B$}
\label{warn:ch41-law-not-determined-on-B}
The uniqueness clause of \Cref{prop:ch41-norm-law} refers to the value on
the polynomial algebra, not to the value on $B$ itself, and the
distinction is real.  Over $A=\mathbf F_p$ and $B=A$, the two
polynomial laws $x\mapsto x$ and $x\mapsto x^p$ agree on
$B\otimes_AA=\mathbf F_p$, where $x^p=x$, and differ on
$B\otimes_AA[T]$.  Thus values on the original module do not determine an arbitrary
polynomial law over every base.  This does not invalidate pointwise
determination when additional hypotheses provide enough test values;
the universal-value proof works without such hypotheses. It is distinct from the trace degeneracy in \Cref{ex:ch41-frobenius-average}.
\end{warning}

\begin{remark}[Where the polynomial-law formulation is doing work]
\label{rem:ch41-law-work}
The book plan singles out ``universal properties of determinant norms'' for
audit, and this is the place.  Three things would go wrong with a naive
formulation.

\emph{(a) A norm defined only on elements of $B$ is not enough.}  In
\Cref{sec:ch41-averaging} the norm is applied to
$\sum_jT_js_j$, an element of $B\otimes_AA[\mathbf T]$ and not of
$B$; the coefficients extracted from the result are the objects that
become sections downstairs.  Without a norm on
$B\otimes_AA[\mathbf T]$, and without knowing that it is compatible with
the norm on $B$, the construction does not even typecheck.

\emph{(b) The degree must be exactly $d$, uniformly.}  Homogeneity of
degree $d$ in \Cref{def:ch41-polynomial-law} says
$\mathrm N(rx)=r^d\mathrm N(x)$ for all $r$ in all base changes; this
is what makes $\mathrm N(\sum_jT_js_j)$ a homogeneous polynomial of degree
exactly $d$ in the $T_j$, and hence what makes the multi-indices $I$
in \eqref{eq:bj-universal-norm-polynomial} run over
$|I|=d$.  The identity $\sum_we(w/z)=d$ of
\eqref{eq:bj-finite-flat-average} is the fibrewise shadow of this
homogeneity.

\emph{(c) Base change must be arbitrary, not just flat or finite.}  The base
changes used below are $A\to\kappa(\xi)\to\mathcal H(z)\to\Omega$: a
residue field, a completion, and an algebraically closed isometric
extension.  The field-extension steps are flat, while the residue-field map need
not be flat.  The completion and algebraic closure need not be finite
or of finite type; a trivial completion can of course be the original
field.  Arbitrary base-change compatibility covers the whole chain.
The fibrewise computation of
\Cref{lem:ch41-fibre-decomposition} takes place after it.
\end{remark}

\begin{lemma}
\label{lem:ch41-universal-polynomial}
Let $A\to B$ be finite locally free of rank $d$ and let
$b_0,\dots,b_r\in B$.  Then there are uniquely determined elements
$n_I\in A$, indexed by the multi-indices
$I=(I_0,\dots,I_r)\in\mathbf N^{r+1}$ with $|I|=d$, such that in
$A[T_0,\dots,T_r]$
\[
 \mathrm N_{B/A}\Bigl(\sum_{j=0}^rT_jb_j\Bigr)
 =\sum_{|I|=d}n_I\,T^I .
\tag{41.16}\label{eq:ch41-universal-coefficients}
\]
The formation of the $n_I$ commutes with arbitrary base change: for every
$A$-algebra $R$, the coefficients of
$\mathrm N_{B_R/R}(\sum_jT_j(b_j\otimes1))$ are the images of the
$n_I$ in $R$.
\end{lemma}

\begin{proof}
Apply \Cref{prop:ch41-norm-law} with $R=A[T_0,\dots,T_r]$ to the element
$\sum_jT_jb_j\in B\otimes_AR$.  The result is an element of $R$, and by
homogeneity of degree $d$ applied to the $A[\mathbf T]$-algebra
$A[\mathbf T][\lambda]$ and the scalar $\lambda$, it is a homogeneous
polynomial of degree $d$ in the $T_j$; this defines the $n_I$
uniquely, polynomial rings being free $A$-modules on the monomials.
Base-change compatibility is the naturality square of
\Cref{def:ch41-polynomial-law} for the map
$A[\mathbf T]\to R[\mathbf T]$, applied coefficientwise.
\end{proof}

\begin{example}[A rank-two determinant norm]
\label{ex:ch41-rank-two}
Let $A$ be a ring and
$B=A[u]/(u^2-a)$, which is finite free of rank two with basis $1,u$.
For $b=x+yu\in B$, multiplication by $b$ has matrix
\[
             \begin{pmatrix}x&ay\\ y&x\end{pmatrix}
\]
in this basis.  Hence
\[
             \mathrm N_{B/A}(x+yu)=x^2-ay^2.
\]
Equivalently, the universal norm polynomial is
\[
             \mathrm N_{B/A}(T_0+T_1u)=T_0^2-aT_1^2.
\]
The calculation remains valid in characteristic two and when the algebra
is inseparable or nonreduced: it uses a determinant, not an average over two
geometric points and not division by two.  It is the rank-two local model
for the polynomial law in \Cref{prop:bj-weighted-coefficient-norm}.

For the following three field regimes, assume that $A$ is a field.
They can all be read from $T_0^2-aT_1^2$.  If $2$ is invertible and $a$ is a unit which is not a
square, the algebra is a separable field extension and the polynomial
factors over a quadratic extension as
$(T_0-\sqrt aT_1)(T_0+\sqrt aT_1)$: two sheets, weights $1,1$.  If
$a=0$, the algebra is $A[u]/(u^2)$, nonreduced, and the polynomial is
$T_0^2$: one sheet, weight $2$, and the coefficient of $T_0T_1$
vanishes.  If $\operatorname{char}A=2$ and $a$ is a nonsquare unit, the
algebra is an inseparable field extension and the polynomial is
$T_0^2+aT_1^2=(T_0+\sqrt aT_1)^2$: again one sheet with weight $2$, and
again the middle coefficient --- which is $\Tr(u)$ --- vanishes.
In all three cases the total weight is $2=d$.  The middle coefficient
vanishes in \emph{all three} cases, since
$\Tr(u)=\operatorname{tr}\begin{psmallmatrix}0&a\\1&0\end{psmallmatrix}=0$
identically, so it is not what separates them.  What degenerates in the last
two is the trace \emph{form} $(b,c)\mapsto\Tr(bc)$, whose Gram matrix in
the basis $1,u$ is $\operatorname{diag}(2,2a)$, of determinant $4a$ ---
nonzero exactly in the first case.  That is the degeneration discussed in
\Cref{rem:ch41-norm-vs-trace}.
\end{example}

\subsection{The norm of a line bundle and of a section}
\label{sec:ch41-norm-line}

\begin{definition}[Determinant norm line]
\label{def:ch41-norm-line}
Let $q:W\to Z$ be finite locally free of rank $d>0$.  For a line
bundle $M$ on $W$, define its determinant norm line by
\[
 \mathrm N_q(M):=
 \det(q_*M)\otimes\det(q_*\cO_W)^{-1}.
\tag{41.17}\label{eq:bj-determinant-norm-line}
\]
This construction is multiplicative in $M$ and commutes with arbitrary
base change.  A section $s:\cO_W\to M$ induces the determinant section
$\mathrm N_q(s)$.  Locally, after trivializing $M$, this is the determinant
of multiplication by the coefficient of $s$ on the finite locally free
algebra $q_*\cO_W$.
For functions this is the canonical norm of
\cite[\href{https://stacks.math.columbia.edu/tag/0BD2}{Tag 0BD2}]
{stacks-project}; the line-bundle formulation used here is the determinant
norm functor of \cite[Proposition~3.3]{FerrandNorm}.
\end{definition}

Definition \eqref{eq:bj-determinant-norm-line} is a definition of a line
bundle; the content is that it behaves like a norm.  We collect what is
needed, with proofs, since the local frame description in
\cref{itm:ch41-nl:frame} is the workhorse of \Cref{sec:ch41-averaging}.

\begin{proposition}
\label{prop:ch41-norm-line-properties}
Let $q:W\to Z$ be finite locally free of rank $d>0$, and let $M,M'$ be
line bundles on $W$.
\begin{enumerate}
\item\label{itm:ch41-nl:unit}
      $\mathrm N_q(\cO_W)\cong\cO_Z$ canonically, and
      $\mathrm N_q(1)=1$.
\item\label{itm:ch41-nl:mult}
      $\mathrm N_q(M\otimes M')\cong\mathrm N_q(M)\otimes\mathrm N_q(M')$
      canonically; in particular
      $\mathrm N_q(M^{\otimes m})\cong\mathrm N_q(M)^{\otimes m}$ for
      $m\in\mathbf Z$.
\item\label{itm:ch41-nl:pullback}
      For a line bundle $L$ on $Z$,
      $\mathrm N_q(q^*L)\cong L^{\otimes d}$.
\item\label{itm:ch41-nl:frame}
      If $U\subseteq Z$ is open, $q^{-1}(U)=\Spec_UB$ with $B$ free
      of rank $d$ over $A=\cO_Z(U)$, and $\tau$ is a frame of $M$
      on $q^{-1}(U)$, then $\mathrm N_q(\tau)$ is a frame of
      $\mathrm N_q(M)$ on $U$; and for
      $s=f\tau$ with $f\in B$ one has
      $\mathrm N_q(s)=\mathrm N_{B/A}(f)\cdot\mathrm N_q(\tau)$.
\item\label{itm:ch41-nl:divisor}
      If $W,Z$ are integral and $Z$ is normal, then for a nonzero
      rational section $s$ of $M$,
      \[
       \operatorname{div}\mathrm N_q(s)=q_*\operatorname{div}(s),
       \qquad
       c_1(\mathrm N_qM)=q_*c_1(M).
      \]
\item\label{itm:ch41-nl:bc}
      For every morphism $Z'\to Z$, writing $q':W'=W\times_ZZ'\to Z'$
      and $M'$ for the pullback of $M$, there is a canonical
      isomorphism $\mathrm N_{q'}(M')\cong$ the pullback of
      $\mathrm N_q(M)$, compatible with
      \cref{itm:ch41-nl:unit}--\cref{itm:ch41-nl:frame} and carrying
      $\mathrm N_q(s)$ to $\mathrm N_{q'}(s')$.
\end{enumerate}
\end{proposition}

\begin{proof}
\cref{itm:ch41-nl:unit}  Immediate from
\eqref{eq:bj-determinant-norm-line}, the two factors being inverse to each
other, and from $\mathrm N_{B/A}(1)=\det(\operatorname{id})=1$.

\cref{itm:ch41-nl:mult}  This is the multiplicativity of Ferrand's norm
functor \cite[Proposition~3.3]{FerrandNorm}.  It is not formal from
\eqref{eq:bj-determinant-norm-line}, because $q_*(M\otimes M')$ is not
$q_*M\otimes q_*M'$; the isomorphism is constructed by descent from the
local case, where after trivializing $M$ and $M'$ both sides become
$\cO_U$ and the identification is the multiplicativity
$\mathrm N_{B/A}(fg)=\mathrm N_{B/A}(f)\mathrm N_{B/A}(g)$ of
\Cref{prop:appA-norm-identities}.  The cocycle condition holds because the
transition functions of $M\otimes M'$ are the products of those of $M$
and $M'$, and $\mathrm N_{B/A}$ is multiplicative.

\cref{itm:ch41-nl:pullback}  By the projection formula
$q_*(q^*L)\cong L\otimes q_*\cO_W$, so
$\det q_*(q^*L)\cong L^{\otimes d}\otimes\det q_*\cO_W$, the rank of
$q_*\cO_W$ being $d$.  Substituting into
\eqref{eq:bj-determinant-norm-line} gives the claim.

\cref{itm:ch41-nl:frame}  With the stated hypotheses,
$q_*M|_U$ is the free $A$-module $B\tau\cong B$, and the isomorphism
$q_*M|_U\cong q_*\cO_W|_U$ given by $\tau$ induces an isomorphism of
determinants, that is, a nowhere vanishing section
$\mathrm N_q(\tau)$ of $\mathrm N_q(M)|_U$.  Under this trivialization a
section $s=f\tau$ corresponds to the endomorphism $m_f$ of $B$, whose
determinant is $\mathrm N_{B/A}(f)$; this is the local description in
\Cref{def:ch41-norm-line} and the formula asserted.  Independence of the
choice of $\tau$ is \cref{itm:ch41-nl:mult} applied to a unit.

\cref{itm:ch41-nl:divisor}  At a prime divisor $P$ of the normal
target, $\cO_{Z,P}$ is a DVR.  Choose a full-preimage frame as in
\Cref{lem:ch41-local-frame}.  For a regular nonzero coefficient $f$,
Smith normal form over this DVR identifies the order of $\det(m_f)$
with the length of $B_P/fB_P$.  Decomposing that finite-length module
at the primes $Q$ over $P$ gives
\[
 \operatorname{ord}_P\mathrm N(f)
 =\sum_{Q\mapsto P}[\kappa(Q):\kappa(P)]
       \ell_{\cO_{W,Q}}(\cO_{W,Q}/f\cO_{W,Q}).
\]
The length is the coefficient of the Cartier divisor of $f$ on the
possibly nonnormal integral source.  For a rational coefficient, write
it as a quotient of two regular nonzero coefficients and subtract
these finite integer orders.  This proves the divisor identity.
The numerical identity uses the norm-induced map proved in
\Cref{lem:ch41-projection-formula}, rather than an unrestricted
proper pushforward on $N^1$.

\cref{itm:ch41-nl:bc}  Finite local freeness, pushforward along an affine
morphism, and formation of determinants of locally free modules all commute
with arbitrary base change; hence so does
\eqref{eq:bj-determinant-norm-line}.  The compatibilities are checked on a
cover as in \cref{itm:ch41-nl:frame}, where they reduce to the base-change
statement of \Cref{prop:ch41-norm-law}.
\end{proof}

Multiplicativity extends the norm to $\mathbf Q$-line bundles:
if $M$ denotes a rational tensor class and $rM$ is an honest line
bundle, put $\mathrm N_q(M)=\frac1r\mathrm N_q(rM)$ in
$\operatorname{Pic}(Z)\otimes\mathbf Q$.  Clearing two denominators
and using \cref{itm:ch41-nl:mult} proves independence of $r$.
Section presentations below are likewise made after clearing
denominators; no actual root of a line bundle is asserted.

\subsection{Base change}
\label{sec:ch41-base-change}

We now state the base-change property in the form in which it is used, and
prove it.  This is item~\cref{itm:ch41-task:law} of
\Cref{sec:ch41-tasks}.

\begin{proposition}
\label{prop:ch41-norm-base-change}
Let $q:W\to Z$ be finite locally free of rank $d>0$, let $M$ be a
line bundle on $W$, let $s_0,\dots,s_r\in H^0(W,M)$, and let
\[
 \mathrm N_q\Bigl(\sum_jT_js_j\Bigr)=\sum_{|I|=d}c_IT^I,
 \qquad
 c_I\in H^0\bigl(Z,\mathrm N_q(M)\bigr),
\]
be the expansion supplied by \Cref{lem:ch41-universal-polynomial} and
\Cref{prop:ch41-norm-line-properties}\,\cref{itm:ch41-nl:frame}.  Let
$\pi:Z'\to Z$ be any morphism of schemes, and form the cartesian square
\[
\begin{tikzcd}
W'\arrow[r,"\pi_W"]\arrow[d,"q'"']&W\arrow[d,"q"]\\
Z'\arrow[r,"\pi"']&Z
\end{tikzcd}
\]
with $M':=\pi_W^*M$ and $s_j':=\pi_W^*s_j$.  Then $q'$ is finite
locally free of rank $d$, and, under the canonical isomorphism
$\mathrm N_{q'}(M')\cong\pi^*\mathrm N_q(M)$ of
\Cref{prop:ch41-norm-line-properties}\,\cref{itm:ch41-nl:bc},
\[
 \mathrm N_{q'}\Bigl(\sum_jT_js_j'\Bigr)=\sum_{|I|=d}\pi^*(c_I)\,T^I .
\tag{41.18}\label{eq:ch41-base-change-coefficients}
\]
In particular the individual coefficients satisfy $c_I'=\pi^*c_I$; no
regrouping or rescaling of coefficients occurs.
\end{proposition}

\begin{proof}
Finite local freeness of constant rank is stable under arbitrary base
change, which gives the first assertion.  For the second, the question is
local on $Z'$ and then on $Z$, so we may assume $Z=\Spec A$,
$Z'=\Spec A'$, $q^{-1}(Z)=\Spec B$ with $B$ free of rank $d$ over
$A$, and $M$ trivialized by a frame $\tau$ on $\Spec B$.  Write
$s_j=f_j\tau$ with $f_j\in B$.  Then
$B'=B\otimes_AA'$ is free of rank $d$ over $A'$, the pullback
$\pi_W^*\tau$ is a frame of $M'$, and $s_j'=(f_j\otimes1)\pi_W^*\tau$.
By \Cref{prop:ch41-norm-line-properties}\,\cref{itm:ch41-nl:frame},
\[
 \mathrm N_q\Bigl(\sum_jT_js_j\Bigr)
 =\mathrm N_{B/A}\Bigl(\sum_jT_jf_j\Bigr)\cdot\mathrm N_q(\tau),
\]
and similarly downstairs with $B'$, $A'$ and $\pi_W^*\tau$.  The
coefficients therefore transform as the coefficients of
$\mathrm N_{B/A}(\sum_jT_jf_j)$, and
\Cref{lem:ch41-universal-polynomial} says that these commute with the base
change $A\to A'$.  Since $\mathrm N_{q'}(\pi_W^*\tau)$ corresponds to
$\pi^*\mathrm N_q(\tau)$ under the canonical isomorphism of
\Cref{prop:ch41-norm-line-properties}\,\cref{itm:ch41-nl:bc}, the identity
\eqref{eq:ch41-base-change-coefficients} follows.  The statement is
independent of the choices because two full-preimage frames differ by a unit, whose norm
is a unit by \Cref{lem:ch41-element-norm}\,\cref{itm:ch41-en:unit}, and both
sides of \eqref{eq:ch41-base-change-coefficients} are multiplied by the same
unit.
\end{proof}

\begin{example}[Base change to a fibre]
\label{ex:ch41-base-change-fibre}
Take $A=k[x]$, $B=A[u]/(u^2-x)$, $d=2$, and let
$\pi:\Spec\kappa\to\Spec A$ be the inclusion of a point.  For
$\kappa=k(x)$, the generic point, $B\otimes_A\kappa=k(x)[u]/(u^2-x)$ is
a field, and the norm polynomial $T_0^2-xT_1^2$ is irreducible over
$\kappa$; its coefficients are $1,0,-x$.  For $\kappa=k$ the residue
field at $x=c^2\ne0$ with $\operatorname{char}k\ne2$, the fibre algebra
is $k[u]/(u^2-c^2)\cong k\times k$, and the polynomial becomes
$T_0^2-c^2T_1^2=(T_0-cT_1)(T_0+cT_1)$: the two factors are the two
sheets.  For $\kappa=k$ at $x=0$, the fibre algebra is $k[u]/(u^2)$
and the polynomial becomes $T_0^2$: one sheet with weight $2$.  In each
case the coefficients are literally the images of $1,0,-x$; this is
\eqref{eq:ch41-base-change-coefficients}, and it is what allows the
computation of \Cref{sec:ch41-averaging} to be performed fibre by fibre and
then reassembled.
\end{example}

\subsection{Frames adapted to an analytic point}
\label{sec:ch41-frames}

To evaluate $\log|c_I|$ at a point $z\in Z^{\An}$ one needs a frame of
$\mathrm N_q(M)$ near $c(z)$, and to evaluate $\log|s_j|$ at the
points $w$ over $z$ one needs a frame of $M$ near each $c(w)$.
\Cref{prop:ch41-norm-line-properties}\,\cref{itm:ch41-nl:frame} produces the
first from the second, provided the second is defined on a full preimage
$q^{-1}(U)$.  The following lemma says that this is always arrangeable.

\begin{lemma}
\label{lem:ch41-local-frame}
Let $q:W\to Z$ be finite locally free of rank $d>0$ with $Z$ a
noetherian scheme, let $M$ be a line bundle on $W$, and let $p\in Z$.
Then there is an affine open $U\subseteq Z$ containing $p$ such that:
\begin{enumerate}
\item $q^{-1}(U)$ is affine, say $q^{-1}(U)=\Spec B$ with
      $B$ free of rank $d$ over $A=\cO_Z(U)$;
\item $M|_{q^{-1}(U)}$ is free.
\end{enumerate}
\end{lemma}

\begin{proof}
Since $q$ is affine and $q_*\cO_W$ is locally free, there is an affine
open $U_1\ni p$ with $q^{-1}(U_1)$ affine and $q_*\cO_W|_{U_1}$ free
of rank $d$.  The fibre $q^{-1}(p)$ is a finite set of points
$\eta_1,\dots,\eta_n$ of $W$, because $q$ is finite.  The semilocal
ring $\cO_{W,\{\eta_1,\dots,\eta_n\}}$ has finitely many maximal ideals,
so every finitely generated projective module of rank one over it is free;
in particular $M$ admits a generator $\tau$ over that semilocal ring.
A generator is defined on some open $V\subseteq q^{-1}(U_1)$ containing
all the $\eta_i$, and the locus where it generates $M$ is open and
contains $\{\eta_1,\dots,\eta_n\}$; shrink $V$ to that locus.

Now use that $q$ is finite, hence closed: the set
$q\bigl(q^{-1}(U_1)\setminus V\bigr)$ is closed in $U_1$ and does not
contain $p$.  Let $U_2$ be its open complement in $U_1$; then
$q^{-1}(U_2)\subseteq V$, so $M$ is free on $q^{-1}(U_2)$.  Finally
choose an affine open $U\subseteq U_2$ containing $p$ with $A$ such
that $q_*\cO_W|_U$ is free; then $q^{-1}(U)$ is affine and both
conditions hold.
\end{proof}

\begin{remark}[Why a frame on $W$ and not on $Z$]
\label{rem:ch41-frame-side}
The frame produced by \Cref{lem:ch41-local-frame} lives upstairs, on
$q^{-1}(U)$, and this is essential: the sections $s_j$ are sections of a
bundle on $W$, and the local element-norm formula requires a trivialization of $M$ on
a full preimage.  The determinant section itself is globally defined
before choosing any frame.  A frame near only one point $\eta_i$ does not trivialize the full
finite algebra needed for this local calculation.  This is the precise technical reason why
\Cref{lem:ch41-local-frame} has to invoke the semilocal ring at the whole
fibre rather than a local ring at one point, and it is the local preparation for the symmetric, globally defined norm of
\Cref{sec:ch41-naive}.
\end{remark}

\subsection{Two explicit norm polynomials}
\label{sec:ch41-norm-examples}

\begin{example}[The $p$-th power map]
\label{ex:ch41-cyclic-degree-p}
Let $\operatorname{char}k=p>0$, $A=k[x]$ and
$B=A[u]/(u^p-x)$, which is free of rank $d=p$ with basis
$1,u,\dots,u^{p-1}$.  Multiplication by $u$ has, in this basis, the
companion matrix of $T^p-x$, so its characteristic polynomial is
$T^p-x$ and
\[
 \mathrm N_{B/A}(u)=(-1)^p\cdot(-x)=x
\]
in every characteristic $p$, the two signs cancelling for odd $p$ and
being irrelevant for $p=2$.  For the universal linear form, its $p$-th power is
$T_0^p+xT_1^p$.  Over the generic field, multiplication by that form
has characteristic polynomial
$T^p-(T_0^p+xT_1^p)$; its determinant is therefore
\[
 \mathrm N_{B/A}(T_0+T_1u)=T_0^p+xT_1^p .
\]
The identity extends to $A[\mathbf T]$ because both sides are
polynomials over the domain $A$.  Every mixed coefficient vanishes.
In particular the coefficient of $T_0^{p-1}T_1$, which is
$\Tr(u)$, is zero.  The surviving extreme coefficients are exactly
$1$ and $x$, with no sign ambiguity in characteristic $p$.

The geometric reading is that $\Spec B\to\Spec A$ is a homeomorphism ---
it is the relative Frobenius --- so every fibre has a single point, of
weight $p$.  Indeed, over the residue field at $x=a$, the fibre algebra
is $k[u]/(u^p-a)=k[u]/((u-a^{1/p})^p)$, local of length $p$, and
\Cref{lem:ch41-element-norm}\,\cref{itm:ch41-en:zero} gives
$\mathrm N(b)=\bar b^{\,p}$.  This is the algebra behind
\Cref{ex:ch41-frobenius-average}.
\end{example}

\begin{example}[A cubic with a nonreduced fibre]
\label{ex:ch41-cubic-norm}
Let $A=k[x]$ and $B=k[u]$, with $x=u^3-u^2$; then $B$ is free over
$A$ with basis $1,u,u^2$, the relation being $u^3=x+u^2$.  In this
basis multiplication by $u$ has matrix
\[
 m_u=\begin{pmatrix}0&0&x\\1&0&0\\0&1&1\end{pmatrix},
\]
whose columns are the coordinates of $u\cdot1=u$, $u\cdot u=u^2$ and
$u\cdot u^2=x+u^2$.  Hence
\[
 \mathrm N_{B/A}(T_0+T_1u)
 =\det\begin{pmatrix}
 T_0&0&xT_1\\ T_1&T_0&0\\ 0&T_1&T_0+T_1
 \end{pmatrix}
 =T_0^3+T_0^2T_1+x\,T_1^3 .
\tag{41.19}\label{eq:ch41-cubic-norm}
\]
In the notation of \eqref{eq:ch41-universal-coefficients} the coefficients
are
\[
 n_{(3,0)}=1,\quad n_{(2,1)}=1,\quad n_{(1,2)}=0,\quad n_{(0,3)}=x .
\]
Setting $T_0=0,T_1=1$ recovers $\mathrm N_{B/A}(u)=x$, as it must, the
minimal polynomial of $u$ over $k(x)$ being $T^3-T^2-x$.

The fibre over $x=0$ is
\[
 B\otimes_Ak[x]/(x)=k[u]/(u^3-u^2)=k[u]/(u^2)\times k ,
\]
a product of local Artinian algebras of lengths $2$ and $1$, the first
of which is nonreduced; the lengths sum to $3=d$.  This example is carried
to explicit values of the averaging operator in
\Cref{ex:ch41-nonreduced-fibre}, and it is the smallest case in which a
nonreduced fibre factor and a reduced one occur over the same point.
\end{example}

\section{Weighted averaging}
\label{sec:ch41-averaging}

We can now define the averaging operator and prove the fibre formula.  The
section is organized so that the two hypotheses the plan forbids ---
separability and reducedness --- can be seen never to be used: the fibre
algebra is treated as an arbitrary finite algebra over a complete valued
field, and the only invariants extracted from it are dimensions of local
factors.

\subsection{The fibre algebra at an analytic point}
\label{sec:ch41-fibre-algebra}

\begin{definition}[Fibre algebra and completed residue field]
\label{def:ch41-fibre-algebra}
Let $q:W\to Z$ be finite locally free of rank $d>0$ and let
$z\in Z^{\An}$, written $(\xi,|\cdot|_z)$ as in
\Cref{lem:ch41-point-pair}.  Write
\[
 \mathcal H(z):=\text{the completion of }\kappa(\xi)
 \text{ for }|\cdot|_z ,
\]
a complete valued field, and
\[
 B_z:=\bigl(q_*\cO_W\bigr)_\xi\otimes_{\cO_{Z,\xi}}\kappa(\xi),
 \qquad
 B_z^{\wedge}:=B_z\otimes_{\kappa(\xi)}\mathcal H(z),
\tag{41.20}\label{eq:ch41-fibre-algebra}
\]
the scheme-theoretic fibre of $q$ over $\xi$ and its completed form.
Both are finite of dimension $d$ over their base fields.
\end{definition}

Note that $|\cdot|_z$ is allowed to be trivial, in which case
$\mathcal H(z)=\kappa(\xi)$; a trivially valued field is complete.  The
algebra $B_z$ is the scheme-theoretic fibre with all its nilpotents.
It is reduced on the étale locus.  Completion may acquire nilpotents
in an inseparable extension, so full support does not rule them out.

\begin{lemma}
\label{lem:ch41-fibre-decomposition}
With the notation of \Cref{def:ch41-fibre-algebra}:
\begin{enumerate}
\item\label{itm:ch41-fd:decomp}
      $B_z^{\wedge}=\prod_{i=1}^{n}C_i$ with each $C_i$ a local
      Artinian $\mathcal H(z)$-algebra, finite over
      $\mathcal H(z)$, with residue field $K_i$ a finite extension of
      $\mathcal H(z)$; and
      $\sum_i\dim_{\mathcal H(z)}C_i=d$.
\item\label{itm:ch41-fd:unique}
      Each $K_i$ carries a unique absolute value extending
      $|\cdot|_z$; its restriction along
      $\kappa(\eta_i)\hookrightarrow K_i$, where $\eta_i\in W$ is the
      point of the fibre corresponding to $C_i$, defines a point
      $w_i\in W^{\An}$ with $q^{\An}(w_i)=z$.
\item\label{itm:ch41-fd:onto}
      Every $w\in W^{\An}$ with $q^{\An}(w)=z$ equals $w_i$ for at
      least one $i$; in particular the fibre of $q^{\An}$ over $z$ is
      finite and nonempty.
\end{enumerate}
Define, for $w\in(q^{\An})^{-1}(z)$,
\[
 e(w/z):=\sum_{i\,:\,w_i=w}\dim_{\mathcal H(z)}C_i
 \ \in\mathbf Z_{\ge1},
 \qquad\text{so that}\qquad
 \sum_{w\mid z}e(w/z)=d .
\tag{41.21}\label{eq:ch41-weights}
\]
\end{lemma}

\begin{proof}
\cref{itm:ch41-fd:decomp}  A finite algebra over a field is Artinian, hence a
finite product of local Artinian algebras; dimensions add.  Finiteness of
$\dim_{\mathcal H(z)}B_z^{\wedge}=d$ is
\Cref{def:ch41-fibre-algebra}, local freeness of rank $d$ being preserved
by the base changes $\cO_{Z,\xi}\to\kappa(\xi)\to\mathcal H(z)$.

\cref{itm:ch41-fd:unique}  A finite extension of a complete non-Archimedean
valued field carries exactly one absolute value extending the given one;
this is the classical uniqueness theorem, valid also when the given absolute
value is trivial, in which case every extension is trivial because every
element of an algebraic extension of a trivially valued field has absolute
value $1$.  See \cite[Chapter~VI, \S11]{ZariskiSamuelII}.  The point
$\eta_i\in W$ is the image of $\Spec C_i\to\Spec B_z\to W$, and
$\kappa(\eta_i)$ is a subfield of $K_i$ whose compositum with
$\mathcal H(z)$ is dense in $K_i$; restricting the unique absolute value
of $K_i$ gives an absolute value on $\kappa(\eta_i)$ restricting to
$|\cdot|_z$ on $\kappa(\xi)$, that is, a point $w_i\in W^{\An}$ with
$q^{\An}(w_i)=z$ by
\Cref{lem:ch41-support-centre-finite}\,\cref{itm:ch41-sc:supp}.

\cref{itm:ch41-fd:onto}  Let $w=(\eta,|\cdot|_w)$ lie over $z$.
The inclusion $\kappa(\xi)\to\kappa(\eta)$ is isometric, hence
extends to $\mathcal H(z)\to\mathcal H(w)$.
There is a map
$B_z^\wedge\to\mathcal H(w)$ induced by evaluation at $w$.
Its image is a finite-dimensional domain over $\mathcal H(z)$,
so it is a field.  This field is complete, as a finite extension
of a complete valued field, and contains the dense image of
$\kappa(\eta)$ in $\mathcal H(w)$.  Therefore its image is all
of $\mathcal H(w)$.  Its kernel is a maximal ideal belonging to
one factor $C_i$, whose residue field is thus $\mathcal H(w)$.
Uniqueness of the extended absolute value gives $w_i=w$.
Nonemptiness follows from $B_z^\wedge\ne0$.

Finally \eqref{eq:ch41-weights} follows: the weights $e(w/z)$ are positive
because each $C_i$ has positive dimension, and their sum is
$\sum_i\dim_{\mathcal H(z)}C_i=d$ by \cref{itm:ch41-fd:decomp} together
with \cref{itm:ch41-fd:onto}, which guarantees that every index $i$ is
counted in exactly one weight.
\end{proof}

\begin{remark}[What the weight is made of]
\label{rem:ch41-weight-anatomy}
The single integer $e(w/z)$ packages several different phenomena, and it
is worth separating them once.  Write $C$ for a local factor with residue
field $K$, and write $\ell_C(C)$ for its composition length as a
module over itself.  Then the unconditional identity is
\[
 \dim_{\mathcal H(z)}C
 =\ell_C(C)\,[K:\mathcal H(z)].
\]
If the valued-field extension is defectless, the second factor is the
product of its ramification index and residue degree.  In general a defect
factor remains, and it must not be suppressed.  Thus the weight records the
Artin length, the full valued-field degree---including ramification, residue
degree, inseparability and possible defect---and any coincidence of local
factors that induce the same analytic lift.  An unweighted count of
analytic lifts loses field degrees, while an unweighted count after
algebraic closure loses Artin lengths whenever the fibre is nonreduced.  In the running examples:
\Cref{ex:ch41-frobenius-average} has
length $1$ and purely inseparable \emph{ramification} of index $p$ at
divisorial points --- there $e=p$ and $f=1$, the value group being
divided by $p$ while the residue field does not grow --- and length
$p$ at trivially valued closed points;
\Cref{ex:ch41-nonreduced-fibre} has a factor of length $2$ over a
trivially valued point and a ramification index $2$ over a divisorial
point.  In both cases the total is $d$.
\end{remark}

\begin{warning}{The weights are not locally constant in $z$}
\label{warn:ch41-weights-not-constant}
The multiset of fibre weights need not be locally constant.
In \Cref{ex:ch41-nonreduced-fibre}, choose a finite $a$ outside the
branch values.  The points $z_t=t\ord_{x=a}$ have three lifts of
weight $1$, and approach $v_{Z,\triv}$ as $t\downarrow0$.
At that trivial generic point there is one lift, of weight $3$.
The point $t\ord_{x=0}$ instead has weights $2,1$; it is not
approached by the other rays at a fixed positive parameter.
Only the total degree is constant.  The extended-real continuity of
an averaged FS stage follows from its coefficient presentation and
\Cref{prop:ch34-generic-finiteness}\,\cref{itm:ch34-gf:cont},
not from a continuous labelling of individual lifts.
\end{warning}

\subsection{Gauss norms}
\label{sec:ch41-gauss}

\begin{lemma}
\label{lem:ch41-gauss-multiplicative}
Let $(\Omega,|\cdot|)$ be a field with a multiplicative non-Archimedean
absolute value, let $\rho_0,\dots,\rho_r>0$, and define on
$\Omega[T_0,\dots,T_r]$
\[
 \Bigl\|\sum_Ia_IT^I\Bigr\|_\rho:=\max_I|a_I|\,\rho^I,
 \qquad
 \rho^I:=\prod_j\rho_j^{I_j}.
\tag{41.22}\label{eq:ch41-gauss-norm}
\]
Then $\|PQ\|_\rho=\|P\|_\rho\|Q\|_\rho$ for all
$P,Q\in\Omega[T_0,\dots,T_r]$, including when one of them is $0$.
\end{lemma}

\begin{proof}
The inequality $\le$ is the ultrametric inequality applied to each
coefficient of $PQ$, which is a sum of products $a_Ib_J$ with
$I+J$ fixed.  For $\ge$ we may assume $P,Q\ne0$.  Fix a monomial order
$\preceq$ on $\mathbf N^{r+1}$, that is, a total order compatible with
addition; the lexicographic order will do.  Let $I_0$ be the
$\preceq$-largest multi-index among those attaining
$\max_I|a_I|\rho^I$, and let $J_0$ be defined analogously for $Q$.
If $I+J=I_0+J_0$ with $(I,J)\ne(I_0,J_0)$ and both $I$ and $J$
attain the respective maxima, then $I\preceq I_0$ and $J\preceq J_0$
with at least one strict, whence $I+J\prec I_0+J_0$, a contradiction.
Therefore in the coefficient
$\sum_{I+J=I_0+J_0}a_Ib_J$ of $T^{I_0+J_0}$ the term
$a_{I_0}b_{J_0}$ has strictly larger absolute value than all the others,
so the coefficient has absolute value $|a_{I_0}||b_{J_0}|$, and
\[
 \|PQ\|_\rho\ge|a_{I_0}b_{J_0}|\rho^{I_0+J_0}
 =\|P\|_\rho\|Q\|_\rho .
\]
The case $P=0$ or $Q=0$ is trivial, both sides being $0$.
\end{proof}

\begin{remark}[Why a Gauss norm and not an evaluation]
\label{rem:ch41-why-gauss}
The variables $T_j$ are formal and are never specialized.  What
\Cref{lem:ch41-gauss-multiplicative} converts is a factorization of the
\emph{polynomial} $\mathrm N(\sum_jT_js_j)$ into a factorization of a
single real number, the Gauss norm with radii
$\rho_j=\rme^{\lambda_j}$.  Specializing the $T_j$ to elements of
$\Omega$ would lose exactly the information that the maximum over $j$ is
being taken: the value of a linear form at a point can be much smaller than
the largest of its coefficients, whereas its Gauss norm is exactly the
largest.  This is the technical reason for introducing the universal
variables, and it explains why the weights $\lambda_j$ enter as radii.
\end{remark}

\subsection{The averaging operator}
\label{sec:ch41-average-operator}

\begin{proposition}\label{prop:bj-weighted-coefficient-norm}
Assume that $q:W\to Z$ is finite faithfully flat of rank $d$, with
$W$ integral and $W,Z$ projective, and with $Z$ normal.  Let $M$ be a line bundle on
$W$, and let
\[
 \varphi=\frac1m\max_j\{\log|s_j|+\lambda_j\},
 \qquad
 s_j\in H^0(W,M^{\otimes m}),\quad
 \lambda_j\in\mathbf Q,
\tag{41.23}\label{eq:bj-weighted-fs-upstairs}
\]
be rational generically finite Fubini--Study data.

Introduce independent variables $T_j$.  The universal norm is a
homogeneous polynomial law of degree $d$:
\[
 \mathrm N_q\left(\sum_jT_js_j\right)
     =\sum_{|I|=d}c_IT^I,
 \qquad
 c_I\in H^0\left(
 Z,\mathrm N_q(M^{\otimes m})\right).
\tag{41.24}\label{eq:bj-universal-norm-polynomial}
\]
Define
\[
 \operatorname{Av}_q(\varphi):=
 \frac1{md}\max_{|I|=d}
 \left\{\log|c_I|+\sum_jI_j\lambda_j\right\}.
\tag{41.25}\label{eq:bj-weighted-average-definition}
\]
Then $\operatorname{Av}_q(\varphi)$ is a generically finite rational FS function, and
for every $z\in Z^{\An}$,
\[
 \operatorname{Av}_q(\varphi)(z)
 =\frac1d\sum_{w\mid z}e(w/z)\varphi(w),
 \qquad
 \sum_{w\mid z}e(w/z)=d.
\tag{41.26}\label{eq:bj-finite-flat-average}
\]
The positive integers $e(w/z)$ include residue, inseparability, and local
Artin-length multiplicities.  The numerical class of this presentation is
\[
 [\operatorname{Av}_q(\varphi)]
   =\frac1d c_1(\mathrm N_qM)
   =\frac1d q_*c_1(M).
\tag{41.27}\label{eq:bj-norm-class}
\]
The operation preserves order, commutes with homogeneous scaling whenever the weights stay in the
specified group, and
sends decreasing GF nets to decreasing GF nets.  These statements remain
valid for inseparable and nonreduced fibres and at points where the value
is $-\infty$.
\end{proposition}

\begin{proof}
\noindent\emph{Step 1: the polynomial law and the coefficient sections.}
The determinant in \eqref{eq:bj-universal-norm-polynomial} is homogeneous
of degree $d$ in the $T_j$.  Formation of the determinant commutes with
base change, so the coefficient sections $c_I$ glue globally.  This
coefficient construction is the universal polynomial-law form of the norm;
it avoids selecting or ordering geometric sheets.

In detail: cover $Z$ by affine opens $U$ as in
\Cref{lem:ch41-local-frame}, so that $q^{-1}(U)=\Spec B$ with $B$ free
of rank $d$ over $A=\cO_Z(U)$ and $M^{\otimes m}$ is trivialized on
$q^{-1}(U)$ by a frame $\tau$.  Writing $s_j=f_j\tau$ with
$f_j\in B$ and applying
\Cref{prop:ch41-norm-line-properties}\,\cref{itm:ch41-nl:frame},
\[
 \mathrm N_q\Bigl(\sum_jT_js_j\Bigr)
 =\mathrm N_{B/A}\Bigl(\sum_jT_jf_j\Bigr)\cdot\mathrm N_q(\tau)
 =\Bigl(\sum_{|I|=d}n_IT^I\Bigr)\mathrm N_q(\tau),
\tag{41.28}\label{eq:ch41-local-coefficients}
\]
with $n_I\in A$ given by \Cref{lem:ch41-universal-polynomial} and
$\mathrm N_q(\tau)$ a frame of $\mathrm N_q(M^{\otimes m})$ on $U$.
Setting $c_I:=n_I\,\mathrm N_q(\tau)$ on $U$, two such descriptions
agree on overlaps: if $\tau'=\upsilon\tau$ with $\upsilon\in B^\times$, then
$f_j'=\upsilon^{-1}f_j$.  Thus $n_I'=\mathrm N_{B/A}(\upsilon)^{-1}n_I$
and $\mathrm N_q(\tau')=\mathrm N_{B/A}(\upsilon)\mathrm N_q(\tau)$, by
\Cref{prop:appA-norm-identities} and
\Cref{lem:ch41-element-norm}\,\cref{itm:ch41-en:unit}.  Therefore the
$c_I$ are well defined global sections of
$\mathrm N_q(M^{\otimes m})$, and
\eqref{eq:bj-universal-norm-polynomial} holds.

\medskip
\noindent\emph{Step 2: $\operatorname{Av}_q(\varphi)$ is Fubini--Study data.}
By
\Cref{prop:ch41-norm-line-properties}\,\cref{itm:ch41-nl:mult} there is a
canonical isomorphism
$\mathrm N_q(M^{\otimes m})\cong\mathrm N_q(M)^{\otimes m}$.  Put
\[
 m':=md,
 \qquad
 L':=\tfrac1d\,\mathrm N_q(M)\ \text{ as a }\mathbf Q\text{-line bundle},
 \qquad
 \lambda'_I:=\sum_jI_j\lambda_j\in\mathbf Q .
\]
Then $m'L'=m\,\mathrm N_q(M)$, so $c_I\in H^0(Z,m'L')$ and
\eqref{eq:bj-weighted-average-definition} reads
\[
 \operatorname{Av}_q(\varphi)
 =\frac1{m'}\max_{|I|=d}\bigl\{\log|c_I|+\lambda'_I\bigr\},
\]
which is a weighted section system in the sense of
\eqref{eq:gf-fs-presentation}, with rational weights.  It therefore defines
an element of $\cH^{\mathrm{gf}}_{\mathbf Q}(L')$ as soon as the $c_I$
do not all vanish identically, which is proved in Step~5.

\medskip
\noindent\emph{Step 3: the fibre calculation.}
Fix $z\in Z^{\An}$, written $(\xi,|\cdot|_z)$, and adopt the notation
of \Cref{def:ch41-fibre-algebra} and
\Cref{lem:ch41-fibre-decomposition}.  Choose $U\ni c(z)$ as in Step~1;
this is legitimate because $\xi$ specializes to $c(z)$, so
$\xi\in U$, and because by
\Cref{lem:ch41-support-centre-finite}\,\cref{itm:ch41-sc:centre} the centre
of every $w$ over $z$ lies in $q^{-1}(c(z))\subseteq q^{-1}(U)$.
With this choice, for every $w$ over $z$,
\[
 \log|s_j|(w)=-w(f_j),
 \qquad
 \log|c_I|(z)=-z(n_I).
\tag{41.29}\label{eq:ch41-frame-values}
\]
Let $\Omega$ be an algebraic closure of $\mathcal H(z)$, equipped with
an absolute value extending that of $\mathcal H(z)$; such an extension
exists by \cite[Chapter~VI, \S4]{ZariskiSamuelII}, and the inclusion
$\kappa(\xi)\hookrightarrow\Omega$ is isometric by construction.  By
\Cref{prop:ch41-norm-base-change}, applied to the composite base change
$A\to\kappa(\xi)\to\mathcal H(z)\to\Omega$, the image of
\eqref{eq:ch41-local-coefficients} in $\Omega[\mathbf T]$ is the norm
polynomial of the $\Omega$-algebra
\[
 B\otimes_A\Omega=B_z\otimes_{\kappa(\xi)}\Omega
 =B_z^{\wedge}\otimes_{\mathcal H(z)}\Omega
 =\prod_{i=1}^n\ \prod_{l}D_{il},
\]
where $C_i\otimes_{\mathcal H(z)}\Omega=\prod_lD_{il}$ with each
$D_{il}$ local Artinian with residue field $\Omega$, $\Omega$ being
algebraically closed.  By
\Cref{lem:ch41-element-norm}\,\cref{itm:ch41-en:product} the determinant is
the product of the block determinants, and
a composition filtration of $D_{il}$ has
$\dim_\Omega D_{il}$ quotients isomorphic to $\Omega$.
Tensoring it with $\Omega[\mathbf T]$, multiplication by
$\sum_jT_jf_j$ acts on every quotient by the same residue linear form.
The resulting block-triangular determinant is therefore
\[
 \Bigl(\sum_jT_j\,\varepsilon_{il}(f_j)\Bigr)^{\dim_\Omega D_{il}},
\]
where $\varepsilon_{il}:B\otimes_A\Omega\to D_{il}\to\Omega$ is the
residue map.  Thus in $\Omega[\mathbf T]$
\[
 \sum_{|I|=d}n_I\,T^I
 =\prod_{i=1}^n\prod_l
 \Bigl(\sum_jT_j\,\varepsilon_{il}(f_j)\Bigr)^{\dim_\Omega D_{il}} .
\tag{41.30}\label{eq:ch41-norm-factorization}
\]
The residue map induced on $C_i$ factors through its residue
field $K_i$ and an embedding $K_i\hookrightarrow\Omega$
over $\mathcal H(z)$.  No coefficient-field inclusion
$K_i\subset C_i$ is being chosen.  Since $K_i$ carries
a \emph{unique} absolute value extending that of $\mathcal H(z)$, by
\Cref{lem:ch41-fibre-decomposition}\,\cref{itm:ch41-fd:unique}, every such
embedding is isometric; consequently
\[
 |\varepsilon_{il}(f_j)|_\Omega=|f_j|_{w_i}=\rme^{\log|s_j|(w_i)}
\]
is independent of $l$.  Give $\Omega[\mathbf T]$ the weighted Gauss
norm \eqref{eq:ch41-gauss-norm} with radii
$\rho_j:=\rme^{\lambda_j}$.  The Gauss norm of the linear form
$\sum_jT_j\varepsilon_{il}(f_j)$ is
\[
 \max_j|f_j|_{w_i}\rme^{\lambda_j}
 =\exp\Bigl(\max_j\{\log|s_j|(w_i)+\lambda_j\}\Bigr)
 =\rme^{\,m\varphi(w_i)} ,
\]
again independent of $l$.  Applying
\Cref{lem:ch41-gauss-multiplicative} to
\eqref{eq:ch41-norm-factorization} and using
$\sum_l\dim_\Omega D_{il}=\dim_{\mathcal H(z)}C_i$,
\[
 \max_{|I|=d}|n_I|_\Omega\,\rho^I
 =\prod_{i=1}^n\rme^{\,m\varphi(w_i)\dim_{\mathcal H(z)}C_i}
 =\prod_{w\mid z}\rme^{\,m\,e(w/z)\varphi(w)} ,
\]
the last equality being the grouping \eqref{eq:ch41-weights}.  Taking
logarithms, using $|n_I|_\Omega=|n_I|_z$ and
\eqref{eq:ch41-frame-values},
\[
 \max_{|I|=d}
 \left\{\log|c_I(z)|+\sum_jI_j\lambda_j\right\}
 =
 \sum_{w\mid z}e(w/z)
 \max_j\{\log|s_j(w)|+\lambda_j\}.
\tag{41.31}\label{eq:bj-norm-gauss-identity}
\]
The norm polynomial and its coefficients commute with the isometric base
extension, so this equality descends to the original completed residue
field.  Division by $md$ proves
\eqref{eq:bj-finite-flat-average}.

\medskip
\noindent\emph{Step 4: common zeros, nilpotents, and the value $-\infty$.}
If all $s_j$ vanish at one lift $w$, the corresponding linear factor
is the zero polynomial.  Hence the entire determinant polynomial is zero
at $z$, all $c_I(z)$ vanish, and both sides of
\eqref{eq:bj-norm-gauss-identity} equal $-\infty$.  If no local factor
vanishes, their product is nonzero.  Nilpotents change the factor length
but not this conclusion.

It is worth being explicit, since this is the extended-real point of the
proof.  \Cref{lem:ch41-gauss-multiplicative} was stated to include the zero
polynomial, so \eqref{eq:ch41-norm-factorization} and the Gauss-norm
identity hold verbatim in $\Rbar$ with the convention
$\log0=-\infty$, and no case distinction is needed.  Both sides of
\eqref{eq:bj-norm-gauss-identity} are then $-\infty$ exactly when some
factor vanishes, that is, exactly when $\varphi(w)=-\infty$ for some
$w$ over $z$.  In particular the polar locus of
$\operatorname{Av}_q(\varphi)$ is the image under $q^{\An}$ of the polar
locus of $\varphi$, which is closed by
\Cref{lem:ch41-finite-surjective} and consistent with
\Cref{prop:ch34-generic-finiteness}\,\cref{itm:ch34-gf:polar}.  No
subtraction of infinities occurs anywhere: the right-hand side of
\eqref{eq:bj-finite-flat-average} is a finite sum of terms in
$[-\infty,+\infty)$ with positive coefficients, which is unambiguous by
ledger entry~(2) of \Cref{sec:extended-arithmetic}.

No trace average and no Reynolds operator occurs here.  In particular,
when the characteristic divides $d$, mixed coefficients may vanish; the
factor $1/d$ is used only in the real-valued weight and the numerical
class.  Thus neither separability nor geometric reducedness is required.

\medskip
\noindent\emph{Step 5: class, generic finiteness, order, and scaling.}
For a rational section $s$ of $M$, the determinant calculation gives
\[
 \operatorname{div}\mathrm N_q(s)=q_*\operatorname{div}(s),\qquad
 c_1(\mathrm N_qM)=q_*c_1(M).
\tag{41.32}\label{eq:bj-norm-divisor}
\]
This is
\Cref{prop:ch41-norm-line-properties}\,\cref{itm:ch41-nl:divisor}.  Since
$\mathrm N_q(M^{\otimes m})\simeq\mathrm N_q(M)^{\otimes m}$, the class of
\eqref{eq:bj-weighted-average-definition} is
\[
 \frac1{md}c_1\left(\mathrm N_q(M^{\otimes m})\right)
 =\frac1dq_*c_1(M).
\]
At the generic point, the universal linear form is nonzero in
$K(W)[\mathbf T]$; its field norm is nonzero.  Hence some $c_I$ is
nonzero, proving generic finiteness.  In detail: $Z$ is integral and
$q$ is finite flat, so the generic fibre $B\otimes_AK(Z)$ is a
$d$-dimensional $K(Z)$-algebra contained in $K(W)$, hence equal to
$K(W)$; the element $\sum_jT_jf_j$ of $K(W)[\mathbf T]$ is nonzero
because $\varphi$ is generically finite, so some $f_j$ is nonzero; and
multiplication by a nonzero element of the domain $K(W)[\mathbf T]$ is an
injective endomorphism of a free $K(Z)[\mathbf T]$-module of rank $d$,
hence has nonzero determinant.  Thus some $n_I\ne0$ in $K(Z)$, and
\Cref{prop:ch34-generic-finiteness}\,\cref{itm:ch34-gf:equiv} applies.

Order preservation and compatibility with decreasing nets follow from the
finite sum in \eqref{eq:bj-finite-flat-average}.  Finite monotone
convergence remains exact when a summand tends to $-\infty$.  Homogeneous
scaling preserves support kernels, permutes the lifts over a scaled point,
and preserves all local Artin lengths.  The same fibre formula proves
compatibility with scaling.

Once more in detail.  Formula \eqref{eq:bj-finite-flat-average} shows that
$\operatorname{Av}_q(\varphi)$ depends only on $\varphi$ as a function
on $W^{\An}$, not on the chosen presentation
\eqref{eq:bj-weighted-fs-upstairs}; so $\operatorname{Av}_q$ is a well
defined operator on rational generically finite Fubini--Study functions,
and $\varphi\le\psi$ implies
$\operatorname{Av}_q(\varphi)\le\operatorname{Av}_q(\psi)$, the weights
being positive.  If $(\varphi_i)$ is a decreasing net with pointwise
limit $\varphi_\infty$, then for each fixed $z$ the finite sum
$\frac1d\sum_{w\mid z}e(w/z)\varphi_i(w)$ decreases to
$\frac1d\sum_{w\mid z}e(w/z)\varphi_\infty(w)$, by monotone convergence
for finite sums in $[-\infty,+\infty)$.  For the scaling action, let
$t>0$; the homeomorphism $v\mapsto tv$ of
\Cref{prop:ch32-scaling-action} raises every absolute value to the power
$t$, leaving the support prime and hence the fibre algebra $B_z$
unchanged, and leaving the topology of $\kappa(\xi)$ and therefore the
completion $\mathcal H(z)$ unchanged as a field.  The decomposition of
\Cref{lem:ch41-fibre-decomposition} is thus unaffected, so $w\mapsto tw$
is a bijection from the lifts of $z$ to the lifts of $tz$ preserving the
weights.  With the action on functions normalized as in
\eqref{eq:ch32-scaling-on-functions} and its calculus
\Cref{lem:ch35-action-basic}, this gives
\[
 \operatorname{Av}_q(t\cdot\varphi)(z)
 =\frac1d\sum_{w\mid z}e(w/z)\,t\,\varphi(t^{-1}w)
 =t\,\operatorname{Av}_q(\varphi)(t^{-1}z)
 =\bigl(t\cdot\operatorname{Av}_q(\varphi)\bigr)(z),
\]
which is
the asserted equivariance.  For rational-weight FS data this stays
within the domain for positive rational $t$.  The real-weight extension
in \Cref{cor:ch41-real-weights} permits all $t>0$.
\end{proof}

For any function $f:W^{\An}\to\mathbf R\cup\{-\infty\}$, the finite
fibre formula also defines $\operatorname{Av}_q(f)$.
This extension is an operation on functions; it asserts no FS or psh
membership without the hypotheses above.  It lets us use the same
notation for the limits and pullback identities below.

\begin{corollary}
\label{cor:ch41-average-pullback}
In the setting of \Cref{prop:bj-weighted-coefficient-norm}, let
$\psi:Z^{\An}\to\mathbf R\cup\{-\infty\}$ be any function and suppose
$\varphi=q^*\psi$ on the fibre of $q^{\An}$ over a point $z$.  Then
\[
 \operatorname{Av}_q(\varphi)(z)=\psi(z).
\]
Consequently, if $\varphi=q^*\psi$ everywhere, then
$\operatorname{Av}_q(\varphi)=\psi$; in particular
$\operatorname{Av}_q\circ\,q^*=\operatorname{id}$ on functions.
\end{corollary}

\begin{proof}
Every $w$ over $z$ has $\varphi(w)=\psi(z)$, so
\eqref{eq:bj-finite-flat-average} gives
\[
 \operatorname{Av}_q(\varphi)(z)
 =\frac1d\psi(z)\sum_{w\mid z}e(w/z)
 =\psi(z),
\]
using $\sum_we(w/z)=d$.  The identity is valid also when
$\psi(z)=-\infty$, the sum then being $-\infty$.
\end{proof}

\begin{remark}[The two roles of $\sum_we(w/z)=d$]
\label{rem:ch41-two-roles-of-d}
The identity $\sum_we(w/z)=d$ is used twice, in genuinely different ways.
In \Cref{cor:ch41-average-pullback} it is what makes the averaging operator
a left inverse of pullback, and hence what turns
\Cref{prop:bj-weighted-coefficient-norm} into a descent statement.  In
\eqref{eq:bj-norm-class} its algebraic counterpart, the constancy of the
rank, is what makes the numerical class of the averaged function equal to
$\frac1dq_*c_1(M)$ and therefore what makes
\eqref{eq:ch41-projection-formula} applicable in
\Cref{sec:ch41-descent}.  If the weights were taken to be numbers of
distinct points without their required degrees and lengths, neither use
would be justified in general.  The total could differ from $d$, and
the resulting unweighted formula would no longer be the coefficient-norm
construction that supplies the stated presentation class.
\end{remark}

\begin{corollary}\label{cor:bj-coefficient-norm-limit}
In the setting of \Cref{prop:bj-weighted-coefficient-norm}, let
$\varphi_j\downarrow\varphi$ be a decreasing net of generically
finite GF functions with real weights on $W$, and assume
\[
             \frac1dq_*[\varphi_j]\longrightarrow
             \eta\in N^1(Z)_{\mathbf R}.
\tag{41.33}\label{eq:bj-norm-limit-classes}
\]
If
\[
 a(z):=\frac1d\sum_{w\mid z}e(w/z)\varphi(w)
\tag{41.34}\label{eq:bj-norm-limit-function}
\]
is not identically $-\infty$, then $a\in\PSH_Z(\eta)$, more
precisely $\operatorname{Av}_q(\varphi_j)\downarrow a$.  If the stages are homogeneous
and vanish at the trivial valuation, then $a$ has the same properties
and is automatically nontrivial.
\end{corollary}

\begin{proof}
\Cref{prop:bj-weighted-coefficient-norm} and
\Cref{cor:ch41-real-weights} give decreasing GF functions
$\operatorname{Av}_q(\varphi_j)$ with real weights and the classes in
\eqref{eq:bj-norm-limit-classes}.  The fibre formula and finite monotone
convergence identify their pointwise limit with
\eqref{eq:bj-norm-limit-function}.  The varying-class definition
\Cref{def:bj-psh} proves the assertion.

We spell out the four verifications required by \Cref{def:bj-psh}.

\emph{The stages are admissible.}  Each $\varphi_j$ is
generically finite Fubini--Study data with real weights for some $\mathbf Q$-line bundle
$M_j$ on $W$, and \Cref{prop:bj-weighted-coefficient-norm}, Step~2 and
Step~5, produces
$\operatorname{Av}_q(\varphi_j)\in\cH^{\mathrm{gf}}_{\mathbf R}(L'_j)$
with $L'_j=\frac1d\mathrm N_q(M_j)$ and
$c_1(L'_j)=\frac1dq_*[\varphi_j]$.

\emph{The stages decrease.}  If $j\le j'$ then
$\varphi_j\ge\varphi_{j'}$, so
$\operatorname{Av}_q(\varphi_j)\ge\operatorname{Av}_q(\varphi_{j'})$ by
order preservation.  The index set is the same directed set, so
$(\operatorname{Av}_q(\varphi_j))$ is a decreasing net.

\emph{The classes converge.}  This is the hypothesis
\eqref{eq:bj-norm-limit-classes}.

\emph{The limit is the asserted function and is generically finite.}
For each $z\in Z^{\An}$, \eqref{eq:bj-finite-flat-average} gives
$\operatorname{Av}_q(\varphi_j)(z)=\frac1d\sum_{w\mid z}e(w/z)\varphi_j(w)$,
a finite sum with positive integer coefficients summing to $d$; letting
$j$ increase and using monotone convergence for finite sums in
$[-\infty,+\infty)$ yields $a(z)$ as in
\eqref{eq:bj-norm-limit-function}.  Generic finiteness of $a$ is the
hypothesis that $a\not\equiv-\infty$, together with
\Cref{prop:ch32-maximum-principle}\,\cref{itm:ch32-mp:gf} and the
integrality of $Z$: a decreasing limit of functions that are decreasing
for \eqref{eq:semivaluation-order} is decreasing, so $a$ is
generically finite as soon as $a(v_{Z,\triv})>-\infty$, which holds since
$a\not\equiv-\infty$ and $a$ attains its supremum at
$v_{Z,\triv}$.  Now \Cref{thm:bj-basic-closure}\,(5), the varying-class
closure statement, gives $a\in\PSH_Z(\eta)$.

For the last sentence, suppose the $\varphi_j$ are homogeneous and
$\varphi_j(v_{W,\triv})=0$.  By the scaling equivariance in
\Cref{prop:bj-weighted-coefficient-norm} each
$\operatorname{Av}_q(\varphi_j)$ is homogeneous, hence so is the limit;
and by \eqref{eq:bj-finite-flat-average} applied at $z=v_{Z,\triv}$,
whose unique lift is $v_{W,\triv}$ with weight $d$ --- the generic
point of $W$ is the only point over the generic point of $Z$, and the
trivial absolute value has a unique extension --- we get
$\operatorname{Av}_q(\varphi_j)(v_{Z,\triv})=0$ for every $j$, so
$a(v_{Z,\triv})=0$ and $a$ is in particular not identically
$-\infty$.
\end{proof}

\begin{remark}[Why the corollary is stated with varying classes]
\label{rem:ch41-varying-classes-needed}
The presented classes $\frac1dq_*[\varphi_j]$ need only converge
to $\eta$; they may already equal it in particular cases.
Here $[\varphi_j]$ denotes the class of the chosen line-bundle
presentation, not a uniquely determined class of the function.
Moving-class \Cref{def:bj-psh} directly justifies the passage without
asserting that finite-stage correction is impossible.
A convention restricted to bounded PL approximants in a fixed class
would instead require that class to be nef by
\Cref{lem:ch34-psh-nonempty}.  The descent argument concerns arbitrary
real classes, with the necessary nontriviality checked separately.
\end{remark}

\subsection{Real weights}
\label{sec:ch41-real-weights}

\begin{corollary}
\label{cor:ch41-real-weights}
\Cref{prop:bj-weighted-coefficient-norm} holds verbatim with
$\lambda_j\in\Lambda$ for an arbitrary additive subgroup
$\Lambda\subseteq\mathbf R$; the conclusion is then
$\operatorname{Av}_q(\varphi)\in\cH^{\mathrm{gf}}_\Lambda
\bigl(\tfrac1d\mathrm N_q(M)\bigr)$, and
\eqref{eq:bj-finite-flat-average}, \eqref{eq:bj-norm-class} and the order,
net properties are unchanged.  Scaling preserves the subgroup
domain whenever $t\Lambda\subseteq\Lambda$; for $\Lambda=\mathbf R$
all positive real $t$ are allowed.  In particular the case
$\Lambda=\mathbf R$ holds.
\end{corollary}

\begin{proof}
Inspect the proof.  Step~1 constructs the coefficient sections $c_I$ and
does not mention the weights.  Step~2 uses only that the numbers
$\lambda'_I=\sum_jI_j\lambda_j$ lie in the same additive subgroup as the
$\lambda_j$, which holds for any subgroup, and that
$m'L'=m\,\mathrm N_q(M)$, which is independent of the weights.  Step~3
uses the weights only through the radii $\rho_j=\rme^{\lambda_j}$ in
\Cref{lem:ch41-gauss-multiplicative}, which is stated for arbitrary positive
radii.  Steps~4 and~5 do not use the weights at all, except through the
already established identity \eqref{eq:bj-finite-flat-average}.
\end{proof}

\subsection{Three covers of curves, computed}
\label{sec:ch41-averaging-examples}

The following three examples carry
\eqref{eq:bj-finite-flat-average} to explicit numbers.  All three are
finite morphisms of smooth projective curves, so $q=h$ is already flat by
miracle flatness and no flattening is needed; the point is to exhibit the
weights.  In each case $X=T=\mathbf P^1$, the affine coordinate
downstairs is $x$ and upstairs is $u$, and we use the global weighted system of the two sections of
$\cO(1)$ vanishing at infinity and at zero.  At points whose centre
lies in the finite affine chart, its expression is
\[
 \varphi:=\max\{\lambda_0,\ \log|u|+\lambda_1\},
 \qquad
 \lambda_0,\lambda_1\in\mathbf Q,
\tag{41.35}\label{eq:ch41-example-data}
\]
which is \eqref{eq:bj-weighted-fs-upstairs} with $m=1$, $M=\cO(1)$,
$s_0$ the section cutting out the point at infinity and $s_1$ the
section cutting out $\{u=0\}$; the section $s_0$ is a frame on that chart, so
$\log|s_0|=0$ and $\log|s_1|=\log|u|$ at centres there.
All subsequent formulas using the rational coordinates $u$ and $x$
are evaluated at points whose centres lie in the finite affine chart
downstairs; their lifts have centres in the finite chart upstairs.
At other centres the global norm-section presentation must be used:
the rational function $\log|u|$ alone is not the section value.

\begin{example}[A separable double cover: the average, checked]
\label{ex:ch41-separable-average}
Let $\operatorname{char}k\ne2$ and $h(u)=u^2=x$, so $d=2$ and
$B=A[u]/(u^2-x)$ with $A=k[x]$.  Take the presentation
\eqref{eq:ch41-example-data} with $s_1$ replaced by the section vanishing
at $u=1$, that is
\[
 \varphi=\max\{0,\ \log|u-1|+\lambda\},
\]
so that the interesting behaviour occurs over $x=1$.  From
\Cref{ex:ch41-rank-two} with the basis $1,u$, and writing
$u-1=-1+u$,
\[
 \mathrm N_{B/A}(T_0+T_1(u-1))
 =\det\begin{pmatrix}T_0-T_1&xT_1\\ T_1&T_0-T_1\end{pmatrix}
 =(T_0-T_1)^2-xT_1^2 .
\]
Expanding, the coefficients are
\[
 n_{(2,0)}=1,\qquad n_{(1,1)}=-2,\qquad n_{(0,2)}=1-x .
\]
Now evaluate at $z=t\ord_{x=1}$ with $t>0$.  Here
$\log|1|=0$, $\log|-2|=0$ since $\operatorname{char}k\ne2$ and $k$
is trivially valued, and $\log|1-x|(z)=-t$.  Therefore
\[
 \operatorname{Av}_q(\varphi)(z)
 =\tfrac12\max\{0,\ \lambda,\ 2\lambda-t\}.
\]
On the other side, $z$ has two lifts, $w_\pm$ lying over $u=\pm1$,
each of weight $1$, and $\log|u-1|(w_+)=-t$ while
$\log|u-1|(w_-)=0$; indeed $x-1=(u-1)(u+1)$ and $u+1$ is a unit at
$u=1$.  Hence
\[
 \tfrac12\bigl(\varphi(w_+)+\varphi(w_-)\bigr)
 =\tfrac12\bigl(\max\{0,\lambda-t\}+\max\{0,\lambda\}\bigr).
\]
The two expressions agree.  For $\lambda\le0$ both are $0$; for
$0\le\lambda\le t$ both equal $\lambda/2$; for $\lambda\ge t$ both
equal $\lambda-t/2$.  With $\lambda=1$ and $t=3$ both give
$\tfrac12$; with $\lambda=4$ and $t=3$ both give $\tfrac52$.

Note that the middle coefficient $n_{(1,1)}=-2$ is a unit here and does
contribute: at $\lambda=t=1$ the maximum
$\max\{0,\lambda,2\lambda-t\}=\max\{0,1,1\}$ is attained by the middle
term as well as by the last.  The next example shows what happens when that
coefficient is not merely inactive but zero.
\end{example}

\begin{example}[Frobenius: the trace dies, the norm does not]
\label{ex:ch41-frobenius-average}
Let $\operatorname{char}k=p>0$ and let $h(u)=u^p=x$, the relative
Frobenius of $\mathbf P^1$, so $d=p$ and $B=A[u]/(u^p-x)$ as in
\Cref{ex:ch41-cyclic-degree-p}.  This morphism is finite, surjective,
purely inseparable, and a universal homeomorphism.

\emph{The trace.}  The trace map $\Tr_{B/A}:B\to A$ is identically zero.
Indeed $\Tr(1)=p\cdot1=0$, and for $1\le i\le p-1$ the matrix of
multiplication by $u^i$ in the basis $1,u,\dots,u^{p-1}$ has all
diagonal entries $0$, because $u^i\cdot u^j$ is $u^{i+j}$ if
$i+j<p$ and $xu^{i+j-p}$ otherwise, and in neither case is the
coefficient of $u^j$ nonzero when $i\ne0$.  Hence
$\Tr(u^i)=0$ for all $0\le i\le p-1$, and $\Tr\equiv0$ by
$A$-linearity.  Any attempted averaging by
$\frac1p\log|\Tr(\cdot)|$ therefore returns $-\infty$ identically, and
carries no information whatsoever.

\emph{The norm.}  By \Cref{ex:ch41-cyclic-degree-p},
\[
 \mathrm N_{B/A}(T_0+T_1u)=T_0^p+xT_1^p ,
\]
with all mixed coefficients zero.  Take \eqref{eq:ch41-example-data}, so
that
\[
 \begin{aligned}
 \operatorname{Av}_q(\varphi)
 &=\frac1p\max\Bigl\{
     p\lambda_0,
     \lambda_0+(p-1)\lambda_1+\log|0|,
     \dots,\\
 &\hspace{43mm}p\lambda_1+\log|x|
   \Bigr\}\\
 &=\max\Bigl\{\lambda_0,
     \lambda_1+\tfrac1p\log|x|\Bigr\}.
 \end{aligned}
\]
since all the intermediate coefficients vanish and contribute
$-\infty$, while the nonzero extreme coefficients have the indicated values.

\emph{The check.}  Let $z=t\ord_{x=0}$ with $t>0$.  Its unique lift is
$w=\frac tp\ord_{u=0}$, because $p\,\ord_{u=0}(u)=\ord_{u=0}(u^p)$, and
the weight is $e(w/z)=p$, the fibre algebra being
$\mathcal H(z)[u]/(u^p-x)$, a purely inseparable field extension of
degree $p$.  Thus
\[
 \frac1p\,e(w/z)\,\varphi(w)
 =\varphi(w)
 =\max\Bigl\{\lambda_0,\ \lambda_1-\tfrac tp\Bigr\},
\]
which is exactly $\operatorname{Av}_q(\varphi)(z)$, since
$\log|x|(z)=-t$.  With $p=2$, $\lambda_0=0$, $\lambda_1=1$ and
$t=1$ both sides equal $\tfrac12$; with $p=3$, $\lambda_0=0$,
$\lambda_1=1$ and $t=6$ both sides equal $0$.

\emph{At a trivially valued point.}  Let $z=v_{\{x=0\},\triv}$, of proper
support.  Then $\mathcal H(z)=k$, the fibre algebra is
$k[u]/(u^p)$, local of length $p$ and nonreduced, and the unique lift is
$w=v_{\{u=0\},\triv}$ with weight $p$.  Both sides of
\eqref{eq:bj-finite-flat-average} equal
$\max\{\lambda_0,\lambda_1-\infty\}=\lambda_0$, since
$\log|x|(z)=-\infty$ and $\log|u|(w)=-\infty$.  Here the weight comes
from a length, not from a residue degree; the point of the example is that
the same integer $p$ appears in both descriptions, as
\Cref{rem:ch41-weight-anatomy} predicts.
\end{example}

\begin{example}[A nonreduced fibre]
\label{ex:ch41-nonreduced-fibre}
Let $h(u)=u^3-u^2=x$, a morphism $\mathbf P^1\to\mathbf P^1$ of degree
$d=3$, with $B=k[u]$ free over $A=k[x]$ on $1,u,u^2$ and norm
polynomial \eqref{eq:ch41-cubic-norm}:
\[
 \mathrm N_{B/A}(T_0+T_1u)=T_0^3+T_0^2T_1+x\,T_1^3,
 \qquad
 n_{(3,0)}=1,\ n_{(2,1)}=1,\ n_{(1,2)}=0,\ n_{(0,3)}=x .
\]
With the data \eqref{eq:ch41-example-data},
\[
 \begin{aligned}
 \operatorname{Av}_q(\varphi)
 &=\tfrac13\max\bigl\{3\lambda_0,\ 2\lambda_0+\lambda_1,\
 \lambda_0+2\lambda_1+\log|0|,\ 3\lambda_1+\log|x|\bigr\}
 \\
 &=\tfrac13\max\bigl\{3\lambda_0,\ 2\lambda_0+\lambda_1,\
 3\lambda_1+\log|x|\bigr\}.
 \end{aligned}
\tag{41.36}\label{eq:ch41-cubic-average}
\]

\emph{At a divisorial point.}  Let $z=t\ord_{x=0}$, $t>0$.  Since
$x=u^2(u-1)$, the lifts are
\[
 w_0=\tfrac t2\ord_{u=0}\ \ \text{with}\ \ e(w_0/z)=2,
 \qquad
 w_1=t\ord_{u=1}\ \ \text{with}\ \ e(w_1/z)=1 ;
\]
indeed over $\mathcal H(z)=k((x))$ the polynomial $u^3-u^2-x$ has a
factor of degree $2$, ramified with $\ord(u)=t/2$, and a factor of
degree $1$ with $\ord(u-1)=t$.  The weights sum to $3$.  Now
$\log|u|(w_0)=-t/2$ and $\log|u|(w_1)=0$, so
\[
 \frac13\sum_{w\mid z}e(w/z)\varphi(w)
 =\frac13\Bigl(2\max\bigl\{\lambda_0,\lambda_1-\tfrac t2\bigr\}
 +\max\{\lambda_0,\lambda_1\}\Bigr).
\]
This agrees with \eqref{eq:ch41-cubic-average}, whose right-hand side is
$\frac13\max\{3\lambda_0,2\lambda_0+\lambda_1,3\lambda_1-t\}$.  To see the
agreement, expand the left-hand side as
\[
 \tfrac13\max\bigl\{3\lambda_0,\ 2\lambda_0+\lambda_1,\
 \lambda_0+2\lambda_1-t,\ 3\lambda_1-t\bigr\},
\]
and observe that the third entry is redundant: if $\lambda_1\le\lambda_0$
then $\lambda_0+2\lambda_1-t\le3\lambda_0$, and if
$\lambda_1>\lambda_0$ then $\lambda_0+2\lambda_1-t<3\lambda_1-t$.  The
redundant entry is exactly the one whose coefficient
$n_{(1,2)}$ vanishes: the vanishing of a coefficient is harmless because
the corresponding term was never going to be the maximum.  Numerically,
with $\lambda_0=0,\lambda_1=1,t=1$ both sides equal $\tfrac23$; with
$\lambda_0=0,\lambda_1=1,t=4$ both equal $\tfrac13$; with
$\lambda_0=2,\lambda_1=0,t=1$ both equal $2$.

\emph{At a point of proper support, where the fibre is nonreduced.}  Let
$z=v_{\{x=0\},\triv}$.  Then $\mathcal H(z)=k$ with the trivial absolute
value, and
\[
 B_z^{\wedge}=k[u]/(u^3-u^2)=k[u]/(u^2)\times k ,
\]
a product of a local Artinian algebra of length $2$, which is
\emph{nonreduced}, and a copy of $k$.  The associated analytic points are
$w_0=v_{\{u=0\},\triv}$ with weight $2$ and
$w_1=v_{\{u=1\},\triv}$ with weight $1$.  Here
$\log|u|(w_0)=-\infty$ and $\log|u|(w_1)=0$, so
\[
 \frac13\sum_{w\mid z}e(w/z)\varphi(w)
 =\frac13\bigl(2\lambda_0+\max\{\lambda_0,\lambda_1\}\bigr),
\]
while \eqref{eq:ch41-cubic-average} with $\log|x|(z)=-\infty$ gives
$\frac13\max\{3\lambda_0,2\lambda_0+\lambda_1\}$.  These agree, both
being $\frac13(2\lambda_0+\max\{\lambda_0,\lambda_1\})$.  With
$\lambda_0=0,\lambda_1=1$ both equal $\tfrac13$; with
$\lambda_0=1,\lambda_1=0$ both equal $1$.

Three features of this last computation deserve emphasis.  First, the
nonreduced factor occurs at a point of \emph{proper support}, illustrating
the claim made in \Cref{sec:ch41-what}: nonreducedness of fibres and the
value $-\infty$ are met at the same points, and a treatment that
restricted attention to $X^{\val}$ would omit this proper-support
calculation.  This does not assert that every inseparable completed
fibre at full support is reduced, or that every psh function is finite
at all full-support points.  Second, the length $2$ of the nonreduced
factor is not visible in the set of geometric points of the fibre, which has
only two elements; it is visible in the exponent of the linear factor in
\eqref{eq:ch41-norm-factorization}.  Third, the identity holds with
$-\infty$ values on both sides without any convention beyond
$\log0=-\infty$.
\end{example}

\section{Finite normal descent}
\label{sec:ch41-descent}

\subsection{One preliminary on birational pullback of section data}
\label{sec:ch41-birational-gf}

The averaging operator is applied to Fubini--Study data pulled back through
$\alpha:W\to T$, and $W$ need not be normal, so
\Cref{thm:ch39-pullback} is not available.  What is needed is only the
elementary half of that theorem, and it holds for any proper dominant
morphism of integral projective varieties.

\begin{lemma}
\label{lem:ch41-proper-gf-pullback}
Let $g:W\to T$ be a proper dominant morphism of integral projective
$k$-varieties, let $L$ be a $\mathbf Q$-line bundle on $T$ and let
$\Lambda\subseteq\mathbf R$ be an additive subgroup.  Then:
\begin{enumerate}
\item\label{itm:ch41-pgf:centre}
      for every $w\in W^{\An}$, writing $y:=g^{\An}(w)$, one has
      $\xi_y=g(\xi_w)$ and $c(y)=g(c(w))$;
\item\label{itm:ch41-pgf:values}
      for every $m\ge1$ and $s\in H^0(T,mL)$,
      $\log|g^*s|(w)=\log|s|(y)$;
\item\label{itm:ch41-pgf:fs}
      $g^*$ maps $\cH^{\mathrm{gf}}_\Lambda(L)$ into
      $\cH^{\mathrm{gf}}_\Lambda(g^*L)$ and
      $\cH_{\mathrm{hom}}(L)$ into $\cH_{\mathrm{hom}}(g^*L)$,
      preserving the presentation and the weights.
\end{enumerate}
\end{lemma}

\begin{proof}
\cref{itm:ch41-pgf:centre}  The support statement is the definition of
$g^{\An}$, as in
\Cref{lem:ch41-support-centre-finite}\,\cref{itm:ch41-sc:supp}.  For the
centre, the argument of
\Cref{lem:ch41-support-centre-finite}\,\cref{itm:ch41-sc:centre} applies
verbatim: it used only that $g$ is a morphism of schemes, that centres
exist and are unique on proper separated $k$-varieties, which is
\Cref{prop:centers-proper-separated}, and that
$\cO_{T,g(c(w))}\to\cO_{W,c(w)}$ is a local homomorphism.  Finiteness of
$g$ was used in that lemma only for the assertion that
$\kappa(\xi_y)\subseteq\kappa(\xi_w)$ is a finite extension, which is not
claimed here.

\cref{itm:ch41-pgf:values}  Choose a frame $\tau$ of $mL$ near
$c(y)$ and write $s=f\tau$.  By \cref{itm:ch41-pgf:centre} the section
$g^*\tau$ is a frame of $mg^*L$ near $c(w)$, and
$g^*s=(f\circ g)g^*\tau$; hence
$\log|g^*s|(w)=-w(f\circ g)=-y(f)=\log|s|(y)$, by
\Cref{lem:ch34-frame-independence}.

\cref{itm:ch41-pgf:fs}  Let
$\phi=\frac1m\max_j\{\log|s_j|+\lambda_j\}\in
\cH^{\mathrm{gf}}_\Lambda(L)$.  Each $g^*s_j$ is a nonzero section of
$mg^*L$, because $g$ is dominant and $W$ is integral, so
$g^*s_j$ is nonzero at the generic point of $W$.  By
\cref{itm:ch41-pgf:values},
$g^*\phi=\frac1m\max_j\{\log|g^*s_j|+\lambda_j\}$, a weighted section
system with the same weights.  Generic finiteness holds because $W$ is
integral.  The homogeneous case is the case $\Lambda=0$.
\end{proof}

\subsection{The theorem}
\label{sec:ch41-descent-theorem}

\begin{theorem}\label{thm:bj-finite-normal-descent}
Let $h:T\to X$ be finite surjective between normal integral projective
varieties, and let $\theta\in N^1(X)_{\mathbf R}$.  Then
\[
 \operatorname{EP}_T(h^*\theta)
 \quad\Longrightarrow\quad
 \operatorname{EP}_X(\theta).
\tag{41.37}\label{eq:bj-finite-normal-descent}
\]
\end{theorem}

\begin{proof}
The proof has four steps.  Step~0 lifts the regularized supremum to $T$
and uses the assumed envelope property there.  Step~1 averages it down to
the birational model $Z$.  Step~2 averages down a homogeneous carrier.
Step~3 reflects birationally from $Z$ to $X$.  The envelope property is
used exactly twice, both times on $T$ and both times as a hypothesis:
once in Step~0 and once in Step~2, through
\Cref{lem:bj-canonical-carrier}.

\medskip
\noindent\emph{Step 0: transport the family to $T$.}
Let $(\varphi_i)_{i\in I}\subset\PSH_X(\theta)$ be nonempty and bounded
above, and put
\[
                         u_X:=\left(\sup_i\varphi_i\right)^*.
\tag{41.38}\label{eq:bj-finite-envelope-downstairs}
\]
Dominant pullback preserves psh functions.  Analytic openness and source
EP give
\[
 u_T:=h^*u_X
     =\left(\sup_i h^*\varphi_i\right)^*
     \in\PSH_T(h^*\theta).
\tag{41.39}\label{eq:bj-finite-envelope-upstairs}
\]
In detail: by
\Cref{prop:ch41-finite-pullback}\,\cref{itm:ch41-pb:psh} each
$h^*\varphi_i$ lies in $\PSH_T(h^*\theta)$, and the family
$(h^*\varphi_i)$ is nonempty and bounded above by the same constant, by
\Cref{prop:ch41-finite-pullback}\,\cref{itm:ch41-pb:sup}.  The hypothesis
$\operatorname{EP}_T(h^*\theta)$ therefore places
$(\sup_ih^*\varphi_i)^*$ in $\PSH_T(h^*\theta)$.  On the other hand
$\sup_ih^*\varphi_i=h^*(\sup_i\varphi_i)$, because pullback commutes with
suprema by
\Cref{prop:ch41-finite-pullback}\,\cref{itm:ch41-pb:ops}; and
\Cref{lem:bj-finite-normal-open} applied to the bounded-above function
$a=\sup_i\varphi_i$ gives
$(h^*a)^*=h^*(a^*)=h^*u_X$.  This is
\eqref{eq:bj-finite-envelope-upstairs}.  Note that both hypotheses of
\Cref{lem:bj-finite-normal-open} are met: $T$ and $X$ are normal
integral projective, and $h$ is finite dominant.  Note also that
$u_X$ itself is \emph{not} yet known to be psh --- that is the conclusion
--- but it is upper semicontinuous, bounded above, and satisfies
$u_X\ge\varphi_{i_0}$ for any fixed $i_0$, so
\[
 -\infty<\varphi_{i_0}(v_{X,\triv})\le u_X(v_{X,\triv})
 =\sup_{X^{\An}}u_X<+\infty,
\tag{41.40}\label{eq:ch41-uX-finite-at-triv}
\]
because, with $C=\sup_i\varphi_i(v_{X,\triv})$, each member of
the family is at most $C$ by the maximum principle.
Thus its raw supremum is bounded by $C$ and equals $C$ at
$v_{X,\triv}$; upper-semicontinuous regularization remains at most
$C$ and majorizes the raw supremum.  This proves the displayed equality
without applying a psh theorem to the uncertified $u_X$.
The lower bound uses \Cref{def:bj-psh}.

\medskip
\noindent\emph{Step 1: descend the psh function to the birational model.}
Choose the square \eqref{eq:bj-flatification-square}.  Let
$(\phi_j)$ be a decreasing GF regularization of $u_T$ with real weights
and presented rational line-bundle classes
\[
                              \eta_j\longrightarrow h^*\theta.
\]
That such a net exists is \Cref{def:bj-psh}\,(3) applied to
$u_T\in\PSH_T(h^*\theta)$: there are $\mathbf Q$-line bundles $L_j$
on $T$ and $\phi_j\in\cH^{\mathrm{gf}}_{\mathbf R}(L_j)$, decreasing
with $\phi_j\searrow u_T$ pointwise and
$\eta_j:=c_1(L_j)\to h^*\theta$.  The weights of the presentations are
real, not necessarily rational, which is why the averaging proposition is
invoked below in the form \Cref{cor:ch41-real-weights}.

Pull the FS data through $\alpha$.  By
\Cref{lem:ch41-proper-gf-pullback}\,\cref{itm:ch41-pgf:fs}, the functions
$\alpha^*\phi_j$ lie in
$\cH^{\mathrm{gf}}_{\mathbf R}(\alpha^*L_j)$ and form a decreasing net
converging pointwise to $\alpha^*u_T$.  Applying
\Cref{prop:bj-weighted-coefficient-norm} together with
\Cref{cor:ch41-real-weights} to the finite faithfully flat morphism $q$ of
rank $d$ produces the decreasing net
$\operatorname{Av}_q(\alpha^*\phi_j)$ of generically finite Fubini--Study
functions on $Z$, of classes $\frac1dq_*\alpha^*\eta_j$.
The normalized norm classes satisfy
\[
 \begin{aligned}
 \frac1dq_*\alpha^*\eta_j
 &\longrightarrow\frac1dq_*\alpha^*h^*\theta\\
 &=\frac1dq_*q^*\beta^*\theta
  =\beta^*\theta.
 \end{aligned}
\tag{41.41}\label{eq:bj-norm-stage-classes}
\]
The last equality is the projection formula for the finite flat map of
rank $d$, that is, \eqref{eq:ch41-projection-formula} of
\Cref{lem:ch41-projection-formula}; the identity
$\alpha^*h^*\theta=q^*\beta^*\theta$ is
\eqref{eq:bj-flatification-class}; and the convergence in the first line
holds because $\frac1dq_*\alpha^*$ is a linear map between the
finite-dimensional real vector spaces $N^1(T)_{\mathbf R}$ and
$N^1(Z)_{\mathbf R}$, hence continuous.

For every $z\in Z^{\An}$, commutativity of the square gives
\[
 \begin{aligned}
 \frac1d\sum_{w\mid z}e(w/z)u_T(\alpha(w))
 &=\frac1d\sum_{w\mid z}e(w/z)\beta^*u_X(z)\\
 &=\beta^*u_X(z).
 \end{aligned}
\tag{41.42}\label{eq:bj-norm-stage-limit}
\]
The first equality is the computation
\[
 u_T\bigl(\alpha^{\An}(w)\bigr)
 =u_X\bigl(h^{\An}\alpha^{\An}(w)\bigr)
 =u_X\bigl(\beta^{\An}q^{\An}(w)\bigr)
 =\beta^*u_X(z),
\]
valid for every $w$ over $z$ and using only that the square commutes;
the second is $\sum_{w\mid z}e(w/z)=d$, that is,
\Cref{cor:ch41-average-pullback} applied to $\psi=\beta^*u_X$.  In other
words, $\alpha^*u_T=q^*\beta^*u_X$ as functions on $W^{\An}$, and the
averaging operator returns the function it was pulled back from.
At the trivial valuation this value is finite: the original family is
nonempty, every member is finite there, and the family is bounded above.
Precisely, $\beta^*u_X(v_{Z,\triv})=u_X(v_{X,\triv})$, which is finite by
\eqref{eq:ch41-uX-finite-at-triv}, so $\beta^*u_X\not\equiv-\infty$.
\Cref{cor:bj-coefficient-norm-limit} therefore yields
\[
                         u_Z:=\beta^*u_X
                         \in\PSH_Z(\beta^*\theta).
\tag{41.43}\label{eq:bj-uZ-psh}
\]
Only GF pullback and coefficient norms were used on the possibly nonnormal
scheme $W$.  This is worth restating as a warning to the reader: at no
point is a psh theory on $W$ invoked, and no closure theorem on $W$ is imported here:
\Cref{def:bj-psh} and every closure statement of
\Cref{chap:bj-fs-pl-psh} are set up for the varieties of
\eqref{eq:ch35-standing}, while the results of
\Cref{chap:bj-reflection} that we use require normality.  What is used on
$W$ is only \Cref{lem:ch41-proper-gf-pullback}, which is a statement about
sections and semivaluations.

\medskip
\noindent\emph{Step 2: descend a homogeneous carrier.}
The psh space on $T$ is nonempty because it contains $u_T$.  By source
EP and \Cref{lem:bj-canonical-carrier}, the function
\[
 P_T=
 \left(\sup\{v\in\PSH_T(h^*\theta):\sup v\le0\}\right)^*
 \in\PSH_{\mathrm{hom},T}(h^*\theta)
\tag{41.44}\label{eq:bj-finite-source-carrier}
\]
satisfies
\[
                              u_T\le P_T+C
\tag{41.45}\label{eq:bj-uT-carrier}
\]
for a finite $C$, for example $C=\sup u_T$.  This is the second use of the assumed envelope property in the main
descent proof, alongside the regularized family supremum in Step~0; the family in
\eqref{eq:bj-finite-source-carrier} is nonempty because
$u_T-\sup u_T$ belongs to it, and it is bounded above by $0$.

Choose a homogeneous rational GF regularization
$\chi_j\downarrow P_T$, pull its FS data through $\alpha$, and apply
\Cref{cor:bj-coefficient-norm-limit}.  Such a regularization exists by
\Cref{prop:homogenization-gf}: every homogeneous psh function on the
integral projective variety $T$ is the pointwise limit of a decreasing net
$\chi_j\in\cH_{\mathrm{hom}}(L_j')$ with $c_1(L_j')\to h^*\theta$, the
weights of \eqref{eq:homogeneous-fs-presentation} being all zero.  By
\Cref{lem:ch41-proper-gf-pullback}\,\cref{itm:ch41-pgf:fs} the pullbacks
$\alpha^*\chi_j$ again lie in $\cH_{\mathrm{hom}}(\alpha^*L_j')$.
The same class calculation as in
\eqref{eq:bj-norm-stage-classes} gives
\[
 p_Z(z):=\frac1d\sum_{w\mid z}e(w/z)P_T(\alpha(w))
 \in\PSH_{\mathrm{hom},Z}(\beta^*\theta).
\tag{41.46}\label{eq:bj-pZ-carrier}
\]
Three things have to be checked here, and each is available.  The classes
$\frac1dq_*\alpha^*c_1(L_j')$ converge to $\beta^*\theta$ exactly as in
\eqref{eq:bj-norm-stage-classes}.  The limit function $p_Z$ is not
identically $-\infty$: by
\Cref{prop:ch35-homogeneous-normalization} a homogeneous psh function
attains its supremum $0$ at the trivial valuation, so
$P_T(v_{T,\triv})=0$, and the unique point of $W^{\An}$ over
$v_{Z,\triv}$ is $v_{W,\triv}$, with weight $d$, whence
$p_Z(v_{Z,\triv})=P_T(v_{T,\triv})=0$.  Finally $p_Z$ is homogeneous,
because $\operatorname{Av}_q$ commutes with the scaling action by the last
paragraph of the proof of \Cref{prop:bj-weighted-coefficient-norm}, and a
decreasing limit of homogeneous functions is homogeneous.  Hence
$p_Z\in\PSH_{\mathrm{hom},Z}(\beta^*\theta)$ by
\eqref{eq:homogeneous-psh-space}.

Averaging \eqref{eq:bj-uT-carrier} and using
$\alpha^*u_T=q^*u_Z$ yields
\[
                              u_Z\le p_Z+C.
\tag{41.47}\label{eq:bj-uZ-carrier}
\]
Explicitly, apply $\alpha^*$ to \eqref{eq:bj-uT-carrier} to get
$\alpha^*u_T\le\alpha^*P_T+C$ on $W^{\An}$; evaluate both sides at the
points $w$ over a fixed $z\in Z^{\An}$, multiply by the positive
weights $e(w/z)$, sum, and divide by $d$.  On the left one obtains
$u_Z(z)$, by Step~1 and \Cref{cor:ch41-average-pullback}; on the right one
obtains $p_Z(z)+C$, the constant $C$ surviving because
$\frac1d\sum_we(w/z)C=C$.
If one value of $P_T$ in the fibre is $-\infty$, then
\eqref{eq:bj-uT-carrier} forces the corresponding value of $u_T$ to be
$-\infty$, and the determinant identity makes both relevant averages
$-\infty$.  No subtraction of infinite values occurs.  In particular
\eqref{eq:bj-uZ-carrier} is an inequality in $[-\infty,+\infty)$ and is
never rewritten as $u_Z-p_Z\le C$, which would be undefined at the points
where both are $-\infty$; this is the standing prohibition of
\Cref{warn:ch39-no-subtraction}.

\medskip
\noindent\emph{Step 3: reflect birationally.}
Apply carrier-dominated reflection,
\Cref{thm:bj-carrier-reflection}, to $\beta:Z\to X$.  Its hypotheses
\eqref{eq:bj-carrier-reflection-hypotheses} are exactly
\eqref{eq:bj-uZ-psh}, \eqref{eq:bj-pZ-carrier} and
\eqref{eq:bj-uZ-carrier}, with $Z$ and $X$ normal integral projective
and $\beta$ projective birational by
\Cref{lem:bj-finite-flatification}.  It produces
$v\in\PSH_X(\theta)$ with
\[
                         \beta^*v=u_Z=\beta^*u_X.
\]
Since $\beta^{\An}$ is surjective, $v=u_X$.  Thus the regularized
supremum of the original family is psh, proving
\eqref{eq:bj-finite-normal-descent}.

The proof assumed neither EP on $Z$ nor EP on $W$.  It did not define
psh functions on the nonnormal $W$, and it did not apply a determinant
norm directly to the generally nonflat morphism $h$.
\end{proof}

\begin{remark}[Where each hypothesis of the theorem is used]
\label{rem:ch41-descent-hypothesis-map}
Normality of $X$ supplies the analytic quotient/radicial argument
and the target hypothesis of carrier reflection.  Normality of $T$
is retained for the source psh, regularization and carrier inputs;
it is not needed for the analytic openness assertion, as
\Cref{rem:ch41-openness-scope} explains.  Flattening gives a normal
$Z$ but leaves $W$ integral, which suffices for the FS pullback and
norm computations.  Projectivity supplies compact analytifications,
centres, numerical classes and projective modifications.
If $\theta$ is not pseudoeffective, $\PSH_X(\theta)$ is empty and EP
is vacuous.  No nefness hypothesis is imposed on the class; moving
classes and the nontrivial limit conditions are checked explicitly.
\end{remark}

\begin{warning}{The implication is one-way as proved}
\label{warn:ch41-one-way}
\Cref{thm:bj-finite-normal-descent} proves
$\operatorname{EP}_T(h^*\theta)\Rightarrow\operatorname{EP}_X(\theta)$
and says nothing about the converse.  The converse is in fact true once the
full envelope theorem is available, since $T$ is then also normal
integral projective and unibranch; but that deduction belongs to
\Cref{chap:bj-envelope-consequences} and may not be used here, by
\Cref{warn:ch41-firewall}.  The identity
$\operatorname{Av}_q\circ q^*=\operatorname{id}$ is a left-inverse
identity; it does not make every upstairs function a pullback.
Fibrewise constancy is still required, as the cone example below
illustrates.
\end{warning}

\subsection{Consequences}
\label{sec:ch41-descent-consequences}

\begin{corollary}
\label{cor:ch41-finite-tower}
Let $T''\xrightarrow{h_2}T'\xrightarrow{h_1}X$ be finite surjective
morphisms of normal integral projective varieties and let
$\theta\in N^1(X)_{\mathbf R}$.  Then
\[
 \operatorname{EP}_{T''}\bigl((h_1h_2)^*\theta\bigr)
 \ \Longrightarrow\
 \operatorname{EP}_{T'}(h_1^*\theta)
 \ \Longrightarrow\
 \operatorname{EP}_X(\theta).
\]
\end{corollary}

\begin{proof}
Apply \Cref{thm:bj-finite-normal-descent} to $h_2$ with the class
$h_1^*\theta$, and then to $h_1$ with the class $\theta$.  Both
morphisms are finite surjective between normal integral projective
varieties, and $(h_1h_2)^*=h_2^*h_1^*$.
\end{proof}

\begin{corollary}
\label{cor:ch41-finite-bijection}
In the setting of \Cref{thm:bj-finite-normal-descent}, assume
$\operatorname{EP}_T(h^*\theta)$.  Then
\[
 h^*:\PSH_X(\theta)\longrightarrow
 \bigl\{u\in\PSH_T(h^*\theta)\ \big|\
 u\text{ is constant on the fibres of }h^{\An}\bigr\}
\]
is a bijection, with inverse induced by
$u\mapsto\bigl(z\mapsto u(w)\ \text{for any }w\mapsto z\bigr)$.
\end{corollary}

\begin{proof}
Injectivity and the fact that the image consists of fibrewise constant
functions are
\Cref{prop:ch41-finite-pullback}\,\cref{itm:ch41-pb:order} and the
definition of $h^*$.  For surjectivity, let $u\in\PSH_T(h^*\theta)$ be
constant on the fibres of $h^{\An}$, and let
$a:X^{\An}\to\mathbf R\cup\{-\infty\}$ be the unique function with
$h^*a=u$; it exists because $h^{\An}$ is surjective, by
\Cref{lem:ch41-finite-surjective}, and $u$ is fibrewise constant.  It is
upper semicontinuous: $\{a\ge c\}=h^{\An}(\{u\ge c\})$ is the image of a
closed, hence compact, set under a continuous map into a Hausdorff space,
hence closed.  It is bounded above by $\sup u$.

Now run the proof of \Cref{thm:bj-finite-normal-descent} with the one-member
family $\{a\}$ replaced as follows.  Since $a$ is usc and $u$ is psh,
we have $a=a^*$ and $h^*a=u=(h^*a)^*$.  Steps~1 to~3 of that proof used
about $u_X$ only that it is usc, bounded above, finite at
$v_{X,\triv}$, and that $h^*u_X\in\PSH_T(h^*\theta)$; all four hold for
$a$, the third because $a(v_{X,\triv})=u(v_{T,\triv})>-\infty$ by
\Cref{prop:ch41-finite-pullback}\,\cref{itm:ch41-pb:sup}.  The conclusion of
Step~3 is $a\in\PSH_X(\theta)$, and $h^*a=u$.
\end{proof}

\begin{example}[Descent along the double cover of the quadric cone]
\label{ex:ch41-descent-cone}
Return to \Cref{ex:ch41-nonflat-cone} and
\Cref{ex:ch41-flatten-cone}: $h:\mathbf P^2\to\mathbf P(1,1,2)$ of degree
$2$, flattened by $q:\mathbf F_1\to\mathbf F_2$.  Let
$H$ be the hyperplane class of $X=\mathbf P(1,1,2)\subset\mathbf P^3$,
which in weighted coordinates is $\cO_X(2)$, with global sections spanned
by $a^2,ab,b^2,z$; it is a genuine line bundle, whereas $\cO_X(1)$ is
only a Weil divisor class, not invertible at the vertex.  Put
$\theta:=c_1(H)$.  Then
$h^*H=\cO_{\mathbf P^2}(2)$, since $h^*a=u$, $h^*b=v$ and
$h^*z=w^2$.  Since $\mathbf P^2$ is
smooth, $\operatorname{EP}_{\mathbf P^2}(h^*\theta)$ is available from
\Cref{thm:bj-smooth-envelope}, and \Cref{thm:bj-finite-normal-descent}
yields $\operatorname{EP}_X(\theta)$ for the singular normal surface
$X=\mathbf P(1,1,2)$.  This is the pattern of
\Cref{chap:bj-alterations}: a singular normal target, a smooth source, and
a finite morphism between them.

It is instructive to see what the averaging does on a concrete function.
Let $s''\in H^0(\mathbf P^2,\cO(1))$ cut out the line
$\{u=w\}$, and put $\varphi'':=\log|s''|\in
\cH^{\mathrm{gf}}(\cO(1))$.  This line is \emph{not} $\iota$-invariant:
$\iota$ carries it to $\{u=-w\}$.  Consequently $\varphi''$ is not
constant on the fibres of $h^{\An}$, and it is not the pullback of any
function on $X^{\An}$.  Its average is nevertheless computable in closed
form.  On the split étale locus, a point $y\in X^{\An}$ has two lifts
$t_1,t_2\in T^{\An}$ of weight $1$, interchanged by $\iota$, so the
fibre average of $\varphi''$ over $y$ is
\[
 \begin{aligned}
 \tfrac12\bigl(\log|u-w|(t_1)+\log|u-w|(t_2)\bigr)
 &=\tfrac12\bigl(\log|u-w|+\log|u+w|\bigr)(t_1)\\
 &=\tfrac12\log\bigl|u^2-w^2\bigr|(t_1),
 \end{aligned}
\]
because $\iota^*(u-w)=u+w$ and $\log|\cdot|$ is multiplicative by
\Cref{lem:ch34-log-s-properties}\,\cref{itm:ch34-logs:product}.  Since
$h^*(a^2-z)=u^2-w^2$, the right-hand side is
$\tfrac12\log|a^2-z|(y)$.  This split calculation is only a check, not a
proof at an arbitrary analytic point: away from the branch divisor an
analytic fibre may still consist of one nonsplit point of weight $2$.

For the global identity, first identify the norm line bundle by its
divisor.  Put $M=\alpha^*\cO(1)$.  The divisor $C$ of $\alpha^*w$
is the strict transform of $\{w=0\}$, which avoids the blowup centre.
The involution fixes $C$ pointwise, and $q$ maps it isomorphically
onto the divisor $B=\{z=0\}$ on $Z$.  The blowup $\beta$ is an
isomorphism along $B$, since the vertex has $z\ne0$.  Hence
\cref{itm:ch41-nl:divisor} gives
\[
 \operatorname{div}\mathrm N_q(\alpha^*w)
 =q_*C=B=\operatorname{div}(\beta^*z).
\]
Since $Z$ is normal, these equal Cartier divisors identify
$\mathrm N_q(M)$ with $\beta^*H$, carrying
$\mathrm N_q(\alpha^*w)$ to $\beta^*z$.  This also rules out an
unaccounted exceptional divisor in the identification.

Now work in the rational frame $w$ of $\cO(1)$ and put $s=u/w$.
The involution sends $s$ to $-s$, so the field norm of the coefficient
of $u-w$ is
\[
 N_{K(T)/K(X)}(s-1)
 =(s-1)(-s-1)=1-s^2=\frac{z-a^2}{z}.
\]
Under the established line-bundle identification the norm section is
therefore $\beta^*(z-a^2)$.  Multiplication by the unit $-1$
does not affect section log values.  The all-point norm formula,
including nonsplit and nonreduced fibres, now yields
\[
 \operatorname{Av}_q(\alpha^*\varphi'')
 =\tfrac12\,\beta^*\log\bigl|a^2-z\bigr| .
\]
This is the degree-two determinant calculation of
\Cref{ex:ch41-rank-two}, applied to the coefficient in an actual
rational frame; homogeneous coordinates themselves are sections,
not elements of the function field.

The classes match, as they must.  On one side,
\eqref{eq:bj-norm-class} gives
\[
 \bigl[\operatorname{Av}_q(\alpha^*\varphi'')\bigr]
 =\tfrac12q_*\alpha^*c_1(\cO(1))
 =\tfrac12q_*\alpha^*\bigl(\tfrac12h^*\theta\bigr)
 =\tfrac14q_*q^*\beta^*\theta
 =\tfrac12\beta^*\theta,
\]
using \eqref{eq:bj-flatification-class} and
\eqref{eq:ch41-projection-formula}.  On the other side,
$\tfrac12\beta^*\log|a^2-z|$ has class
$\tfrac12\beta^*c_1(H)=\tfrac12\beta^*\theta$.  The example shows that the
operator does not simply ``restrict invariant data'': it produces sections
of the norm bundle, whose degree is the pushforward, and the resulting
numerical class is $\frac1dq_*$, not the naive pullback-inverse.
\end{example}

\section{Proper support, the extended-real audit, and checkpoints}
\label{sec:ch41-audit}

\subsection{Values at proper-support semivaluations}
\label{sec:ch41-proper-support-audit}

The book plan lists ``values at proper-support semivaluations'' and ``every
use of subtraction involving an infinite value'' among the points to be
audited in the envelope proof.  They interact in the norm calculation and carrier reflection, so
we record five checkpoints.

The following are five places where a proper-support point could have been
mishandled.

\emph{(1) The description of analytic points.}
\Cref{lem:ch41-point-pair} presents a point of $X^{\An}$ as a pair
consisting of a scheme point and an absolute value on its residue field, not
as a valuation of $K(X)$.  Every subsequent argument about lifting points
--- the orbit computation in \Cref{lem:bj-finite-normal-open}, the fibre
decomposition in \Cref{lem:ch41-fibre-decomposition} --- begins with
$\kappa(\xi)$ for an arbitrary $\xi$, not with $K(X)$.  Had we
restricted to $X^{\val}$, the openness lemma would have been proved on a
dense subset only, which is not enough for
\eqref{eq:bj-finite-open-usc}: upper-semicontinuous regularization is
computed with neighbourhoods in $X^{\An}$, and the regularized value at a
proper-support point is not determined by the values on $X^{\val}$
without a separate density argument.

\emph{(2) The fibre algebra.}
Over a proper-support point the fibre algebra $B_z$ of
\Cref{def:ch41-fibre-algebra} is the scheme-theoretic fibre, with all its
nilpotents; over a full-support point the scheme fibre is the function field
of $W$, and in the generically étale case its completed algebra is
a product of fields.  This is the precise sense in which nonreduced fibres and
proper support meet, and it is exhibited in
\Cref{ex:ch41-nonreduced-fibre}.

\emph{(3) The polar locus of the average.}
By Step~4 of the proof of \Cref{prop:bj-weighted-coefficient-norm}, the
average takes the value $-\infty$ at $z$ if and only if $\varphi$
takes it at some point over $z$.  Nothing is discarded: the coefficient
sections $c_I$ vanish simultaneously exactly on the image of the polar
locus upstairs.

\emph{(4) The carrier inequality.}
Inequality \eqref{eq:bj-uT-carrier} and its consequence
\eqref{eq:bj-uZ-carrier} are inequalities in $[-\infty,+\infty)$.  They
are never converted into a statement about a difference.  Where
$P_T=-\infty$, the inequality forces $u_T=-\infty$, and the averaged
inequality \eqref{eq:bj-uZ-carrier} degenerates to
$-\infty\le-\infty+C$, which is true and harmless.

\emph{(5) The final identification.}
Step~3 of the proof concludes $v=u_X$ from $\beta^*v=\beta^*u_X$ and
surjectivity of $\beta^{\An}$.  This is an equality of functions on all of
$X^{\An}$, including proper support, and it does not pass through
$X^{\divis}$; in particular it does not use
\Cref{prop:ch39-determination-package}.  A divisorial equality with
an uncertified psh function would not by itself establish equality
at every point.

\begin{warning}{Two identities that must not be written}
\label{warn:ch41-forbidden-identities}
The following manipulations are not justified in the required
extended-real domain.

First, $\operatorname{Av}_q(\varphi)-\operatorname{Av}_q(\psi)
=\operatorname{Av}_q(\varphi-\psi)$.  Upstairs the difference is
undefined where both functions are $-\infty$; for FS stages these
are their common polar points, described by
\Cref{prop:ch34-generic-finiteness}\,\cref{itm:ch34-gf:polar}.
Downstairs the two averages can both be $-\infty$ even when the
upstairs poles occur at different lifts of the same point.
Thus neither side is an unrestricted identity in the extended-real
domain.  What is used instead is order-preservation and the
identity $\operatorname{Av}_q(\varphi+C)=\operatorname{Av}_q(\varphi)+C$
for a real constant $C$, as in \eqref{eq:bj-uZ-carrier}.
For finite-valued arbitrary functions the fibre average is linear;
the problem here is that differences at common poles are undefined,
and a difference of FS functions need not be FS data.

Second, $\log\bigl|\sum_jc_j\bigr|=\max_j\log|c_j|$.  Only $\le$ holds,
by \Cref{lem:ch34-log-s-properties}\,\cref{itm:ch34-logs:sum}, and equality
can fail; this is precisely why the trace, which is a sum, cannot implement
an average, and why the identity actually used is the multiplicativity of
the Gauss norm, \Cref{lem:ch41-gauss-multiplicative}, which is an equality.
\end{warning}

\subsection{What fails if a safeguard is dropped}
\label{sec:ch41-safeguard-failures}

\emph{Drop the flattening.}  In the quadric-cone example, the finite
algebra at the vertex has generic rank $2$, but its fibre is
$k[s,t]/(s^2,st,t^2)$ of dimension $3$.
The generic field norm of a scalar is its square, while the determinant
on this fibre algebra is its cube.  Thus the rank-two determinant
package cannot be transported to this fibre by arbitrary base change.
A generic norm on the normal target can exist, and the norm of
$\cO_{\mathbf P^2}(1)$ in this example even extends as $H$.
There is no claimed line-bundle obstruction from the local class group.
Flattening supplies the constant-rank, base-change-compatible package
actually used in this proof.

\emph{Drop the local Artin lengths.}  At the proper-support point
$z=v_{\{x=0\},\triv}$ of the cubic example, the two lifts have weights
$2,1$.  With $\lambda_0=0,\lambda_1=1$, their function values are
$0,1$, so the actual average is $1/3$.
An average over the two distinct points normalized by $2$ gives $1/2$.
If one keeps the denominator $d=3$ but assigns weight $1$ to each point,
even constants are not preserved.  Both attempts lose the Artin lengths
that the determinant exponent records.

\emph{Drop the openness of $h^{\An}$.}  Then Step~0 of the proof of
\Cref{thm:bj-finite-normal-descent} breaks: one still knows
$(\sup_ih^*\varphi_i)^*\in\PSH_T(h^*\theta)$, but no longer that it equals
$h^*u_X$.  \Cref{ex:ch41-node-not-open} shows that the two can differ, and
in fact differ by more than a null set.  Everything after Step~0 would then
prove a statement about a function on $T$ that is not the pullback of the
object under study.

\emph{Drop the carrier.}  The established
\Cref{thm:bj-carrier-reflection} then cannot be applied in Step~3.
The unconditional homogeneous reflection theorem does not address a
function $u_Z$ that is not assumed homogeneous.
This identifies a missing hypothesis of the available proof; it does
not disprove a stronger reflection theorem.
In particular an example of a function that is not fibrewise constant
cannot refute descent for fibrewise constant functions.
Step~2 constructs the carrier needed here from source EP.

\emph{Drop normality of $X$.}  Analytic openness and commutation
with regularization can fail, by \Cref{ex:ch41-node-not-open},
so the present proof no longer applies.  That bounded-function
example does not by itself disprove EP;
\Cref{sec:ch37-node-cusp} likewise distinguishes a normalization
reflection obstruction from a converse envelope theorem.
The extension to unibranch nonnormal targets is proved in
\Cref{chap:bj-universal-homeomorphisms,chap:bj-subintegral-normalization}.

\subsection{Checkpoints}
\label{sec:ch41-checkpoints}

\begin{enumerate}
\item Analytic openness is proved on all support strata.  Replacing
      analytic points by valuations of $K(X)$ would omit the points
      needed for usc regularization.
\item The purely inseparable part is handled by uniqueness of extension,
      not by an unproved separability reduction.
\item The flattening step is performed before taking determinant norms.
      Finite flatness supplies a constant rank and a norm line compatible
      with base change.
\item The universal polynomial
      $\mathrm N_q(\sum T_js_j)$ packages all sheets at once.  This is why
      the construction makes sense for inseparable and nonreduced fibres.
\item Local Artin lengths appear in $e(w/z)$.  Dropping them would make
      both the degree identity and the fibre-average formula false.
\item Common zeros are retained as $-\infty$.  The determinant
      polynomial vanishes exactly when a fibre factor vanishes, so no
      proper-support point is silently discarded.
\item The number $1/d$ normalizes a real weight and a numerical class;
      it is not division inside the finite algebra.
\item The final reflection is birational and uses the carrier from source
      EP.  No envelope property on either space in the flattening square
      is smuggled into the proof.
\end{enumerate}

\section{Chapter dictionary}
\label{sec:ch41-dictionary}

The correspondences established or used in this chapter are the following.
Each row records a statement with the hypotheses under which it was proved.
Rows one to three are the analytic dictionary of
\Cref{sec:ch41-finite-analytic}; rows four and five are the algebraic
constructions; rows six and seven are the averaging identity and the
descent.

\begin{center}
\begin{tabular}{p{.29\textwidth}p{.31\textwidth}p{.29\textwidth}}
\toprule
analytic datum & algebraic datum & hypothesis\\
\midrule
a point of $X^{\An}$ &
a pair (scheme point, absolute value on its residue field) trivial on
$k$ &
none (\Cref{lem:ch41-point-pair})\\
the fibre of $h^{\An}$ over $x\in X^{\An}$ &
a $G$-orbit for the Galois closure, or a single point after a radicial
step &
$T,X$ normal integral projective, $h$ finite dominant
(\Cref{lem:bj-finite-normal-open})\\
openness of $h^{\An}$, hence
$(h^*a)^*=h^*(a^*)$ &
$B^G=A$ on $G$-stable charts, plus uniqueness of the extension through
a purely inseparable step &
same; it fails for non-normal targets
(\Cref{ex:ch41-node-not-open})\\
the flattened cover $q:W\to Z$ &
the strict transform of $h_*\cO_T$ under a $U$-admissible blowup &
$U$ the flat locus; Raynaud--Gruson
(\Cref{lem:bj-finite-flatification}, \Cref{thm:appB-flattening})\\
the coefficient sections $c_I$ &
the coefficients of the universal norm polynomial
$\mathrm N_q(\sum_jT_js_j)$ &
$q$ finite locally free of constant rank $d$
(\Cref{lem:ch41-universal-polynomial},
\Cref{prop:ch41-norm-base-change})\\
$\operatorname{Av}_q(\varphi)(z)$ &
$\frac1d\sum_{w\mid z}e(w/z)\varphi(w)$, with $e(w/z)$ the dimension
over $\mathcal H(z)$ of the local factors &
no separability, no reducedness, no restriction on the characteristic
(\Cref{prop:bj-weighted-coefficient-norm})\\
$\operatorname{EP}_T(h^*\theta)$ &
$\operatorname{EP}_X(\theta)$ &
one-way as proved here
(\Cref{thm:bj-finite-normal-descent}, \Cref{warn:ch41-one-way})\\
\bottomrule
\end{tabular}
\end{center}

Five separations should be carried forward.

\emph{Weight versus number of sheets.}  The integer $e(w/z)$ is a
dimension of a local factor of a fibre algebra over a complete valued
field.  It coincides with a sum of geometric Artin lengths after algebraic closure.
When the resulting fibre is reduced, these lengths are all one;
several geometric points can still determine one analytic lift.  In
\Cref{ex:ch41-frobenius-average} it is $p$ with a single sheet; in
\Cref{ex:ch41-nonreduced-fibre} it is $2$ at a nonreduced factor and
$1$ at a reduced one.  The invariant identity is
$\sum_we(w/z)=d$, never a count.

\emph{Determinant versus trace.}  The determinant is multiplicative and
converts the ultrametric identity for products of Gauss norms into the
weighted average; the trace is additive and converts the ultrametric
\emph{inequality} for sums into a bound in the wrong direction, and it
vanishes identically for the Frobenius.  The two are not variants of one
another (\Cref{rem:ch41-norm-vs-trace},
\Cref{ex:ch41-frobenius-average}).

\emph{Flat versus finite.}  Finiteness of $h$ is a hypothesis of the
theorem; flatness is not, and generally fails
(\Cref{ex:ch41-nonflat-cone}).  The determinant construction requires
flatness, so it is applied to $q$ and not to $h$; the two are related by
a birational square, not by an isomorphism
(\Cref{lem:bj-finite-flatification},
\Cref{warn:ch41-flattening-scope}).

\emph{Polynomial law versus polynomial function.}  The norm is a natural
transformation of functors of $A$-algebras
(\Cref{def:ch41-polynomial-law}), whose local basis formulas agree under changes of basis and
under all base changes.  The distinction is what makes the coefficients $c_I$ glue, what
makes the degree exactly $d$, and what makes the fibrewise computation
legitimate after the composite
$A\to\kappa(\xi)\to\mathcal H(z)\to\Omega$;
the residue-field step need not be flat, while the field extensions are
flat and need not be finite (\Cref{rem:ch41-law-work}).

\emph{Envelope property as hypothesis versus as theorem.}  The descent proof consumes EP on the source $T$ as a displayed
hypothesis and proves it on $X$ as its conclusion.  The conclusion carries no
hypothesis on $X$ beyond normality and projectivity.  Reading
\Cref{thm:bj-final-envelope} back into Step~0 or Step~2 would make
\Crefrange{chap:bj-cartier-realization}{chap:bj-subintegral-normalization}
circular (\Cref{warn:ch41-firewall}).

\section{Source note}
\label{sec:ch41-source-note}

The openness, determinant-norm, flattening, and finite-descent arguments
follow the corresponding section of the supplied unpublished envelope
manuscript, with the provenance boundary recorded in
\Cref{sec:appL-envelope}.
The algebraic normalization and flattening facts are cited there by stable
Stacks Project tags under \cite{stacks-project}; the determinant norm is
also treated by \cite{FerrandNorm}.  The presentation here keeps the
universal polynomial law and its base-change property explicit in order to
show precisely why neither separability nor reduced fibres are required.

The precise external attributions are as follows.  Flattening by blowup is
Raynaud--Gruson \cite{RaynaudGruson1971}, imported through
\Cref{thm:appB-flattening} in the form
\cite[\href{https://stacks.math.columbia.edu/tag/0815}{Tag 0815}]
{stacks-project}; the finite-flat criterion is
\cite[\href{https://stacks.math.columbia.edu/tag/02K9}{Tag 02K9}]
{stacks-project}, integrality of blowups is
\cite[\href{https://stacks.math.columbia.edu/tag/01OF}{Tag 01OF}]
{stacks-project}, and finiteness of normalization over a Nagata scheme is
\cite[\href{https://stacks.math.columbia.edu/tag/035S}{Tag 035S}]
{stacks-project}.  The group-theoretic inputs to
\Cref{lem:bj-finite-normal-open} are
\cite[\href{https://stacks.math.columbia.edu/tag/0BRI}{Tag 0BRI}]
{stacks-project} and
\cite[\href{https://stacks.math.columbia.edu/tag/0BRJ}{Tag 0BRJ}]
{stacks-project}; the characteristic-$p$ criterion for a universal
homeomorphism used in \Cref{lem:ch41-radicial-homeo} is
\cite[\href{https://stacks.math.columbia.edu/tag/0CNF}{Tag 0CNF}]
{stacks-project}, imported through
\Cref{thm:appA-universal-homeomorphism-criterion}.  Each of these seven
tags is audited in \Cref{app:stacks}, with the exact statement, the
hypotheses to be checked, and the first place of application; the chain
diagram in \Cref{sec:stacks-dependency-chains} shows how they combine.  The
classical valuation-theoretic facts --- existence of an extension of an
absolute value to an algebraic extension, uniqueness over a complete field,
and the structure of the integral closure of a valuation ring --- are
\cite[Chapter~VI]{ZariskiSamuelII}.  The determinant norm of a line bundle
and its multiplicativity are Ferrand's norm functor
\cite[Proposition~3.3]{FerrandNorm}; the
norm of a function is
\cite[\href{https://stacks.math.columbia.edu/tag/0BD2}{Tag 0BD2}]
{stacks-project}.  The distinct existence theorem for a norm under a
finite surjective morphism to a normal integral target is
\cite[\href{https://stacks.math.columbia.edu/tag/0BD3}{Tag 0BD3}]
{stacks-project}; this explains why lack of flatness alone is not a
line-bundle obstruction, while it does obstruct the constant-rank
base-change argument used here.
Miracle flatness, used in \Cref{ex:ch41-flatten-cone} and in the
smooth curve covers of \Cref{sec:ch41-averaging-examples}, is \cite[Theorem~23.1]{Mat89}; the projection
formula for proper pushforward is \cite[Propositions~1.4 and~2.3]{Ful}.  The
psh formalism, the regularization statements, and the Legendre and carrier
technology are those of \cite{BJ22}, proved in
\Cref{chap:bj-fs-pl-psh,chap:bj-homogeneity-bdivisors,chap:bj-reflection};
the envelope property was formulated in \cite{BJ22} as Boucksom--Jonsson's
Conjecture~5.14; \cite{BJ22b} records partial results available in that
literature.  The proof developed in this book is
\Cref{thm:bj-final-envelope}, with historical context in
\Cref{chap:bj-envelope-history}.

The general theory of polynomial laws in the sense of
\Cref{def:ch41-polynomial-law} is due to Roby \cite{Roby1963}, and the
statement that degree-$d$ homogeneous laws $V\rightsquigarrow U$ correspond to
$A$-linear maps from the $d$-th divided power $\Gamma^d_A(V)$ to
$U$ is his theorem.  We use only the definition and the elementary
uniqueness assertion proved in \Cref{prop:ch41-norm-law}; the classification
is recorded solely for provenance.  Ferrand \cite{FerrandNorm}, who develops
the norm functor in this language, gives the corresponding exposition for
determinant norms.

Six aspects of the argument are isolated here.  The presentation of an analytic point as a pair,
\Cref{lem:ch41-point-pair}, and the compatibility of supports and centres
with a finite morphism, \Cref{lem:ch41-support-centre-finite}, are isolated
so that every later fibre computation can begin at an arbitrary support
prime; this is what makes \Cref{rem:ch41-openness-scope} a checkable
statement rather than a slogan.  The conjugacy lemma
\Cref{lem:ch41-value-transitivity} is proved from Tag 0BRI rather than
quoted, so that the reduction to the purely inseparable case is visible.
The fibre decomposition \Cref{lem:ch41-fibre-decomposition}, with its
definition \eqref{eq:ch41-weights} of the weights and the proof that every
analytic lift is realized, is stated separately because the weight is the
one object in the chapter that a reader is likely to misidentify.  The
Gauss-norm lemma \Cref{lem:ch41-gauss-multiplicative} is stated with
arbitrary positive radii and including the zero polynomial, which is what
makes the $-\infty$ bookkeeping of Step~4 automatic.  The base-change
proposition \Cref{prop:ch41-norm-base-change} and the discussion
\Cref{rem:ch41-law-work} of what would go wrong without it answer the audit
item on universal properties of determinant norms.  Finally the worked
computations \Cref{ex:ch41-nonflat-cone}, \Cref{ex:ch41-flatten-cone},
\Cref{ex:ch41-separable-average}, \Cref{ex:ch41-frobenius-average},
\Cref{ex:ch41-cubic-norm}, \Cref{ex:ch41-nonreduced-fibre} and
\Cref{ex:ch41-node-not-open} are included as worked checks because the two
generality claims that this chapter exists to establish --- no separability,
no reducedness --- are exactly the claims that a plausible argument can
appear to establish without doing so, and because a reader who has verified
the identity \eqref{eq:bj-finite-flat-average} on three explicit covers of
curves will not later mistake the weights for a count of points.
